\input{preamble}

% OK, start here
%
\begin{document}

\title{Cohomologie de de Rham}


\maketitle

\phantomsection
\label{section-phantom}

\tableofcontents

\section{Introduction}
\label{section-introduction}

\noindent
Dans ce chapitre, nous commençons par étudier le complexe de de Rham
d'un morphisme de schémas et nous terminons en démontrant que la cohomologie
de de Rham définit une théorie cohomologique de Weil lorsque le corps de base est de caractéristique nulle.




\section{Le complexe de de Rham}
\label{section-de-rham-complex}

\noindent
Soit $p : X \to S$ un morphisme de schémas. Il existe un complexe
$$
\Omega^\bullet_{X/S} =
\mathcal{O}_{X/S} \to \Omega^1_{X/S} \to \Omega^2_{X/S} \to \ldots
$$
de $p^{-1}\mathcal{O}_S$-modules, où
$\Omega^i_{X/S} = \wedge^i(\Omega_{X/S})$
est placé en degré $i$ et où la différentielle est déterminée par la règle
$\text{d}(g_0 \text{d}g_1 \wedge \ldots \wedge \text{d}g_p) =
\text{d}g_0 \wedge \text{d}g_1 \wedge \ldots \wedge \text{d}g_p$
sur les sections locales.
Voir Modules, section \ref{modules-section-de-rham-complex}.

\medskip\noindent
Étant donné un diagramme commutatif
$$
\xymatrix{
X' \ar[r]_f \ar[d] & X \ar[d] \\
S' \ar[r] & S
}
$$
de schémas, on dispose de morphismes canoniques de complexes
$f^{-1}\Omega_{X/S}^\bullet \to \Omega^\bullet_{X'/S'}$ et
$\Omega_{X/S}^\bullet \to f_*\Omega^\bullet_{X'/S'}$.
Voir Modules, section \ref{modules-section-de-rham-complex}.
Par linéarisation, on obtient, pour tout $p$, un morphisme linéaire
$f^*\Omega^p_{X/S} \to \Omega^p_{X'/S'}$.

\medskip\noindent
En particulier, si $f : Y \to X$ est un morphisme de schémas au-dessus
d'un schéma de base $S$, on dispose d'un morphisme de complexes
$$
\Omega^\bullet_{X/S} \longrightarrow f_*\Omega^\bullet_{Y/S}
$$
Par linéarisation, on voit que, pour tout $p \geq 0$, on obtient un morphisme canonique
$$
\Omega^p_{X/S} \otimes_{\mathcal{O}_X} f_*\mathcal{O}_Y
\longrightarrow
f_*\Omega^p_{Y/S}
$$

\begin{lemma}
\label{lemma-base-change-de-rham}
Soit
$$
\xymatrix{
X' \ar[r]_f \ar[d] & X \ar[d] \\
S' \ar[r] & S
}
$$
un diagramme cartésien de schémas. Les morphismes considérés
ci-dessus induisent alors des isomorphismes
$f^*\Omega^p_{X/S} \to \Omega^p_{X'/S'}$.
\end{lemma}

\begin{proof}
Combiner Morphismes, lemme \ref{morphisms-lemma-base-change-differentials},
avec le fait que la formation des puissances extérieures commute au changement de base.
\end{proof}

\begin{lemma}
\label{lemma-etale}
Considérons un diagramme commutatif de schémas
$$
\xymatrix{
X' \ar[r]_f \ar[d] & X \ar[d] \\
S' \ar[r] & S
}
$$
Si $X' \to X$ et $S' \to S$ sont étales, les morphismes considérés
ci-dessus induisent des isomorphismes
$f^*\Omega^p_{X/S} \to \Omega^p_{X'/S'}$.
\end{lemma}

\begin{proof}
On a $\Omega_{S'/S} = 0$ et $\Omega_{X'/X} = 0$ ; voir par exemple
Morphismes, lemme \ref{morphisms-lemma-etale-at-point}. Les suites exactes courtes de
Morphismes, lemmes
\ref{morphisms-lemma-triangle-differentials} et
\ref{morphisms-lemma-triangle-differentials-smooth},
montrent alors que $\Omega_{X'/S'} = \Omega_{X'/S} = f^*\Omega_{X/S}$.
On conclut en prenant les puissances extérieures.
\end{proof}






\section{Cohomologie de de Rham}
\label{section-de-rham-cohomology}

\noindent
Soit $p : X \to S$ un morphisme de schémas. On appelle
{\it cohomologie de de Rham de $X$ sur $S$} les groupes de cohomologie
$$
H^i_{dR}(X/S) = H^i(R\Gamma(X, \Omega^\bullet_{X/S}))
$$
Puisque $\Omega^\bullet_{X/S}$ est un complexe de $p^{-1}\mathcal{O}_S$-modules,
ces groupes de cohomologie sont naturellement des modules sur $H^0(S, \mathcal{O}_S)$.

\medskip\noindent
Étant donné un diagramme commutatif
$$
\xymatrix{
X' \ar[r]_f \ar[d] & X \ar[d] \\
S' \ar[r] & S
}
$$
de schémas, les morphismes canoniques de la section \ref{section-de-rham-complex}
fournissent des morphismes d'image inverse
$$
f^* :
R\Gamma(X, \Omega^\bullet_{X/S})
\longrightarrow
R\Gamma(X', \Omega^\bullet_{X'/S'})
$$
et
$$
f^* : H^i_{dR}(X/S) \longrightarrow H^i_{dR}(X'/S')
$$
Ces images inverses satisfont à une loi de composition évidente.
En particulier, si l'on travaille sur un schéma de base fixé $S$, la cohomologie
de de Rham est un foncteur contravariant sur la catégorie des schémas sur $S$.

\begin{lemma}
\label{lemma-de-rham-affine}
Soit $X \to S$ un morphisme de schémas affines défini par le morphisme d'anneaux
$R \to A$. Alors $R\Gamma(X, \Omega^\bullet_{X/S}) = \Omega^\bullet_{A/R}$
dans $D(R)$ et $H^i_{dR}(X/S) = H^i(\Omega^\bullet_{A/R})$.
\end{lemma}

\begin{proof}
Cela résulte de Cohomologie des schémas, lemme
\ref{coherent-lemma-quasi-coherent-affine-cohomology-zero},
et du lemme d'acyclicité de Leray
(Catégories dérivées, lemme \ref{derived-lemma-leray-acyclicity}).
\end{proof}

\begin{lemma}
\label{lemma-quasi-coherence-relative}
Soit $p : X \to S$ un morphisme de schémas. Si $p$ est quasi-compact
et quasi-séparé, alors $Rp_*\Omega^\bullet_{X/S}$ est un objet
de $D_\QCoh(\mathcal{O}_S)$.
\end{lemma}

\begin{proof}
Il existe une suite spectrale de première page
$E_1^{a, b} = R^bp_*\Omega^a_{X/S}$ convergeant vers
la cohomologie de $Rp_*\Omega^\bullet_{X/S}$
(voir Catégories dérivées, lemme \ref{derived-lemma-two-ss-complex-functor}).
Ainsi, d'après Homologie, lemme \ref{homology-lemma-first-quadrant-ss},
il suffit de montrer que $R^bp_*\Omega^a_{X/S}$ est quasi-cohérent.
Cela résulte de Cohomologie des schémas, lemme
\ref{coherent-lemma-quasi-coherence-higher-direct-images}.
\end{proof}

\begin{lemma}
\label{lemma-coherence-relative}
Soit $p : X \to S$ un morphisme propre de schémas, où $S$ est localement
noethérien. Alors $Rp_*\Omega^\bullet_{X/S}$ est un objet
de $D_{\textit{Coh}}(\mathcal{O}_S)$.
\end{lemma}

\begin{proof}
Dans ce cas, d'après Morphismes, lemme \ref{morphisms-lemma-finite-type-differentials},
les modules $\Omega^i_{X/S}$ sont cohérents. On peut donc reprendre exactement
l'argument de la démonstration du lemme \ref{lemma-quasi-coherence-relative}
en utilisant Cohomologie des schémas, proposition
\ref{coherent-proposition-proper-pushforward-coherent}.
\end{proof}

\begin{lemma}
\label{lemma-finite-de-Rham}
Soit $A$ un anneau noethérien. Soit $X$ un schéma propre sur $S = \Spec(A)$.
Alors $H^i_{dR}(X/S)$ est un $A$-module de type fini pour tout $i$.
\end{lemma}

\begin{proof}
C'est un cas particulier du lemme \ref{lemma-coherence-relative}.
\end{proof}

\begin{lemma}
\label{lemma-proper-smooth-de-Rham}
Soit $f : X \to S$ un morphisme propre et lisse de schémas. Alors
$Rf_*\Omega^p_{X/S}$, pour $p \geq 0$, et $Rf_*\Omega^\bullet_{X/S}$ sont
des objets parfaits de $D(\mathcal{O}_S)$ dont la formation commute
à tout changement de base.
\end{lemma}

\begin{proof}
Puisque $f$ est lisse, les modules $\Omega^p_{X/S}$ sont des
$\mathcal{O}_X$-modules localement libres de rang fini ; voir Morphismes, lemme
\ref{morphisms-lemma-smooth-omega-finite-locally-free}. Leur
formation commute à tout changement de base d'après
le lemme \ref{lemma-base-change-de-rham}. Par conséquent,
$Rf_*\Omega^p_{X/S}$ est un objet parfait de $D(\mathcal{O}_S)$
dont la formation commute à tout changement de base ; voir
Catégories dérivées des schémas, lemme
\ref{perfect-lemma-flat-proper-perfect-direct-image-general}.
Cela démontre la première assertion du lemme.

\medskip\noindent
Pour démontrer que $Rf_*\Omega^\bullet_{X/S}$ est parfait sur $S$, on peut travailler
localement sur $S$. On peut donc supposer $S$ quasi-compact. Cela signifie
que l'on peut supposer $\Omega^n_{X/S}$ nul pour $n$ assez grand.
Pour tout $p \geq 0$, nous affirmons que
$Rf_*\sigma_{\geq p}\Omega^\bullet_{X/S}$ est un
objet parfait de $D(\mathcal{O}_S)$ dont la formation commute
à tout changement de base. D'après ce qui précède,
c'est vrai pour $p \gg 0$. Supposons l'assertion vraie pour $p$
et considérons le triangle distingué
$$
\sigma_{\geq p}\Omega^\bullet_{X/S} \to
\sigma_{\geq p - 1}\Omega^\bullet_{X/S} \to
\Omega^{p - 1}_{X/S}[-(p - 1)] \to
(\sigma_{\geq p}\Omega^\bullet_{X/S})[1]
$$
dans $D(f^{-1}\mathcal{O}_S)$.
En appliquant le foncteur exact $Rf_*$, on obtient un triangle distingué
dans $D(\mathcal{O}_S)$.
Comme la propriété d'être parfait satisfait à la propriété des deux sur trois
(Cohomologie, lemme \ref{cohomology-lemma-two-out-of-three-perfect}),
on en déduit que $Rf_*\sigma_{\geq p - 1}\Omega^\bullet_{X/S}$ est un
objet parfait de $D(\mathcal{O}_S)$. Il en va de même de la
commutation à tout changement de base.
\end{proof}





\section{Cup-produit}
\label{section-cup-product}

\noindent
Considérons les applications
$\Omega^p_{X/S} \times \Omega^q_{X/S} \to \Omega^{p + q}_{X/S}$
définies par $(\omega , \eta) \longmapsto \omega \wedge \eta$.
En utilisant la formule pour $\text{d}$ donnée à la section \ref{section-de-rham-complex}
et la règle de Leibniz pour $\text{d} : \mathcal{O}_X \to \Omega_{X/S}$,
on voit que $\text{d}(\omega \wedge \eta) = \text{d}(\omega) \wedge \eta +
(-1)^{\deg(\omega)} \omega \wedge \text{d}(\eta)$. Cela signifie que
$\wedge$ définit un morphisme
\begin{equation}
\label{equation-wedge}
\wedge :
\text{Tot}(
\Omega^\bullet_{X/S} \otimes_{p^{-1}\mathcal{O}_S} \Omega^\bullet_{X/S})
\longrightarrow
\Omega^\bullet_{X/S}
\end{equation}
de complexes de $p^{-1}\mathcal{O}_S$-modules.

\medskip\noindent
En combinant le cup-produit de
Cohomologie, section \ref{cohomology-section-cup-product},
avec (\ref{equation-wedge}), on obtient un cup-produit
$H^0(S, \mathcal{O}_S)$-bilinéaire
$$
\cup : H^i_{dR}(X/S) \times H^j_{dR}(X/S) \longrightarrow H^{i + j}_{dR}(X/S)
$$
Par exemple, si $\omega \in \Gamma(X, \Omega^i_{X/S})$ et
$\eta \in \Gamma(X, \Omega^j_{X/S})$ sont fermées,
le cup-produit des classes de cohomologie de de Rham de
$\omega$ et de $\eta$ est la classe de cohomologie de de Rham de $\omega \wedge \eta$ ;
voir la discussion dans Cohomologie, section \ref{cohomology-section-cup-product}.

\medskip\noindent
Étant donné un diagramme commutatif
$$
\xymatrix{
X' \ar[r]_f \ar[d] & X \ar[d] \\
S' \ar[r] & S
}
$$
de schémas, les morphismes d'image inverse
$f^* : R\Gamma(X, \Omega^\bullet_{X/S}) \to R\Gamma(X', \Omega^\bullet_{X'/S'})$
et
$f^* : H^i_{dR}(X/S) \longrightarrow H^i_{dR}(X'/S')$
sont compatibles avec le cup-produit défini ci-dessus.

\begin{lemma}
\label{lemma-cup-product-graded-commutative}
Soit $p : X \to S$ un morphisme de schémas.
Le cup-produit sur $H^*_{dR}(X/S)$ est associatif et
gradué commutatif.
\end{lemma}

\begin{proof}
Cela résulte de
Cohomologie, lemmes \ref{cohomology-lemma-cup-product-associative} et
\ref{cohomology-lemma-cup-product-commutative},
et du fait que $\wedge$ est associatif et gradué commutatif.
\end{proof}

\begin{remark}
\label{remark-relative-cup-product}
Soit $p : X \to S$ un morphisme de schémas. On peut considérer
$\Omega^\bullet_{X/S}$ comme un faisceau de
$p^{-1}\mathcal{O}_S$-algèbres différentielles graduées ; voir
Faisceaux différentiels gradués, définition \ref{sdga-definition-dga}.
En particulier, la discussion de
Faisceaux différentiels gradués, section \ref{sdga-section-misc},
s'applique. Par exemple, pour tout diagramme commutatif
$$
\xymatrix{
X \ar[d]_p \ar[r]_f & Y \ar[d]^q \\
S \ar[r]^h & T
}
$$
de schémas, il existe un cup-produit relatif canonique
$$
\mu :
Rf_*\Omega^\bullet_{X/S}
\otimes_{q^{-1}\mathcal{O}_T}^\mathbf{L}
Rf_*\Omega^\bullet_{X/S}
\longrightarrow
Rf_*\Omega^\bullet_{X/S}
$$
dans $D(Y, q^{-1}\mathcal{O}_T)$, qui est associatif et
qui redonne en cohomologie le cup-produit considéré ci-dessus.
\end{remark}

\begin{remark}
\label{remark-cup-product-as-a-map}
Soit $f : X \to S$ un morphisme de schémas. Soit $\xi \in H_{dR}^n(X/S)$.
D'après la discussion de
Faisceaux différentiels gradués, section \ref{sdga-section-misc},
il existe un morphisme canonique
$$
\xi' : \Omega^\bullet_{X/S} \to \Omega^\bullet_{X/S}[n]
$$
dans $D(f^{-1}\mathcal{O}_S)$, caractérisé de manière unique par
(1) et (2) dans la liste de propriétés suivante :
\begin{enumerate}
\item $\xi'$ se relève en un morphisme dans la catégorie dérivée des
$\Omega^\bullet_{X/S}$-modules différentiels gradués à droite, et
\item $\xi'(1) = \xi$ dans
$H^0(X, \Omega^\bullet_{X/S}[n]) = H^n_{dR}(X/S)$,
\item le morphisme $\xi'$ envoie $\eta \in H^m_{dR}(X/S)$
sur $\xi \cup \eta$ dans $H^{n + m}_{dR}(X/S)$,
\item la construction de $\xi'$ commute aux restrictions aux
ouverts : pour un ouvert $U \subset X$, la restriction $\xi'|_U$ est
le morphisme correspondant à l'image $\xi|_U \in H^n_{dR}(U/S)$,
\item pour tout diagramme comme dans la remarque \ref{remark-relative-cup-product},
on obtient un diagramme commutatif
$$
\xymatrix{
Rf_*\Omega^\bullet_{X/S}
\otimes_{q^{-1}\mathcal{O}_T}^\mathbf{L}
Rf_*\Omega^\bullet_{X/S} \ar[d]_{\xi' \otimes \text{id}}
\ar[r]_-\mu &
Rf_*\Omega^\bullet_{X/S} \ar[d]^{\xi'} \\
Rf_*\Omega^\bullet_{X/S}[n]
\otimes_{q^{-1}\mathcal{O}_T}^\mathbf{L}
Rf_*\Omega^\bullet_{X/S}
\ar[r]^-\mu &
Rf_*\Omega^\bullet_{X/S}[n]
}
$$
dans $D(Y, q^{-1}\mathcal{O}_T)$.
\end{enumerate}
\end{remark}




\section{Cohomologie de Hodge}
\label{section-hodge-cohomology}

\noindent
Soit $p : X \to S$ un morphisme de schémas. On appelle
{\it cohomologie de Hodge de $X$ sur $S$} les groupes de cohomologie
$$
H^n_{Hodge}(X/S) = \bigoplus\nolimits_{n = p + q} H^q(X, \Omega^p_{X/S})
$$
considérés comme un $H^0(X, \mathcal{O}_X)$-module gradué. Le produit extérieur
des formes, combiné au cup-produit de
Cohomologie, section \ref{cohomology-section-cup-product},
définit un cup-produit $H^0(X, \mathcal{O}_X)$-bilinéaire
$$
\cup :
H^i_{Hodge}(X/S) \times H^j_{Hodge}(X/S)
\longrightarrow
H^{i + j}_{Hodge}(X/S)
$$
Bien entendu, si $\xi \in H^q(X, \Omega^p_{X/S})$ et
$\xi' \in H^{q'}(X, \Omega^{p'}_{X/S})$, alors $\xi \cup \xi' \in
H^{q + q'}(X, \Omega^{p + p'}_{X/S})$.

\begin{lemma}
\label{lemma-cup-product-hodge-graded-commutative}
Soit $p : X \to S$ un morphisme de schémas.
Le cup-produit sur $H^*_{Hodge}(X/S)$ est associatif et gradué commutatif.
\end{lemma}

\begin{proof}
La démonstration est identique à celle du
lemme \ref{lemma-cup-product-graded-commutative}.
\end{proof}

\noindent
Étant donné un diagramme commutatif
$$
\xymatrix{
X' \ar[r]_f \ar[d] & X \ar[d] \\
S' \ar[r] & S
}
$$
de schémas, on dispose de morphismes d'image inverse
$f^* : H^i_{Hodge}(X/S) \longrightarrow H^i_{Hodge}(X'/S')$
compatibles avec les graduations et avec le cup-produit défini ci-dessus.







\section{Deux suites spectrales}
\label{section-hodge-to-de-rham}

\noindent
Soit $p : X \to S$ un morphisme de schémas. Puisque la catégorie
des $p^{-1}\mathcal{O}_S$-modules sur $X$ a suffisamment d'injectifs,
il existe une résolution de Cartan--Eilenberg de $\Omega^\bullet_{X/S}$.
Voir Catégories dérivées, lemme \ref{derived-lemma-cartan-eilenberg}.
On peut donc appliquer Catégories dérivées, lemme
\ref{derived-lemma-two-ss-complex-functor}, pour obtenir deux suites spectrales
convergeant toutes deux vers la cohomologie de de Rham de $X$ sur $S$.

\medskip\noindent
La première est habituellement appelée {\it suite spectrale de Hodge vers de Rham}.
Sa première page est donnée par
$$
E_1^{p, q} = H^q(X, \Omega^p_{X/S})
$$
qui sont les groupes de cohomologie de Hodge de $X/S$ (d'où son nom). La
différentielle $d_1$ de cette page est donnée par les applications
$d_1^{p, q} : H^q(X, \Omega^p_{X/S}) \to H^q(X, \Omega^{p + 1}_{X/S})$
induites par la différentielle
$\text{d} : \Omega^p_{X/S} \to \Omega^{p + 1}_{X/S}$.
Voici une représentation
$$
\xymatrix{
H^2(X, \mathcal{O}_X) \ar[r] \ar@{-->}[rrd] \ar@{..>}[rrrdd] &
H^2(X, \Omega^1_{X/S}) \ar[r] \ar@{-->}[rrd] &
H^2(X, \Omega^2_{X/S}) \ar[r] &
H^2(X, \Omega^3_{X/S}) \\
H^1(X, \mathcal{O}_X) \ar[r] \ar@{-->}[rrd] &
H^1(X, \Omega^1_{X/S}) \ar[r] \ar@{-->}[rrd] &
H^1(X, \Omega^2_{X/S}) \ar[r] &
H^1(X, \Omega^3_{X/S}) \\
H^0(X, \mathcal{O}_X) \ar[r] &
H^0(X, \Omega^1_{X/S}) \ar[r] &
H^0(X, \Omega^2_{X/S}) \ar[r] &
H^0(X, \Omega^3_{X/S})
}
$$
où les flèches en tirets indiquent la source et le but
des différentielles de la page $E_2$, et la flèche en pointillés une différentielle
de la page $E_3$. En considérant le degré $0$, on en déduit que
$$
H^0_{dR}(X/S) =
\Ker(\text{d} : H^0(X, \mathcal{O}_X) \to H^0(X, \Omega^1_{X/S}))
$$
Bien entendu, cela résulte aussi immédiatement du fait que le
complexe de de Rham commence en degré $0$ par $\mathcal{O}_X \to \Omega^1_{X/S}$.

\medskip\noindent
La seconde suite spectrale est habituellement appelée
{\it suite spectrale conjuguée}. Sa deuxième page est donnée par
$$
E_2^{p, q} = H^p(X, H^q(\Omega^\bullet_{X/S})) = H^p(X, \mathcal{H}^q)
$$
où $\mathcal{H}^q = H^q(\Omega^\bullet_{X/S})$ est le faisceau de cohomologie de degré $q$
du complexe de de Rham de $X/S$. Les différentielles
de cette page sont données par $E_2^{p, q} \to E_2^{p + 2, q - 1}$.
Voici une représentation
$$
\xymatrix{
H^0(X, \mathcal{H}^2) \ar[rrd] \ar@{..>}[rrrdd] &
H^1(X, \mathcal{H}^2) \ar[rrd] &
H^2(X, \mathcal{H}^2) &
H^3(X, \mathcal{H}^2) \\
H^0(X, \mathcal{H}^1) \ar[rrd] &
H^1(X, \mathcal{H}^1) \ar[rrd] &
H^2(X, \mathcal{H}^1) &
H^3(X, \mathcal{H}^1) \\
H^0(X, \mathcal{H}^0) &
H^1(X, \mathcal{H}^0) &
H^2(X, \mathcal{H}^0) &
H^3(X, \mathcal{H}^0)
}
$$
En considérant le degré $0$, on en déduit que
$$
H^0_{dR}(X/S) = H^0(X, \mathcal{H}^0)
$$
ce qui est évident à la réflexion. En degré $1$, on obtient une suite exacte
$$
0 \to H^1(X, \mathcal{H}^0) \to H^1_{dR}(X/S) \to
H^0(X, \mathcal{H}^1) \to H^2(X, \mathcal{H}^0) \to H^2_{dR}(X/S)
$$
Il se trouve que, si $X \to S$ est lisse et si $S$ est en caractéristique $p$,
les faisceaux $\mathcal{H}^q$ peuvent être calculés (en termes de certains
faisceaux de différentielles), et la suite spectrale conjuguée est un outil
précieux (insérer ici une référence ultérieure).



\section{La filtration de Hodge}
\label{section-hodge-filtration}

\noindent
Soit $X \to S$ un morphisme de schémas. La filtration de Hodge sur $H^n_{dR}(X/S)$
est la filtration induite par la suite spectrale de Hodge vers de Rham
(Homologie, définition
\ref{homology-definition-filtration-cohomology-filtered-complex}).
Pour éviter toute ambiguïté, nous la définissons explicitement comme suit.

\begin{definition}
\label{definition-hodge-filtration}
Soit $X \to S$ un morphisme de schémas. La {\it filtration de Hodge}
sur $H^n_{dR}(X/S)$ est la filtration dont les termes sont
$$
F^pH^n_{dR}(X/S) = \Im\left(H^n(X, \sigma_{\geq p}\Omega^\bullet_{X/S})
\longrightarrow H^n_{dR}(X/S)\right)
$$
où $\sigma_{\geq p}\Omega^\bullet_{X/S}$ est défini dans
Homologie, section \ref{homology-section-truncations}.
\end{definition}

\noindent
Bien entendu, $\sigma_{\geq p}\Omega^\bullet_{X/S}$ est un sous-complexe du
complexe de de Rham relatif, et l'on obtient une filtration
$$
\Omega^\bullet_{X/S} = \sigma_{\geq 0}\Omega^\bullet_{X/S} \supset
\sigma_{\geq 1}\Omega^\bullet_{X/S} \supset
\sigma_{\geq 2}\Omega^\bullet_{X/S} \supset
\sigma_{\geq 3}\Omega^\bullet_{X/S} \supset \ldots
$$
du complexe de de Rham relatif, avec
$\text{gr}^p(\Omega^\bullet_{X/S}) = \Omega^p_{X/S}[-p]$.
La suite spectrale construite dans
Cohomologie, lemme \ref{cohomology-lemma-spectral-sequence-filtered-object},
pour $\Omega^\bullet_{X/S}$ considéré comme un complexe filtré de faisceaux
coïncide avec la suite spectrale de Hodge vers de Rham construite à la
section \ref{section-hodge-to-de-rham}, d'après
Cohomologie, exemple \ref{cohomology-example-spectral-sequence-bis}. De plus, le
produit extérieur (\ref{equation-wedge}) envoie
$\text{Tot}(\sigma_{\geq i}\Omega^\bullet_{X/S} \otimes_{p^{-1}\mathcal{O}_S}
\sigma_{\geq j}\Omega^\bullet_{X/S})$ dans
$\sigma_{\geq i + j}\Omega^\bullet_{X/S}$. On obtient donc des
diagrammes commutatifs
$$
\xymatrix{
H^n(X, \sigma_{\geq i}\Omega^\bullet_{X/S}))
\times 
H^m(X, \sigma_{\geq j}\Omega^\bullet_{X/S}))
\ar[r] \ar[d] &
H^{n + m}(X, \sigma_{\geq i + j}\Omega^\bullet_{X/S})) \ar[d] \\
H^n_{dR}(X/S) \times
H^m_{dR}(X/S)
\ar[r]^\cup &
H^{n + m}_{dR}(X/S)
}
$$
En particulier, on trouve que
$$
F^iH^n_{dR}(X/S) \cup F^jH^m_{dR}(X/S) \subset F^{i + j}H^{n + m}_{dR}(X/S)
$$






\section{Formule de K\"unneth}
\label{section-kunneth}

\noindent
Une propriété importante de la cohomologie de de Rham est l'existence
d'une formule de K\"unneth.

\medskip\noindent
Soient $a : X \to S$ et $b : Y \to S$ des morphismes de schémas de même
but. Soient $p : X \times_S Y \to X$ et $q : X \times_S Y \to Y$ les
projections, et posons $f = a \circ p = b \circ q$. Voici le diagramme correspondant
$$
\xymatrix{
& X \times_S Y \ar[ld]^p \ar[rd]_q \ar[dd]^f \\
X \ar[rd]_a & & Y \ar[ld]^b \\
& S
}
$$
Dans cette section, pour un $\mathcal{O}_X$-module $\mathcal{F}$
et un $\mathcal{O}_Y$-module $\mathcal{G}$, posons
$$
\mathcal{F} \boxtimes \mathcal{G} =
p^*\mathcal{F} \otimes_{\mathcal{O}_{X \times_S Y}} q^*\mathcal{G}
$$
Le bifoncteur
$(\mathcal{F}, \mathcal{G}) \mapsto \mathcal{F} \boxtimes \mathcal{G}$
sur les modules quasi-cohérents s'étend en un bifoncteur sur les modules quasi-cohérents
et les opérateurs différentiels d'ordre fini sur $S$ ; voir
Morphismes, remarque \ref{morphisms-remark-base-change-differential-operators}.
Les différentielles des complexes de de Rham $\Omega^\bullet_{X/S}$ et
$\Omega^\bullet_{Y/S}$ sont des opérateurs différentiels d'ordre $1$
sur $S$, d'après Modules, lemme
\ref{modules-lemma-differentials-relative-de-rham-complex-order-1}.
Il est donc possible de considérer le complexe
$$
\text{Tot}(\Omega^\bullet_{X/S} \boxtimes \Omega^\bullet_{Y/S})
$$
Voir la discussion dans Catégories dérivées des schémas, section
\ref{perfect-section-kunneth-complexes}.

\begin{lemma}
\label{lemma-de-rham-complex-product}
Dans la situation ci-dessus, il existe un isomorphisme canonique
$$
\text{Tot}(\Omega^\bullet_{X/S} \boxtimes \Omega^\bullet_{Y/S})
\longrightarrow
\Omega^\bullet_{X \times_S Y/S}
$$
de complexes de $f^{-1}\mathcal{O}_S$-modules.
\end{lemma}

\begin{proof}
On sait que
$
\Omega_{X \times_S Y/S} = p^*\Omega_{X/S} \oplus q^*\Omega_{Y/S}
$
d'après Morphismes, lemme \ref{morphisms-lemma-differential-product}.
En prenant les puissances extérieures, on obtient
$$
\Omega^n_{X \times_S Y/S} =
\bigoplus\nolimits_{i + j = n}
p^*\Omega^i_{X/S} \otimes_{\mathcal{O}_{X \times_S Y}} q^*\Omega^j_{Y/S} =
\bigoplus\nolimits_{i + j = n}
\Omega^i_{X/S} \boxtimes \Omega^j_{Y/S}
$$
par les propriétés élémentaires des puissances extérieures. Ces identifications déterminent
des isomorphismes entre les termes des complexes à gauche et à droite
de la flèche du lemme. Nous omettons de vérifier que ces morphismes
sont compatibles avec les différentielles.
\end{proof}

\noindent
Posons $A = \Gamma(S, \mathcal{O}_S)$. En combinant le résultat du
lemme \ref{lemma-de-rham-complex-product} avec le morphisme de
Catégories dérivées des schémas, équation
(\ref{perfect-equation-de-rham-kunneth}),
on obtient un cup-produit
$$
R\Gamma(X, \Omega^\bullet_{X/S})
\otimes_A^\mathbf{L}
R\Gamma(Y, \Omega^\bullet_{Y/S})
\longrightarrow
R\Gamma(X \times_S Y, \Omega^\bullet_{X \times_S Y/S})
$$
Au niveau de la cohomologie, en utilisant la discussion de
Compléments d'algèbre, section \ref{more-algebra-section-products-tor},
on obtient une application canonique
$$
H^i_{dR}(X/S) \otimes_A H^j_{dR}(Y/S)
\longrightarrow
H^{i + j}_{dR}(X \times_S Y/S),\quad
(\xi, \zeta) \longmapsto p^*\xi \cup q^*\zeta
$$
Notons que la construction ci-dessus consiste bien à
prendre d'abord les images inverses, puis leur cup-produit.

\begin{lemma}
\label{lemma-kunneth-de-rham}
Supposons $X$ et $Y$ lisses, quasi-compacts et à diagonale affine sur
$S = \Spec(A)$. Alors le morphisme
$$
R\Gamma(X, \Omega^\bullet_{X/S})
\otimes_A^\mathbf{L}
R\Gamma(Y, \Omega^\bullet_{Y/S})
\longrightarrow
R\Gamma(X \times_S Y, \Omega^\bullet_{X \times_S Y/S})
$$
est un isomorphisme dans $D(A)$.
\end{lemma}

\begin{proof}
D'après Morphismes, lemme \ref{morphisms-lemma-smooth-omega-finite-locally-free},
les faisceaux $\Omega^n_{X/S}$ et $\Omega^m_{Y/S}$ sont respectivement des
$\mathcal{O}_X$-modules et des $\mathcal{O}_Y$-modules localement libres de rang fini. D'autre part, $X$ et $Y$
sont plats sur $S$ (Morphismes, lemme \ref{morphisms-lemma-smooth-flat}),
d'où la platitude de $\Omega^n_{X/S}$ et de $\Omega^m_{Y/S}$ sur $S$.
Observons aussi que $\Omega^\bullet_{X/S}$ est localement borné. Le résultat
découle donc du lemme \ref{lemma-de-rham-complex-product} et de
Catégories dérivées des schémas, lemme \ref{perfect-lemma-kunneth-special}.
\end{proof}

\noindent
Il existe une version relative du cup-produit, à savoir un morphisme
$$
Ra_*\Omega^\bullet_{X/S}
\otimes_{\mathcal{O}_S}^\mathbf{L}
Rb_*\Omega^\bullet_{Y/S}
\longrightarrow
Rf_*\Omega^\bullet_{X \times_S Y/S}
$$
dans $D(\mathcal{O}_S)$. Sa construction combine
le lemme \ref{lemma-de-rham-complex-product} avec le morphisme de
Catégories dérivées des schémas, équation
(\ref{perfect-equation-relative-de-rham-kunneth}).
La construction montre que ce morphisme est donné par le diagramme
$$
\xymatrix{
Ra_*\Omega^\bullet_{X/S}
\otimes_{\mathcal{O}_S}^\mathbf{L}
Rb_*\Omega^\bullet_{Y/S}
\ar[d]^{\text{units of adjunction}}  \\
Rf_*(p^{-1}\Omega^\bullet_{X/S})
\otimes_{\mathcal{O}_S}^\mathbf{L}
Rf_*(q^{-1}\Omega^\bullet_{Y/S}) \ar[r] \ar[d]^{\text{relative cup product}} &
Rf_*(\Omega^\bullet_{X \times_S Y/S})
\otimes_{\mathcal{O}_S}^\mathbf{L}
Rf_*(\Omega^\bullet_{X \times_S Y/S}) \ar[d]^{\text{relative cup product}} \\
Rf_*(p^{-1}\Omega^\bullet_{X/S}
\otimes_{f^{-1}\mathcal{O}_S}^\mathbf{L}
q^{-1}\Omega^\bullet_{Y/S})
\ar[d]^{\text{from derived to usual}} \ar[r] &
Rf_*(\Omega^\bullet_{X \times_S Y/S}
\otimes_{f^{-1}\mathcal{O}_S}^\mathbf{L}
\Omega^\bullet_{X \times_S Y/S})
\ar[d]^{\text{from derived to usual}} \\
Rf_*\text{Tot}(p^{-1}\Omega^\bullet_{X/S}
\otimes_{f^{-1}\mathcal{O}_S}
q^{-1}\Omega^\bullet_{Y/S}) \ar[r] \ar[d]^{\text{canonical map}} &
Rf_*\text{Tot}(\Omega^\bullet_{X \times_S Y/S}
\otimes_{f^{-1}\mathcal{O}_S}
\Omega^\bullet_{X \times_S Y/S})
\ar[d]^{\eta \otimes \omega \mapsto \eta \wedge \omega} \\
Rf_*\text{Tot}(\Omega^\bullet_{X/S} \boxtimes \Omega^\bullet_{Y/S})
\ar@{=}[r]
&
Rf_*\Omega^\bullet_{X \times_S Y/S}
}
$$
Ici, la première flèche utilise les unités $\text{id} \to Rp_* p^{-1}$
et $\text{id} \to Rq_* q^{-1}$ des adjonctions, ainsi que les
identifications $Rf_* p^{-1} = Ra_* Rp_* p^{-1}$ et
$Rf_* q^{-1} = Rb_* Rq_* q^{-1}$.
La deuxième flèche est le cup-produit relatif de
Cohomologie, remarque \ref{cohomology-remark-cup-product}.
La troisième flèche est le morphisme du produit tensoriel dérivé
de complexes vers le complexe total de leur produit tensoriel.
L'égalité finale est le lemme \ref{lemma-de-rham-complex-product}.
Sur les sections globales, cette construction redonne celle présentée plus haut.

\begin{lemma}
\label{lemma-kunneth-de-rham-relative}
Supposons $X \to S$ et $Y \to S$ lisses et quasi-compacts,
et les morphismes $X \to X \times_S X$ et $Y \to Y \times_S Y$ affines.
Alors le cup-produit relatif
$$
Ra_*\Omega^\bullet_{X/S}
\otimes_{\mathcal{O}_S}^\mathbf{L}
Rb_*\Omega^\bullet_{Y/S}
\longrightarrow
Rf_*\Omega^\bullet_{X \times_S Y/S}
$$
est un isomorphisme dans $D(\mathcal{O}_S)$.
\end{lemma}

\begin{proof}
Conséquence immédiate du lemme \ref{lemma-kunneth-de-rham}.
\end{proof}







\section{Première classe de Chern en cohomologie de de Rham}
\label{section-first-chern-class}

\noindent
Soit $X \to S$ un morphisme de schémas. Il existe un morphisme de complexes
$$
\text{d}\log : \mathcal{O}_X^*[-1] \longrightarrow \Omega^\bullet_{X/S}
$$
qui envoie la section $g \in \mathcal{O}_X^*(U)$ sur la section
$\text{d}\log(g) = g^{-1}\text{d}g$ de $\Omega^1_{X/S}(U)$.
On peut donc considérer l'application
$$
\Pic(X) = H^1(X, \mathcal{O}_X^*) =
H^2(X, \mathcal{O}_X^*[-1]) \longrightarrow H^2_{dR}(X/S)
$$
où la première égalité est donnée par
Cohomologie, lemme \ref{cohomology-lemma-h1-invertible}.
L'image de la classe d'isomorphisme du module inversible
$\mathcal{L}$ est notée $c^{dR}_1(\mathcal{L}) \in H^2_{dR}(X/S)$.

\medskip\noindent
On peut également utiliser le morphisme $\text{d}\log : \mathcal{O}_X^* \to \Omega^1_{X/S}$
pour définir une classe de Chern en cohomologie de Hodge
$$
c_1^{Hodge} : \Pic(X) \longrightarrow H^1(X, \Omega^1_{X/S})
\subset H^2_{Hodge}(X/S)
$$
Ces constructions sont compatibles avec les images inverses.

\begin{lemma}
\label{lemma-pullback-c1}
Étant donné un diagramme commutatif
$$
\xymatrix{
X' \ar[r]_f \ar[d] & X \ar[d] \\
S' \ar[r] & S
}
$$
de schémas, les diagrammes
$$
\xymatrix{
\Pic(X') \ar[d]_{c_1^{dR}} &
\Pic(X) \ar[d]^{c_1^{dR}} \ar[l]^{f^*} \\
H^2_{dR}(X'/S') &
H^2_{dR}(X/S) \ar[l]_{f^*}
}
\quad
\xymatrix{
\Pic(X') \ar[d]_{c_1^{Hodge}} &
\Pic(X) \ar[d]^{c_1^{Hodge}} \ar[l]^{f^*} \\
H^1(X', \Omega^1_{X'/S'}) &
H^1(X, \Omega^1_{X/S}) \ar[l]_{f^*}
}
$$
sont commutatifs.
\end{lemma}

\begin{proof}
Démonstration omise.
\end{proof}

\noindent
« Calculons » l'élément $c^{dR}_1(\mathcal{L})$ en cohomologie de {\v C}ech
(avec les conventions de signe pour les différentielles de {\v C}ech
de Cohomologie, section
\ref{cohomology-section-cech-cohomology-of-complexes}).
Choisissons un recouvrement ouvert
$\mathcal{U} : X = \bigcup_{i \in I} U_i$ tel que
l'on dispose d'une section trivialisante $s_i$ de $\mathcal{L}|_{U_i}$ pour tout $i$.
Sur les intersections $U_{i_0i_1} = U_{i_0} \cap U_{i_1}$,
on dispose d'une fonction inversible $f_{i_0i_1}$ telle que
$f_{i_0i_1} = s_{i_1}|_{U_{i_0i_1}} s_{i_0}|_{U_{i_0i_1}}^{-1}$\footnote{La
différentielle de {\v C}ech d'un cycle de degré $0$, $\{a_{i_0}\}$,
vaut $a_{i_1} - a_{i_0}$ sur $U_{i_0i_1}$.}.
Bien entendu, on a
$$
f_{i_1i_2}|_{U_{i_0i_1i_2}}
f_{i_0i_2}^{-1}|_{U_{i_0i_1i_2}}
f_{i_0i_1}|_{U_{i_0i_1i_2}} = 1
$$
La classe de cohomologie de $\mathcal{L}$ dans $H^1(X, \mathcal{O}_X^*)$ est
l'image de la classe de cohomologie de {\v C}ech du cocycle $\{f_{i_0i_1}\}$ dans
$\check{\mathcal{C}}^\bullet(\mathcal{U}, \mathcal{O}_X^*)$.
On voit donc que $c_1^{dR}(\mathcal{L})$ est l'image
de la classe de cohomologie associée au cocycle de {\v C}ech
$\{\alpha_{i_0 \ldots i_p}\}$ dans
$\text{Tot}(\check{\mathcal{C}}^\bullet(\mathcal{U}, \Omega_{X/S}^\bullet))$,
de degré $2$, donné par
\begin{enumerate}
\item $\alpha_{i_0} = 0$ dans $\Omega^2_{X/S}(U_{i_0})$,
\item $\alpha_{i_0i_1} = f_{i_0i_1}^{-1}\text{d}f_{i_0i_1}$ dans
$\Omega^1_{X/S}(U_{i_0i_1})$, et
\item $\alpha_{i_0i_1i_2} = 0$ dans $\mathcal{O}_{X/S}(U_{i_0i_1i_2})$.
\end{enumerate}
Supposons donnés des modules inversibles $\mathcal{L}_k$, $k = 1, \ldots, a$,
chacun trivialisé sur $U_i$ pour tout $i \in I$, donnant des cocycles
$f_{k, i_0i_1}$ et $\alpha_k = \{\alpha_{k, i_0 \ldots i_p}\}$ comme ci-dessus.
En utilisant la règle de
Cohomologie, section \ref{cohomology-section-cech-cohomology-of-complexes},
on calcule
$$
\beta = \alpha_1 \cup \alpha_2 \cup \ldots \cup \alpha_a
$$
comme le cocycle $\beta = \{\beta_{i_0 \ldots i_p}\}$
décrit comme suit :
\begin{enumerate}
\item $\beta_{i_0 \ldots i_p} = 0$ dans
$\Omega^{2a - p}_{X/S}(U_{i_0 \ldots i_p})$ sauf si $p = a$, et
\item $\beta_{i_0 \ldots i_a} = (-1)^{a(a - 1)/2}
\alpha_{1, i_0i_1} \wedge \alpha_{2, i_1 i_2} \wedge \ldots \wedge
\alpha_{a, i_{a - 1}i_a}$ dans
$\Omega^a_{X/S}(U_{i_0 \ldots i_a})$.
\end{enumerate}
C'est donc un cocycle représentant
$c_1^{dR}(\mathcal{L}_1) \cup \ldots \cup c_1^{dR}(\mathcal{L}_a)$.
Bien entendu, le même calcul montre que le cocycle
$\{\beta_{i_0 \ldots i_a}\}$ dans
$\check{\mathcal{C}}^a(\mathcal{U}, \Omega_{X/S}^a))$
représente la classe de cohomologie
$c_1^{Hodge}(\mathcal{L}_1) \cup \ldots \cup c_1^{Hodge}(\mathcal{L}_a)$

\begin{remark}
\label{remark-truncations}
Voici une reformulation des calculs ci-dessus en termes plus abstraits.
Soit $p : X \to S$ un morphisme de schémas. Soit $\mathcal{L}$ un
$\mathcal{O}_X$-module inversible. Si l'on considère $\text{d}\log$ comme un morphisme
$$
\mathcal{O}_X^*[-1] \to \sigma_{\geq 1}\Omega^\bullet_{X/S}
$$,
l'identification $\Pic(X) = H^1(X, \mathcal{O}_X^*)$ utilisée plus haut fournit une
classe de cohomologie
$$
\gamma_1(\mathcal{L}) \in H^2(X, \sigma_{\geq 1}\Omega^\bullet_{X/S})
$$
L'image de $\gamma_1(\mathcal{L})$ par le morphisme
$\sigma_{\geq 1}\Omega^\bullet_{X/S} \to \Omega^\bullet_{X/S}$
redonne $c_1^{dR}(\mathcal{L})$. En particulier, on voit que
$c_1^{dR}(\mathcal{L}) \in F^1H^2_{dR}(X/S)$ ; voir la
section \ref{section-hodge-filtration}. L'image de $\gamma_1(\mathcal{L})$
par le morphisme $\sigma_{\geq 1}\Omega^\bullet_{X/S} \to \Omega^1_{X/S}[-1]$
redonne $c_1^{Hodge}(\mathcal{L})$. En prenant le cup-produit
(voir section \ref{section-hodge-filtration}), on obtient
$$
\xi = \gamma_1(\mathcal{L}_1) \cup \ldots \cup \gamma_1(\mathcal{L}_a) \in
H^{2a}(X, \sigma_{\geq a}\Omega^\bullet_{X/S})
$$
Les diagrammes commutatifs de la section \ref{section-hodge-filtration}
montrent que $\xi$ a pour image
$c_1^{dR}(\mathcal{L}_1) \cup \ldots \cup c_1^{dR}(\mathcal{L}_a)$
dans $H^{2a}_{dR}(X/S)$ par le morphisme
$\sigma_{\geq a}\Omega^\bullet_{X/S} \to \Omega^\bullet_{X/S}$.
Il en résulte aussi que
$c_1^{dR}(\mathcal{L}_1) \cup \ldots \cup c_1^{dR}(\mathcal{L}_a)$
appartient à $F^a H^{2a}_{dR}(X/S)$. De même, le morphisme
$\sigma_{\geq a}\Omega^\bullet_{X/S} \to \Omega^a_{X/S}[-a]$
envoie $\xi$ sur
$c_1^{Hodge}(\mathcal{L}_1) \cup \ldots \cup c_1^{Hodge}(\mathcal{L}_a)$
dans $H^a(X, \Omega^a_{X/S})$.
\end{remark}

\begin{remark}
\label{remark-log-forms}
Soit $p : X \to S$ un morphisme de schémas. Pour $i > 0$,
notons $\Omega^i_{X/S, log} \subset \Omega^i_{X/S}$ le sous-faisceau abélien
engendré par les sections locales de la forme
$$
\text{d}\log(u_1) \wedge \ldots \wedge \text{d}\log(u_i)
$$
où $u_1, \ldots, u_n$ sont des sections locales inversibles de $\mathcal{O}_X$.
Pour $i = 0$, le sous-faisceau $\Omega^0_{X/S, log} \subset \mathcal{O}_X$
est l'image de $\mathbf{Z} \to \mathcal{O}_X$. Pour tout $i \geq 0$,
on dispose d'un morphisme de complexes
$$
\Omega^i_{X/S, log}[-i] \longrightarrow \Omega^\bullet_{X/S}
$$
car la différentielle d'une forme logarithmique est nulle. De plus, le produit extérieur
de formes logarithmiques est encore logarithmique ; on obtient donc des applications bilinéaires
$$
\wedge :  \Omega^i_{X/S, log} \times
\Omega^j_{X/S, log} \longrightarrow \Omega^{i + j}_{X/S, log}
$$
compatibles avec (\ref{equation-wedge}) et avec les morphismes ci-dessus.
Soit $\mathcal{L}$ un $\mathcal{O}_X$-module inversible.
En utilisant le morphisme de faisceaux abéliens
$\text{d}\log : \mathcal{O}_X^* \to \Omega^1_{X/S, log}$
et l'identification $\Pic(X) = H^1(X, \mathcal{O}_X^*)$,
on obtient une classe de cohomologie canonique
$$
\tilde \gamma_1(\mathcal{L}) \in H^1(X, \Omega^1_{X/S, log})
$$
Ces classes possèdent les propriétés suivantes :
\begin{enumerate}
\item l'image de $\tilde \gamma_1(\mathcal{L})$ par le morphisme canonique
$\Omega^1_{X/S, log}[-1] \to \sigma_{\geq 1}\Omega^\bullet_{X/S}$
s'obtient en envoyant $\tilde \gamma_1(\mathcal{L})$ sur la classe
$\gamma_1(\mathcal{L}) \in 
H^2(X, \sigma_{\geq 1}\Omega^\bullet_{X/S})$
de la remarque \ref{remark-truncations},
\item l'image de $\tilde \gamma_1(\mathcal{L})$ par le morphisme canonique
$\Omega^1_{X/S, log}[-1] \to \Omega^\bullet_{X/S}$
s'obtient en envoyant $\tilde \gamma_1(\mathcal{L})$ sur $c_1^{dR}(\mathcal{L})$ dans
$H^2_{dR}(X/S)$,
\item l'image de $\tilde \gamma_1(\mathcal{L})$ par le morphisme canonique
$\Omega^1_{X/S, log} \to \Omega^1_{X/S}$
s'obtient en envoyant $\tilde \gamma_1(\mathcal{L})$ sur $c_1^{Hodge}(\mathcal{L})$ dans
$H^1(X, \Omega^1_{X/S})$,
\item la construction de ces classes est compatible avec les images inverses,
\item ajouter d'autres propriétés ici.
\end{enumerate}
\end{remark}





\section{Cohomologie de de Rham d'un fibré en droites}
\label{section-line-bundle}

\noindent
Un fibré en droites est un cas particulier de fibré vectoriel, lequel est
un cône muni d'une structure supplémentaire. Pour étudier convenablement
leur complexe de de Rham, il convient d'examiner le complexe de de Rham
d'un anneau gradué.

\begin{remark}[Complexe de de Rham d'un anneau gradué]
\label{remark-de-rham-complex-graded}
Soit $G$ un monoïde abélien noté additivement, d'élément neutre $0$.
Soit $R \to A$ un morphisme d'anneaux, et supposons $A$ muni d'une graduation
$A = \bigoplus_{g \in G} A_g$ par des $R$-modules telle que $R$ s'envoie dans $A_0$
et que $A_g \cdot A_{g'} \subset A_{g + g'}$. Le module des différentielles
est alors muni d'une graduation
$$
\Omega_{A/R} = \bigoplus\nolimits_{g \in G} \Omega_{A/R, g}
$$
où $\Omega_{A/R, g}$ est le $R$-sous-module de $\Omega_{A/R}$
engendré par les $a_0 \text{d}a_1$ avec $a_i \in A_{g_i}$ tels que
$g = g_0 + g_1$. De même, on obtient
$$
\Omega^p_{A/R} = \bigoplus\nolimits_{g \in G} \Omega^p_{A/R, g}
$$
où $\Omega^p_{A/R, g}$ est le $R$-sous-module de $\Omega^p_{A/R}$
engendré par les $a_0 \text{d}a_1 \wedge \ldots \wedge \text{d}a_p$
avec $a_i \in A_{g_i}$ tels que $g = g_0 + g_1 + \ldots + g_p$.
Bien entendu, les différentielles respectent la graduation et le produit
extérieur est compatible avec les graduations de manière évidente.
\end{remark}

\noindent
Soit $f : X \to S$ un morphisme de schémas. Soit $\pi : C \to X$ un cône ; voir
Constructions, définition \ref{constructions-definition-abstract-cone}.
Rappelons que cela signifie que $\pi$ est affine et que l'on dispose d'une graduation
$\pi_*\mathcal{O}_C = \bigoplus_{n \geq 0} \mathcal{A}_n$ avec
$\mathcal{A}_0 = \mathcal{O}_X$.
En utilisant la discussion de la remarque \ref{remark-de-rham-complex-graded}
sur les ouverts affines, on trouve que\footnote{Avec nos excuses pour la notation !}
$$
\pi_*(\Omega^\bullet_{C/S}) =
\bigoplus\nolimits_{n \geq 0} \Omega^\bullet_{C/S, n}
$$
est canoniquement une somme directe de sous-complexes. De plus, on dispose d'une factorisation
$$
\Omega^\bullet_{X/S} \to \Omega^\bullet_{C/S, 0} \to
\pi_*(\Omega^\bullet_{C/S})
$$
et l'on sait que $\omega \wedge \eta \in \Omega^{p + q}_{C/S, n + m}$
si $\omega \in \Omega^p_{C/S, n}$ et $\eta \in \Omega^q_{C/S, m}$.

\medskip\noindent
Soit $f : X \to S$ un morphisme de schémas. Soit $\pi : L \to X$ le
fibré en droites associé au $\mathcal{O}_X$-module inversible $\mathcal{L}$.
Cela signifie que $\pi$ est l'unique morphisme affine tel que
$$
\pi_*\mathcal{O}_L = 
\bigoplus\nolimits_{n \geq 0} \mathcal{L}^{\otimes n}
$$
comme $\mathcal{O}_X$-algèbres. Ainsi, $L$ est un cône sur $X$.
La discussion ci-dessus fournit une
décomposition canonique en somme directe
$$
\pi_*(\Omega^\bullet_{L/S}) =
\bigoplus\nolimits_{n \geq 0} \Omega^\bullet_{L/S, n}
$$
compatible avec le produit extérieur et avec la décomposition
de $\pi_*\mathcal{O}_L$ ci-dessus, et telle que $\Omega_{X/S}$
s'envoie dans la composante $\Omega_{L/S, 0}$ de degré $0$.

\medskip\noindent
Un autre cas nous sera utile. Considérons en effet le
complémentaire\footnote{Le schéma $L^\star$ est le $\mathbf{G}_m$-torseur
sur $X$ associé à $L$. C'est pourquoi la graduation obtenue ci-dessous est
une $\mathbf{Z}$-graduation ; comparer avec Groupoïdes,
exemple \ref{groupoids-example-Gm-on-affine} et
lemmes \ref{groupoids-lemma-complete-reducibility-Gm} et
\ref{groupoids-lemma-Gm-equivariant-module}.}
$L^\star \subset L$ de la section nulle $o : X \to L$ dans notre fibré
en droites $L$. Un calcul local donne un isomorphisme canonique
$$
(L^\star \to X)_*\mathcal{O}_{L^\star} =
\bigoplus\nolimits_{n \in \mathbf{Z}} \mathcal{L}^{\otimes n}
$$
de $\mathcal{O}_X$-algèbres. Le membre de droite est une algèbre $\mathbf{Z}$-graduée
quasi-cohérente sur $\mathcal{O}_X$. En utilisant la discussion de la
remarque \ref{remark-de-rham-complex-graded} sur les ouverts affines, on obtient
$$
(L^\star \to X)_*(\Omega^\bullet_{L^\star/S}) =
\bigoplus\nolimits_{n \in \mathbf{Z}} \Omega^\bullet_{L^\star/S, n}
$$
de manière compatible avec le produit extérieur et avec la décomposition
de $(L^\star \to X)_*\mathcal{O}_{L^\star}$ ci-dessus, et telle que
$\Omega_{X/S}$ s'envoie dans la composante $\Omega_{L^\star/S, 0}$ de degré $0$.
Le complexe $\Omega^\bullet_{L^\star/S, 0}$ nous intéressera
tout particulièrement.

\begin{lemma}
\label{lemma-the-complex-for-L-star}
Avec les notations ci-dessus, il existe une suite exacte courte de complexes
$$
0 \to \Omega^\bullet_{X/S} \to
\Omega^\bullet_{L^\star/S, 0} \to
\Omega^\bullet_{X/S}[-1] \to 0
$$
\end{lemma}

\begin{proof}
Nous avons construit le morphisme
$\Omega^\bullet_{X/S} \to \Omega^\bullet_{L^\star/S, 0}$ ci-dessus.

\medskip\noindent
Construction de
$\text{Res} : \Omega^\bullet_{L^\star/S, 0} \to \Omega^\bullet_{X/S}[-1]$.
Soit $U \subset X$ un ouvert, et soient $s \in \mathcal{L}(U)$
et $s' \in \mathcal{L}^{\otimes -1}(U)$ des sections telles que
$s' s = 1$. Alors $s$ donne une section inversible du faisceau
d'algèbres $(L^\star \to X)_*\mathcal{O}_{L^\star}$ sur $U$,
d'inverse $s' = s^{-1}$. On peut alors considérer la forme de degré $1$
$\text{d}\log(s) = s' \text{d}(s)$, qui appartient à
$\Omega^1_{L^\star/S, 0}(U)$ d'après notre construction de la graduation sur
$\Omega^1_{L^\star/S}$. Les calculs sur les ouverts affines effectués ci-dessous
montreront que $1$ et $\text{d}\log(s)$ engendrent librement
$\Omega^\bullet_{L^\star/S, 0}|_U$ comme module à droite sur
$\Omega^\bullet_{X/S}|_U$.
On peut donc définir $\text{Res}$ sur $U$ par la règle
$$
\text{Res}(\omega' + \text{d}\log(s) \wedge \omega) = \omega
$$
pour tous $\omega', \omega \in \Omega^\bullet_{X/S}(U)$. Ce
morphisme est indépendant du choix du générateur local $s$ ; il se
recolle donc en un morphisme global. En effet, un autre choix de $s$
serait de la forme $gs$ pour un certain élément inversible $g \in \mathcal{O}_X(U)$,
et l'on aurait $\text{d}\log(gs) = g^{-1}\text{d}(g) + \text{d}\log(s)$,
d'où l'indépendance résulte aisément.
Enfin, observons que notre règle définissant $\text{Res}$
est compatible avec les différentielles,
puisque $\text{d}(\omega' + \text{d}\log(s) \wedge \omega) =
\text{d}(\omega') - \text{d}\log(s) \wedge \text{d}(\omega)$
et que la différentielle sur $\Omega^\bullet_{X/S}[-1]$
envoie $\omega'$ sur $-\text{d}(\omega')$ selon notre convention de signe dans
Homologie, définition \ref{homology-definition-shift-cochain}.

\medskip\noindent
Calcul local. On peut recouvrir $X$ par des ouverts affines $U \subset X$
tels que $\mathcal{L}|_U \cong \mathcal{O}_U$, qui s'envoient en outre
dans un ouvert affine $V \subset S$. Écrivons $U = \Spec(A)$, $V = \Spec(R)$,
et choisissons un générateur $s$ de $\mathcal{L}$. On obtient
$$
L^\star \times_X U = \Spec(A[s, s^{-1}])
$$
Le calcul des différentielles donne
$$
\Omega^1_{A[s, s^{-1}]/R} =
A[s, s^{-1}] \otimes_A \Omega^1_{A/R} \oplus A[s, s^{-1}] \text{d}\log(s)
$$
et, en prenant les puissances extérieures, on obtient donc
$$
\Omega^p_{A[s, s^{-1}]/R} =
A[s, s^{-1}] \otimes_A \Omega^p_{A/R}
\oplus
A[s, s^{-1}] \text{d}\log(s) \otimes_A \Omega^{p - 1}_{A/R}
$$
En prenant les composantes de degré $0$, on trouve
$$
\Omega^p_{A[s, s^{-1}]/R, 0} =
\Omega^p_{A/R} \oplus \text{d}\log(s) \otimes_A \Omega^{p - 1}_{A/R}
$$
ce qui achève la démonstration du lemme.
\end{proof}

\begin{lemma}
\label{lemma-the-complex-for-L-star-gives-chern-class}
Le morphisme de « bord »
$\delta : \Omega^\bullet_{X/S} \to \Omega^\bullet_{X/S}[2]$
dans $D(X, f^{-1}\mathcal{O}_S)$ provenant de
la suite exacte courte du lemme \ref{lemma-the-complex-for-L-star}
est le morphisme de la remarque \ref{remark-cup-product-as-a-map}
pour $\xi = c_1^{dR}(\mathcal{L})$.
\end{lemma}

\begin{proof}
Plus précisément, considérons le décalage
$$
0 \to \Omega^\bullet_{X/S}[1] \to
\Omega^\bullet_{L^\star/S, 0}[1] \to
\Omega^\bullet_{X/S} \to 0
$$
de la suite exacte courte du lemme \ref{lemma-the-complex-for-L-star}.
Comme composante de degré zéro d'une graduation sur
$(L^\star \to X)_*\Omega^\bullet_{L^\star/S}$,
$\Omega^\bullet_{L^\star/S, 0}$ est une
$\mathcal{O}_X$-algèbre différentielle graduée, et le morphisme
$\Omega^\bullet_{X/S} \to \Omega^\bullet_{L^\star/S, 0}$
est un homomorphisme de $\mathcal{O}_X$-algèbres différentielles graduées.
On peut donc considérer $\Omega^\bullet_{X/S}[1] \to
\Omega^\bullet_{L^\star/S, 0}[1]$ comme un morphisme de
$\Omega^\bullet_{X/S}$-modules différentiels gradués à droite sur $X$. Le morphisme
$\text{Res} : \Omega^\bullet_{L^\star/S, 0}[1] \to \Omega^\bullet_{X/S}$
est un morphisme de $\Omega^\bullet_{X/S}$-modules différentiels gradués à droite,
car il est défini localement par la règle
$\text{Res}(\omega' + \text{d}\log(s) \wedge \omega) = \omega$ ; voir
la démonstration du lemme \ref{lemma-the-complex-for-L-star}.
Ainsi, d'après la discussion de
Faisceaux différentiels gradués, section \ref{sdga-section-misc},
$\delta$ provient d'un morphisme
$\delta' : \Omega^\bullet_{X/S} \to \Omega^\bullet_{X/S}[2]$
dans la catégorie dérivée $D(\Omega^\bullet_{X/S}, \text{d})$
des modules différentiels gradués à droite sur le complexe de de Rham.
L'unicité énoncée dans la remarque \ref{remark-cup-product-as-a-map}
montre qu'il suffit de prouver que $\delta(1) = c_1^{dR}(\mathcal{L})$.

\medskip\noindent
Nous affirmons qu'il existe un diagramme commutatif
$$
\xymatrix{
0 \ar[r] &
\mathcal{O}_X^* \ar[r] \ar[d]_{\text{d}\log} &
E \ar[r] \ar[d] &
\underline{\mathbf{Z}} \ar[d] \ar[r] &
0 \\
0 \ar[r] &
\Omega^\bullet_{X/S}[1] \ar[r] &
\Omega^\bullet_{L^\star/S, 0}[1] \ar[r] &
\Omega^\bullet_{X/S} \ar[r] &
0
}
$$
dont la ligne supérieure est une suite exacte courte de faisceaux abéliens dont le
morphisme de bord envoie $1$ sur la classe de $\mathcal{L}$ dans
$H^1(X, \mathcal{O}_X^*)$. Il suffit de démontrer cette assertion,
par compatibilité des morphismes de bord avec les morphismes de suites
exactes courtes. On définit $E$ comme le faisceau associé à la règle
$$
U \longmapsto \{(s, n) \mid
n \in \mathbf{Z},\ s \in \mathcal{L}^{\otimes n}(U)\text{ generator}\}
$$,
la structure de groupe étant donnée par $(s, n) \cdot (t, m) = (s \otimes t, n + m)$.
Le morphisme vertical du milieu envoie $(s, n)$ sur $\text{d}\log(s)$. On obtient ainsi
un morphisme de suites exactes courtes,
car le morphisme $Res : \Omega^1_{L^\star/S, 0} \to \mathcal{O}_X$
construit dans la démonstration du lemme \ref{lemma-the-complex-for-L-star} envoie
$\text{d}\log(s)$ sur $1$ si $s$ est un générateur local de $\mathcal{L}$.
Pour calculer le bord de $1$ dans la ligne supérieure, choisissons des trivialisations locales
$s_i$ de $\mathcal{L}$ sur des ouverts $U_i$ comme à la
section \ref{section-first-chern-class}. Sur les intersections
$U_{i_0i_1} = U_{i_0} \cap U_{i_1}$,
on dispose d'une fonction inversible $f_{i_0i_1}$ telle que
$f_{i_0i_1} = s_{i_1}|_{U_{i_0i_1}} s_{i_0}|_{U_{i_0i_1}}^{-1}$,
et la classe de cohomologie de $\mathcal{L}$ est donnée par le cocycle de {\v C}ech
$\{f_{i_0i_1}\}$. On a alors, bien entendu,
$$
(f_{i_0i_1}, 0) = (s_{i_1}, 1)|_{U_{i_0i_1}} \cdot
(s_{i_0}, 1)|_{U_{i_0i_1}}^{-1}
$$
comme sections de $E$, ce qui achève la démonstration.
\end{proof}

\begin{lemma}
\label{lemma-push-omega-a}
Avec les notations ci-dessus, on a
\begin{enumerate}
\item $\Omega^p_{L^\star/S, n} =
\Omega^p_{L^\star/S, 0} \otimes_{\mathcal{O}_X} \mathcal{L}^{\otimes n}$
pour tout $n \in \mathbf{Z}$, comme $\mathcal{O}_X$-modules quasi-cohérents,
\item $\Omega^\bullet_{X/S} = \Omega^\bullet_{L/X, 0}$
comme complexes, et
\item pour $n > 0$ et $p \geq 0$, on a
$\Omega^p_{L/X, n} = \Omega^p_{L^\star/S, n}$.
\end{enumerate}
\end{lemma}

\begin{proof}
Dans chaque cas, il existe un morphisme canonique défini globalement,
qui est un isomorphisme par des calculs locaux que nous omettons.
\end{proof}

\begin{lemma}
\label{lemma-line-bundle-characteristic-zero}
Dans la situation ci-dessus, supposons qu'il existe un morphisme $S \to \Spec(\mathbf{Q})$.
Alors $\Omega^\bullet_{X/S} \to \pi_*\Omega^\bullet_{L/S}$ est un
quasi-isomorphisme et $H_{dR}^*(X/S) = H_{dR}^*(L/S)$.
\end{lemma}

\begin{proof}
Soit $R$ une $\mathbf{Q}$-algèbre. Soit $A$ une $R$-algèbre.
L'énoncé local affine affirme que le morphisme
$$
\Omega^\bullet_{A/R} \longrightarrow \Omega^\bullet_{A[t]/R}
$$
est un quasi-isomorphisme de complexes de $R$-modules. Il s'agit en fait d'une
équivalence d'homotopie dont un inverse à homotopie près est le morphisme envoyant
$g \omega + g' \text{d}t \wedge \omega'$ sur $g(0)\omega$ pour
$g, g' \in A[t]$ et $\omega, \omega' \in \Omega^\bullet_{A/R}$.
L'homotopie envoie $g \omega + g' \text{d}t \wedge \omega'$
sur $(\int g') \omega'$, où $\int g' \in A[t]$ est le polynôme
de terme constant nul dont la dérivée par rapport à $t$
est $g'$. Bien entendu, on utilise ici le fait que $R$ contient $\mathbf{Q}$,
puisque $\int t^n = (1/n)t^{n + 1}$.
\end{proof}

\begin{example}
\label{example-affine-line}
Le lemme \ref{lemma-line-bundle-characteristic-zero} est
faux en caractéristique positive. Le complexe de de Rham de
$\mathbf{A}^1_k = \Spec(k[x])$ sur un corps $k$ se présente comme une somme directe
$$
k \oplus
\bigoplus\nolimits_{n \geq 1}
(k \cdot t^n \xrightarrow{n}
k \cdot t^{n - 1} \text{d}t)
$$
Ainsi, si la caractéristique de $k$ est $p > 0$,
on voit que $H^0_{dR}(\mathbf{A}^1_k/k)$ et
$H^1_{dR}(\mathbf{A}^1_k/k)$
sont tous deux de dimension infinie sur $k$.
\end{example}








\section{Cohomologie de de Rham de l'espace projectif}
\label{section-projective-space}

\noindent
Soit $A$ un anneau. Soit $n \geq 1$. Le morphisme structural
$\mathbf{P}^n_A \to \Spec(A)$ est propre et lisse de dimension relative
$n$. Il est lisse de dimension relative $n$ et de type fini,
car $\mathbf{P}^n_A$ possède un recouvrement ouvert affine fini par des schémas tous
isomorphes à $\mathbf{A}^n_A$ ; voir Constructions, lemme
\ref{constructions-lemma-standard-covering-projective-space}.
Il est propre, car il est aussi séparé et universellement fermé,
d'après Constructions, lemme \ref{constructions-lemma-projective-space-separated}.
Notons $\mathcal{O}$ et $\mathcal{O}(d)$ respectivement le faisceau structural
$\mathcal{O}_{\mathbf{P}^n_A}$ et les faisceaux tordus de Serre
$\mathcal{O}_{\mathbf{P}^n_A}(d)$.
Notons $\Omega = \Omega_{\mathbf{P}^n_A/A}$ le faisceau
des différentielles relatives et $\Omega^p$ ses puissances extérieures.

\begin{lemma}
\label{lemma-euler-sequence}
Il existe une suite exacte courte
$$
0 \to \Omega \to \mathcal{O}(-1)^{\oplus n + 1} \to \mathcal{O} \to 0
$$
\end{lemma}

\begin{proof}
Pour l'expliciter, rappelons que
$\mathbf{P}^n_A = \text{Proj}(A[T_0, \ldots, T_n])$,
et écrivons symboliquement
$$
\mathcal{O}(-1)^{\oplus n + 1} =
\bigoplus\nolimits_{j = 0, \ldots, n} \mathcal{O}(-1) \text{d}T_j
$$
La première flèche
$$
\Omega \to
\bigoplus\nolimits_{j = 0, \ldots, n} \mathcal{O}(-1) \text{d}T_j
$$
de la suite exacte courte ci-dessus
est donnée, sur chacun des ouverts standards
$D_+(T_i) = \Spec(A[T_0/T_i, \ldots, T_n/T_i])$
mentionnés plus haut, par la règle
$$
\sum\nolimits_{j \not = i} g_j \text{d}(T_j/T_i)
\longmapsto
\sum\nolimits_{j \not = i} g_j/T_i \text{d}T_j
- (\sum\nolimits_{j \not = i} g_jT_j/T_i^2) \text{d}T_i
$$
Cette expression a un sens, puisque $1/T_i$ est une section de $\mathcal{O}(-1)$
sur $D_+(T_i)$. Le morphisme
$$
\bigoplus\nolimits_{j = 0, \ldots, n} \mathcal{O}(-1) \text{d}T_j
\to
\mathcal{O}
$$
est défini en envoyant $\text{d}T_j$ sur $T_j$ ; plus précisément, sur
$D_+(T_i)$, on envoie la section $\sum g_j \text{d}T_j$ sur
$\sum T_jg_j$. Nous omettons de vérifier que l'on obtient ainsi
une suite exacte courte.
\end{proof}

\noindent
Étant donnés un entier $k \in \mathbf{Z}$ et un
$\mathcal{O}_{\mathbf{P}^n_A}$-module quasi-cohérent $\mathcal{F}$,
notons comme d'habitude $\mathcal{F}(k)$ le tordu de Serre d'indice $k$ de $\mathcal{F}$.
Voir Constructions, définition \ref{constructions-definition-twist}.

\begin{lemma}
\label{lemma-twisted-hodge-cohomology-projective-space}
Dans la situation ci-dessus, les groupes de cohomologie vérifient
\begin{enumerate}
\item $H^q(\mathbf{P}^n_A, \Omega^p) = 0$
sauf si $0 \leq p = q \leq n$,
\item pour $0 \leq p \leq n$, le $A$-module
$H^p(\mathbf{P}^n_A, \Omega^p)$ est libre de rang $1$.
\item pour $q > 0$, $k > 0$ et $p$ quelconque, on a
$H^q(\mathbf{P}^n_A, \Omega^p(k)) = 0$, et
\item ajouter d'autres propriétés ici.
\end{enumerate}
\end{lemma}

\begin{proof}
Nous utiliserons les résultats de Cohomologie des schémas, lemme
\ref{coherent-lemma-cohomology-projective-space-over-ring},
sans plus les mentionner. En particulier, les assertions sont vraies
pour $H^q(\mathbf{P}^n_A, \mathcal{O}(k))$.

\medskip\noindent
Démonstration pour $p = 1$. Considérons la suite exacte courte
$$
0 \to \Omega \to \mathcal{O}(-1)^{\oplus n + 1} \to \mathcal{O} \to 0
$$
du lemme \ref{lemma-euler-sequence}. Puisque la cohomologie de $\mathcal{O}(-1)$
est nulle en tout degré, il en résulte que
$H^q(\mathbf{P}^n_A, \Omega)$ est nul sauf en degré $1$,
où il est librement engendré par le bord de $1$ dans
$H^0(\mathbf{P}^n_A, \mathcal{O})$.

\medskip\noindent
Supposons $p > 1$. Considérons la suite exacte courte
ci-dessus comme définissant une filtration à $2$ crans sur $\mathcal{O}(-1)^{\oplus n + 1}$.
La filtration induite sur $\wedge^p\mathcal{O}(-1)^{\oplus n + 1}$ se présente
comme suit :
$$
0 \to \Omega^p \to \wedge^p\left(\mathcal{O}(-1)^{\oplus n + 1}\right)
\to \Omega^{p - 1} \to 0
$$
Observons que $\wedge^p\mathcal{O}(-1)^{\oplus n + 1}$ est isomorphe
à une somme directe comptant autant de facteurs qu'il existe de façons de
choisir, parmi $n + 1$ éléments, $p$ éléments, chaque facteur étant
isomorphe à $\mathcal{O}(-p)$,
et que sa cohomologie est donc nulle en tout degré.
Par l'hypothèse de récurrence, cela montre que $H^q(\mathbf{P}^n_A, \Omega^p)$
est nul sauf si $q = p$, et que $H^p(\mathbf{P}^n_A, \Omega^p)$ est libre
de rang $1$, engendré par le bord du générateur de
$H^{p - 1}(\mathbf{P}^n_A, \Omega^{p - 1})$.

\medskip\noindent
Soit $k > 0$. Observons que $\Omega^n = \mathcal{O}(-n - 1)$, par exemple
en utilisant la suite exacte courte ci-dessus pour $p = n + 1$.
La cohomologie de $\Omega^n(k)$ est donc nulle en degrés strictement positifs.
En utilisant les suites exactes courtes
$$
0 \to \Omega^p(k) \to \wedge^p\left(\mathcal{O}(-1)^{\oplus n + 1}\right)(k)
\to \Omega^{p - 1}(k) \to 0
$$
et une récurrence {\it descendante} sur $p$, on obtient l'annulation de la
cohomologie de $\Omega^p(k)$ en degrés strictement positifs pour tout $p$.
\end{proof}

\begin{lemma}
\label{lemma-hodge-cohomology-projective-space}
On a $H^q(\mathbf{P}^n_A, \Omega^p) = 0$
sauf si $0 \leq p = q \leq n$. Pour $0 \leq p \leq n$, le $A$-module
$H^p(\mathbf{P}^n_A, \Omega^p)$ est libre de rang $1$, de base
$c_1^{Hodge}(\mathcal{O}(1))^p$.
\end{lemma}

\begin{proof}
L'annulation et la liberté résultent du
lemme \ref{lemma-twisted-hodge-cohomology-projective-space}.
Pour $p = 0$, il est bien vrai que
$1 \in H^0(\mathbf{P}^n_A, \mathcal{O})$ est un générateur.

\medskip\noindent
Démonstration pour $p = 1$. Considérons la suite exacte courte
$$
0 \to \Omega \to \mathcal{O}(-1)^{\oplus n + 1} \to \mathcal{O} \to 0
$$
du lemme \ref{lemma-euler-sequence}. Dans la démonstration du
lemme \ref{lemma-twisted-hodge-cohomology-projective-space},
nous avons vu que le générateur de $H^1(\mathbf{P}^n_A, \Omega)$
est le bord $\xi$ de $1 \in H^0(\mathbf{P}^n_A, \mathcal{O})$.
Comme dans la démonstration du lemme \ref{lemma-euler-sequence}, identifions
$\mathcal{O}(-1)^{\oplus n + 1}$ à
$\bigoplus_{j = 0, \ldots, n} \mathcal{O}(-1)\text{d}T_j$.
Considérons le recouvrement ouvert
$$
\mathcal{U} : 
\mathbf{P}^n_A =
\bigcup\nolimits_{i = 0, \ldots, n} D_{+}(T_i)
$$
La restriction de la section globale $1$ de $\mathcal{O}$
à $U_i = D_+(T_i)$ se relève en la section $T_i^{-1} \text{d}T_i$ de
$\bigoplus \mathcal{O}(-1)\text{d}T_j$ sur $U_i$. Ainsi, le cocycle
représentant $\xi$ est donné par
$$
T_{i_1}^{-1} \text{d}T_{i_1} - T_{i_0}^{-1} \text{d}T_{i_0} =
\text{d}\log(T_{i_1}/T_{i_0}) \in \Omega(U_{i_0i_1})
$$
D'autre part, pour tout $i$, la section $T_i$ trivialise
$\mathcal{O}(1)$ sur $U_i$. On voit donc que
$f_{i_0i_1} = T_{i_1}/T_{i_0} \in \mathcal{O}^*(U_{i_0i_1})$
est le cocycle représentant $\mathcal{O}(1)$ dans $\Pic(\mathbf{P}^n_A)$ ;
voir la section \ref{section-first-chern-class}.
Par conséquent, $c_1^{Hodge}(\mathcal{O}(1))$
est donné par le cocycle $\text{d}\log(T_{i_1}/T_{i_0})$,
qui coïncide avec celui obtenu pour $\xi$ ci-dessus.

\medskip\noindent
Démonstration par récurrence pour $p$ quelconque. Les cas initiaux $p = 0, 1$ ont été traités
ci-dessus. Supposons $p > 1$. Dans la démonstration du
lemme \ref{lemma-twisted-hodge-cohomology-projective-space},
nous avons vu que le générateur de $H^p(\mathbf{P}^n_A, \Omega^p)$
est le bord de $c_1^{Hodge}(\mathcal{O}(1))^{p - 1}$
dans la suite exacte longue de cohomologie associée à
$$
0 \to \Omega^p \to \wedge^p\left(\mathcal{O}(-1)^{\oplus n + 1}\right)
\to \Omega^{p - 1} \to 0
$$
D'après le calcul de la section \ref{section-first-chern-class},
la classe de cohomologie $c_1^{Hodge}(\mathcal{O}(1))^{p - 1}$
est représentée, au signe près, par le cocycle de composantes
$$
\beta_{i_0 \ldots i_{p - 1}} =
\text{d}\log(T_{i_1}/T_{i_0}) \wedge
\text{d}\log(T_{i_2}/T_{i_1}) \wedge \ldots \wedge
\text{d}\log(T_{i_{p - 1}}/T_{i_{p - 2}})
$$
dans $\Omega^{p - 1}(U_{i_0 \ldots i_{p - 1}})$. Ces
$\beta_{i_0 \ldots i_{p - 1}}$ se relèvent en les sections
$\tilde \beta_{i_0 \ldots i_{p -1}} =
T_{i_0}^{-1}\text{d}T_{i_0} \wedge \beta_{i_0 \ldots i_{p - 1}}$
de $\wedge^p(\bigoplus \mathcal{O}(-1) \text{d}T_j)$ sur
$U_{i_0 \ldots i_{p - 1}}$. On en déduit que le générateur de
$H^p(\mathbf{P}^n_A, \Omega^p)$ est donné par le cocycle dont
les composantes sont
\begin{align*}
\sum\nolimits_{a = 0}^p (-1)^a
\tilde \beta_{i_0 \ldots \hat{i_a} \ldots i_p}
& =
T_{i_1}^{-1}\text{d}T_{i_1} \wedge \beta_{i_1 \ldots i_p}
+ \sum\nolimits_{a = 1}^p (-1)^a
T_{i_0}^{-1}\text{d}T_{i_0} \wedge
\beta_{i_0 \ldots \hat{i_a} \ldots i_p} \\
& =
(T_{i_1}^{-1}\text{d}T_{i_1} - T_{i_0}^{-1}\text{d}T_{i_0}) \wedge
\beta_{i_1 \ldots i_p} +
T_{i_0}^{-1}\text{d}T_{i_0} \wedge \text{d}(\beta)_{i_0 \ldots i_p} \\
& =
\text{d}\log(T_{i_1}/T_{i_0}) \wedge \beta_{i_1 \ldots i_p}
\end{align*}
considérées comme sections de $\Omega^p$ sur $U_{i_0 \ldots i_p}$.
C'est, au signe près, le cocycle représentant
$c_1^{Hodge}(\mathcal{O}(1))^p$, ce qui achève la démonstration.
\end{proof}

\begin{lemma}
\label{lemma-de-rham-cohomology-projective-space}
Pour $0 \leq i \leq n$, le groupe de cohomologie de de Rham
$H^{2i}_{dR}(\mathbf{P}^n_A/A)$ est un $A$-module libre de rang $1$,
de base $c_1^{dR}(\mathcal{O}(1))^i$.
En tout autre degré, la cohomologie de de Rham de $\mathbf{P}^n_A$
sur $A$ est nulle.
\end{lemma}

\begin{proof}
Considérons la suite spectrale de Hodge vers de Rham de la
section \ref{section-hodge-to-de-rham}.
Le calcul de la cohomologie de Hodge de $\mathbf{P}^n_A$ sur $A$
effectué dans le lemme \ref{lemma-hodge-cohomology-projective-space}
montre que cette suite spectrale dégénère à la page $E_1$.
On voit ainsi que $H^{2i}_{dR}(\mathbf{P}^n_A/A)$ est un
$A$-module libre de rang $1$ pour $0 \leq i \leq n$, et qu'il est nul sinon.
Observons que $c_1^{dR}(\mathcal{O}(1))^i \in H^{2i}_{dR}(\mathbf{P}^n_A/A)$
pour $i = 0, \ldots, n$ et que, pour $i = n$, cet élément est
l'image de $c_1^{Hodge}(\mathcal{L})^n$ par le morphisme de complexes
$$
\Omega^n_{\mathbf{P}^n_A/A}[-n]
\longrightarrow
\Omega^\bullet_{\mathbf{P}^n_A/A}
$$
Cela résulte, par exemple, de la discussion de la remarque \ref{remark-truncations}
ou de la description explicite des cocycles représentant ces classes à la
section \ref{section-first-chern-class}.
La suite spectrale montre que le morphisme induit
$$
H^n(\mathbf{P}^n_A, \Omega^n_{\mathbf{P}^n_A/A}) \longrightarrow
H^{2n}_{dR}(\mathbf{P}^n_A/A)
$$
est un isomorphisme et, puisque $c_1^{Hodge}(\mathcal{L})^n$ engendre
la source (lemme \ref{lemma-hodge-cohomology-projective-space}),
on en déduit que $c_1^{dR}(\mathcal{L})^n$ engendre
le but. La $A$-bilinéarité des cup-produits
entraîne que $c_1^{dR}(\mathcal{L})^i$
engendre également $H^{2i}_{dR}(\mathbf{P}^n_A/A)$ pour
$0 \leq i \leq n$.
\end{proof}









\section{La suite spectrale d'un morphisme lisse}
\label{section-relative-spectral-sequence}

\noindent
Considérons un diagramme commutatif de schémas
$$
\xymatrix{
X \ar[rr]_f \ar[rd]_p & & Y \ar[ld]^q \\
& S
}
$$
où $f$ est un morphisme lisse. On obtient alors une suite exacte courte
localement scindée
$$
0 \to f^*\Omega_{Y/S} \to \Omega_{X/S} \to \Omega_{X/Y} \to 0
$$
d'après Morphismes, lemme \ref{morphisms-lemma-triangle-differentials-smooth}.
Considérons-la comme une filtration décroissante $F$ sur $\Omega_{X/S}$,
avec $F^0\Omega_{X/S} = \Omega_{X/S}$, $F^1\Omega_{X/S} = f^*\Omega_{Y/S}$ et
$F^2\Omega_{X/S} = 0$. En appliquant le foncteur $\wedge^p$, on obtient,
pour tout $p$, une filtration induite
$$
\Omega^p_{X/S} = F^0\Omega^p_{X/S} \supset
F^1\Omega^p_{X/S} \supset
F^2\Omega^p_{X/S} \supset \ldots \supset F^{p + 1}\Omega^p_{X/S} = 0
$$
dont les quotients successifs sont
$$
\text{gr}^k\Omega^p_{X/S} =
F^k\Omega^p_{X/S}/F^{k + 1}\Omega^p_{X/S} =
f^*\Omega^k_{Y/S} \otimes_{\mathcal{O}_X} \Omega^{p - k}_{X/Y} =
f^{-1}\Omega^k_{Y/S} \otimes_{f^{-1}\mathcal{O}_Y} \Omega^{p - k}_{X/Y}
$$
pour $k = 0, \ldots, p$. En fait, le lecteur peut vérifier, à l'aide de la
règle de Leibniz, que $F^k\Omega^\bullet_{X/S}$ est un sous-complexe de
$\Omega^\bullet_{X/S}$. On munit ainsi $\Omega^\bullet_{X/S}$
d'une structure de complexe filtré. On peut aussi le voir en observant
que
$$
F^k\Omega^\bullet_{X/S} =
\Im\left(\wedge :
\text{Tot}(
f^{-1}\sigma_{\geq k}\Omega^\bullet_{Y/S} \otimes_{p^{-1}\mathcal{O}_S}
\Omega^\bullet_{X/S})
\longrightarrow
\Omega^\bullet_{X/S}\right)
$$
est l'image d'un morphisme de complexes sur $X$. Le complexe filtré
$$
\Omega^\bullet_{X/S} = F^0\Omega^\bullet_{X/S} \supset
F^1\Omega^\bullet_{X/S} \supset F^2\Omega^\bullet_{X/S} \supset \ldots
$$
a pour composantes graduées associées
$$
\text{gr}^k\Omega^\bullet_{X/S} = 
f^{-1}\Omega^k_{Y/S}[-k] \otimes_{f^{-1}\mathcal{O}_Y} \Omega^\bullet_{X/Y}
$$
d'après ce qui précède.

\begin{lemma}
\label{lemma-spectral-sequence-smooth}
Soit $f : X \to Y$ un morphisme quasi-compact, quasi-séparé et lisse
de schémas sur un schéma de base $S$. Il existe une suite spectrale bornée
de première page
$$
E_1^{p, q} =
H^q(\Omega^p_{Y/S} \otimes_{\mathcal{O}_Y}^\mathbf{L} Rf_*\Omega^\bullet_{X/Y})
$$
convergeant vers $R^{p + q}f_*\Omega^\bullet_{X/S}$.
\end{lemma}

\begin{proof}
Considérons $\Omega^\bullet_{X/S}$ comme un complexe filtré muni de la
filtration introduite ci-dessus. La suite spectrale est celle de
Cohomologie, lemme
\ref{cohomology-lemma-relative-spectral-sequence-filtered-object}.
D'après Catégories dérivées des schémas, lemme
\ref{perfect-lemma-cohomology-de-rham-base-change}, on a
$$
Rf_*\text{gr}^k\Omega^\bullet_{X/S} =
\Omega^k_{Y/S}[-k] \otimes_{\mathcal{O}_Y}^\mathbf{L} Rf_*\Omega^\bullet_{X/Y}
$$,
d'où la conclusion.
\end{proof}

\begin{remark}
\label{remark-gauss-manin}
Dans le lemme \ref{lemma-spectral-sequence-smooth}, considérons les faisceaux de cohomologie
$$
\mathcal{H}^q_{dR}(X/Y) = H^q(Rf_*\Omega^\bullet_{X/Y})
$$
Si $f$ est propre en plus d'être lisse et si $S$ est un schéma sur
$\mathbf{Q}$, alors $\mathcal{H}^q_{dR}(X/Y)$ est localement libre de rang fini (insérer
ici une référence ultérieure). Si l'on suppose seulement que les $\mathcal{H}^q_{dR}(X/Y)$
sont des $\mathcal{O}_Y$-modules plats, on obtient (nous omettons le bref argument)
$$
E_1^{p, q} =
\Omega^p_{Y/S} \otimes_{\mathcal{O}_Y} \mathcal{H}^q_{dR}(X/Y)
$$,
et les différentielles de la suite spectrale sont des morphismes
$$
d_1^{p, q} :
\Omega^p_{Y/S} \otimes_{\mathcal{O}_Y} \mathcal{H}^q_{dR}(X/Y)
\longrightarrow
\Omega^{p + 1}_{Y/S} \otimes_{\mathcal{O}_Y} \mathcal{H}^q_{dR}(X/Y)
$$
En particulier, pour $p = 0$, on obtient un morphisme
$d_1^{0, q} : \mathcal{H}^q_{dR}(X/Y) \to
\Omega^1_{Y/S} \otimes_{\mathcal{O}_Y} \mathcal{H}^q_{dR}(X/Y)$
qui est une connexion intégrable
$\nabla$ (insérer ici une référence ultérieure),
et le complexe
$$
\mathcal{H}^q_{dR}(X/Y) \to
\Omega^1_{Y/S} \otimes_{\mathcal{O}_Y} \mathcal{H}^q_{dR}(X/Y) \to
\Omega^2_{Y/S} \otimes_{\mathcal{O}_Y} \mathcal{H}^q_{dR}(X/Y) \to \ldots
$$
dont les différentielles sont données par $d_1^{\bullet, q}$
est le complexe de de Rham de $\nabla$.
La connexion $\nabla$ est appelée {\it connexion de Gauss--Manin}.
\end{remark}






\section{Théorèmes de type Leray--Hirsch}
\label{section-leray-hirsch}

\noindent
Dans cette section, nous montrons que, pour un morphisme lisse et propre,
on peut parfois exprimer la cohomologie de de Rham de la source en fonction
de celle du but.

\begin{lemma}
\label{lemma-relative-global-generation-on-fibres}
Soit $f : X \to Y$ un morphisme lisse et propre de schémas.
Soient $N$ et $n_1, \ldots, n_N \geq 0$ des entiers, et soient
$\xi_i \in H^{n_i}_{dR}(X/Y)$, $1 \leq i \leq N$.
Supposons que, pour tout point $y \in Y$, les images de $\xi_1, \ldots, \xi_N$
dans $H^*_{dR}(X_y/y)$ forment une base sur $\kappa(y)$. Alors le morphisme
$$
\bigoplus\nolimits_{i = 1}^N \mathcal{O}_Y[-n_i]
\longrightarrow
Rf_*\Omega^\bullet_{X/Y}
$$
associé à $\xi_1, \ldots, \xi_N$ est un isomorphisme.
\end{lemma}

\begin{proof}
D'après le lemme \ref{lemma-proper-smooth-de-Rham},
$Rf_*\Omega^\bullet_{X/Y}$ est un objet parfait de $D(\mathcal{O}_Y)$
dont la formation commute à tout changement de base.
Le morphisme du lemme est donc un morphisme $a : K \to L$
entre objets parfaits de $D(\mathcal{O}_Y)$
dont la restriction dérivée à tout point est un isomorphisme,
d'après notre hypothèse sur les fibres. Le cône $C$ de $a$ est alors un objet parfait
de $D(\mathcal{O}_Y)$ (Cohomologie, lemme
\ref{cohomology-lemma-two-out-of-three-perfect}) dont
la restriction dérivée à tout point est nulle. Il en résulte que $C$
est nul, d'après Compléments d'algèbre, lemme
\ref{more-algebra-lemma-lift-perfect-from-residue-field},
et que $a$ est un isomorphisme. (On utilise aussi Catégories dérivées des schémas,
lemmes \ref{perfect-lemma-affine-compare-bounded} et
\ref{perfect-lemma-perfect-affine}, pour la traduction en termes algébriques.)
\end{proof}

\noindent
Nous démontrons d'abord le résultat principal de cette section dans le
cas particulier suivant.

\begin{lemma}
\label{lemma-global-generation-on-fibres}
Soit $f : X \to Y$ un morphisme lisse et propre de schémas sur une base $S$.
Supposons que
\begin{enumerate}
\item $Y$ et $S$ sont affines, et
\item il existe des entiers $N$ et $n_1, \ldots, n_N \geq 0$ et des classes
$\xi_i \in H^{n_i}_{dR}(X/S)$, $1 \leq i \leq N$, tels que,
pour tout point $y \in Y$, les images de $\xi_1, \ldots, \xi_N$
dans $H^*_{dR}(X_y/y)$ forment une base sur $\kappa(y)$.
\end{enumerate}
Alors l'application
$$
\bigoplus\nolimits_{i = 1}^N H^*_{dR}(Y/S) \longrightarrow
H^*_{dR}(X/S), \quad
(a_1, \ldots, a_N) \longmapsto  \sum \xi_i \cup f^*a_i
$$
est un isomorphisme.
\end{lemma}

\begin{proof}
Écrivons $Y = \Spec(A)$ et $S = \Spec(R)$.
Dans ce cas, $\Omega^\bullet_{A/R}$ calcule
$R\Gamma(Y, \Omega^\bullet_{Y/S})$ d'après le lemme \ref{lemma-de-rham-affine}.
Choisissons un recouvrement ouvert affine fini $\mathcal{U} : X = \bigcup_{i \in I} U_i$.
Considérons le complexe
$$
K^\bullet =
\text{Tot}(\check{\mathcal{C}}^\bullet(\mathcal{U}, \Omega_{X/S}^\bullet))
$$
comme dans
Cohomologie, section \ref{cohomology-section-cech-cohomology-of-complexes}.
Rassemblons quelques propriétés de ce complexe, dont la plupart
figurent dans la référence que nous venons de donner :
\begin{enumerate}
\item $K^\bullet$ est un complexe de $R$-modules dont les termes sont des
$A$-modules,
\item $K^\bullet$ représente $R\Gamma(X, \Omega^\bullet_{X/S})$ dans $D(R)$
(Cohomologie des schémas, lemme
\ref{coherent-lemma-quasi-coherent-affine-cohomology-zero} et
Cohomologie, lemme \ref{cohomology-lemma-cech-complex-complex-computes}),
\item il existe un morphisme naturel $\Omega^\bullet_{A/R} \to K^\bullet$
de complexes de $R$-modules, $A$-linéaire sur les termes, qui
induit le morphisme d'image inverse $H^*_{dR}(Y/S) \to H^*_{dR}(X/S)$
en cohomologie,
\item $K^\bullet$ possède une multiplication notée $\wedge$
qui en fait une $R$-algèbre différentielle graduée,
\item la multiplication sur $K^\bullet$
induit le cup-produit sur $H^*_{dR}(X/S)$
(Cohomologie, section \ref{cohomology-section-cup-product}),
\item la filtration $F$ sur $\Omega^*_{X/S}$ induit une filtration
$$
K^\bullet =
F^0K^\bullet \supset F^1K^\bullet \supset F^2K^\bullet \supset \ldots
$$
par des sous-complexes de $K^\bullet$ telle que
\begin{enumerate}
\item $F^kK^n \subset K^n$ est un $A$-sous-module,
\item $F^kK^\bullet \wedge F^lK^\bullet \subset F^{k + l}K^\bullet$,
\item $\text{gr}^kK^\bullet$ est un complexe de $A$-modules,
\item $\text{gr}^0K^\bullet =
\text{Tot}(\check{\mathcal{C}}^\bullet(\mathcal{U}, \Omega_{X/Y}^\bullet))$
et représente $R\Gamma(X, \Omega^\bullet_{X/Y})$ dans $D(A)$,
\item la multiplication induit un isomorphisme
$\Omega^k_{A/R}[-k] \otimes_A \text{gr}^0K^\bullet \to \text{gr}^kK^\bullet$
\end{enumerate}
\end{enumerate}
Nous omettons les démonstrations détaillées de ces assertions ; voir la discussion
qui précède la construction de la suite spectrale du
lemme \ref{lemma-spectral-sequence-smooth}.

\medskip\noindent
Pour tout $i = 1, \ldots, N$, choisissons un cocycle $x_i \in K^{n_i}$
représentant $\xi_i$. Considérons ensuite le morphisme de complexes
$$
\tilde x :
M^\bullet = \bigoplus\nolimits_{i = 1, \ldots, N}
\Omega^\bullet_{A/R}[-n_i]
\longrightarrow
K^\bullet
$$
qui envoie $\omega$ dans le facteur d'indice $i$ sur $x_i \wedge \omega$.
Il reste seulement à montrer que ce morphisme est un quasi-isomorphisme.
Munissons $M^\bullet$ d'une structure de complexe filtré
par la règle
$$
F^kM^\bullet =
\bigoplus\nolimits_{i = 1, \ldots, N}
(\sigma_{\geq k}\Omega^\bullet_{A/R})[-n_i]
$$
Avec ce choix, $\tilde x$ est un morphisme de complexes filtrés.
Observons que $\text{gr}^0M^\bullet = \bigoplus A[-n_i]$
et que la multiplication induit un isomorphisme
$\Omega^k_{A/R}[-k] \otimes_A \text{gr}^0M^\bullet \to \text{gr}^kM^\bullet$.
Par construction et d'après le lemme \ref{lemma-relative-global-generation-on-fibres},
$$
\text{gr}^0\tilde x :
\text{gr}^0M^\bullet \longrightarrow
\text{gr}^0K^\bullet
$$
est un isomorphisme dans $D(A)$. Il en résulte que, pour tout $k \geq 0$,
on obtient des isomorphismes
$$
\text{gr}^k \tilde x :
\text{gr}^kM^\bullet = \Omega^k_{A/R}[-k] \otimes_A \text{gr}^0M^\bullet
\longrightarrow
\Omega^k_{A/R}[-k] \otimes_A \text{gr}^0K^\bullet =
\text{gr}^kK^\bullet
$$
dans $D(A)$. En effet, le complexe
$\text{gr}^0K^\bullet =
\text{Tot}(\check{\mathcal{C}}^\bullet(\mathcal{U}, \Omega_{X/Y}^\bullet))$
est K-plat comme complexe de $A$-modules, d'après Catégories dérivées des schémas,
lemme \ref{perfect-lemma-K-flat}.
Le produit tensoriel du membre de droite est donc le
produit tensoriel dérivé, ce qui est aussi vrai par inspection pour le membre de gauche.
Enfin, le produit tensoriel dérivé
$\Omega^k_{A/R}[-k] \otimes_A^\mathbf{L} -$ est un foncteur sur $D(A)$
et envoie donc les isomorphismes sur des isomorphismes.
Par récurrence sur $k$, on en déduit que
$$
\tilde x : M^\bullet/F^kM^\bullet \to K^\bullet/F^kK^\bullet
$$
est un isomorphisme dans $D(R)$, puisqu'on dispose des suites exactes courtes
$$
0 \to F^kM^\bullet/F^{k + 1}M^\bullet \to
M^\bullet/F^{k + 1}M^\bullet \to
\text{gr}^kM^\bullet \to 0
$$
et de même pour $K^\bullet$. Cela prouve que $\tilde x$ est un
quasi-isomorphisme, les filtrations étant finies en chaque degré.
\end{proof}

\begin{proposition}
\label{proposition-global-generation-on-fibres}
Soit $f : X \to Y$ un morphisme lisse et propre de schémas sur une base $S$.
Soient $N$ et $n_1, \ldots, n_N \geq 0$ des entiers, et soient
$\xi_i \in H^{n_i}_{dR}(X/S)$, $1 \leq i \leq N$.
Supposons que, pour tout point $y \in Y$, les images de $\xi_1, \ldots, \xi_N$
dans $H^*_{dR}(X_y/y)$ forment une base sur $\kappa(y)$. Le morphisme
$$
\tilde \xi = \bigoplus \tilde \xi_i[-n_i] :
\bigoplus \Omega^\bullet_{Y/S}[-n_i]
\longrightarrow
Rf_*\Omega^\bullet_{X/S}
$$
(voir la démonstration) est un isomorphisme dans $D(Y, (Y \to S)^{-1}\mathcal{O}_S)$, et
l'application correspondante
$$
\bigoplus\nolimits_{i = 1}^N H^*_{dR}(Y/S) \longrightarrow
H^*_{dR}(X/S), \quad
(a_1, \ldots, a_N) \longmapsto  \sum \xi_i \cup f^*a_i
$$
est un isomorphisme.
\end{proposition}

\begin{proof}
Notons $p : X \to S$ et $q : Y \to S$ les morphismes structuraux.
Soit $\xi'_i : \Omega^\bullet_{X/S} \to \Omega^\bullet_{X/S}[n_i]$
le morphisme de la remarque \ref{remark-cup-product-as-a-map} correspondant
à $\xi_i$. Notons
$$
\tilde \xi_i :
\Omega^\bullet_{Y/S} \to Rf_*\Omega^\bullet_{X/S}[n_i]
$$
la composée de $\xi'_i$ avec le morphisme canonique
$\Omega^\bullet_{Y/S} \to Rf_*\Omega^\bullet_{X/S}$.
En utilisant
$$
R\Gamma(Y, Rf_*\Omega^\bullet_{X/S}) = R\Gamma(X, \Omega^\bullet_{X/S})
$$,
le morphisme $\tilde \xi_i$ induit en cohomologie l'application $\eta \mapsto \xi_i \cup f^*\eta$
de $H^m_{dR}(Y/S)$ dans $H^{m + n}_{dR}(X/S)$.
De plus, puisque la formation de $\xi'_i$ commute aux
restrictions aux ouverts, il en va de même de celle de $\tilde \xi_i$.

\medskip\noindent
On peut donc considérer le morphisme
$$
\tilde \xi = \bigoplus \tilde \xi_i[-n_i] :
\bigoplus \Omega^\bullet_{Y/S}[-n_i]
\longrightarrow
Rf_*\Omega^\bullet_{X/S}
$$
Pour démontrer le lemme, il suffit de montrer que c'est un isomorphisme dans
$D(Y, q^{-1}\mathcal{O}_S)$. Si l'on pouvait montrer que $\tilde \xi$
provient d'un morphisme de complexes filtrés (pour des filtrations convenables),
on pourrait appliquer la suite spectrale du
lemme \ref{lemma-spectral-sequence-smooth} pour conclure.
Cela demanderait plus de travail qu'il n'est nécessaire ; notre méthode
consistera plutôt à nous ramener au cas affine (dont la démonstration utilise,
en un certain sens, la suite spectrale).

\medskip\noindent
En effet, si $Y' \subset Y$ est un ouvert quelconque d'image inverse
$X' \subset X$, alors $\tilde \xi|_{X'}$ induit l'application
$$
\bigoplus\nolimits_{i = 1}^N H^*_{dR}(Y'/S) \longrightarrow
H^*_{dR}(X'/S), \quad
(a_1, \ldots, a_N) \longmapsto  \sum \xi_i|_{X'} \cup f^*a_i
$$
en cohomologie sur $Y'$ ; voir la discussion ci-dessus.
Il suffit donc de trouver une base de la topologie
de $Y$ telle que la proposition soit vraie pour les ouverts de cette base
(en particulier, on peut alors oublier le morphisme $\tilde \xi$).
On se ramène ainsi au cas où $Y$ et $S$
sont affines, traité par le lemme \ref{lemma-global-generation-on-fibres},
ce qui achève la démonstration.
\end{proof}





\section{Formule du fibré projectif}
\label{section-projective-space-bundle-formula}

\noindent
Le titre dit tout.

\begin{proposition}
\label{proposition-projective-space-bundle-formula}
Soit $X \to S$ un morphisme de schémas. Soit $\mathcal{E}$ un
$\mathcal{O}_X$-module localement libre de rang constant $r$. Considérons le morphisme
$p : P = \mathbf{P}(\mathcal{E}) \to X$.
Alors l'application
$$
\bigoplus\nolimits_{i = 0, \ldots, r - 1} H^*_{dR}(X/S)
\longrightarrow
H^*_{dR}(P/S)
$$
définie par la règle
$$
(a_0, \ldots, a_{r - 1}) \longmapsto
\sum\nolimits_{i = 0, \ldots, r - 1} c_1^{dR}(\mathcal{O}_P(1))^i \cup p^*(a_i)
$$
est un isomorphisme.
\end{proposition}

\begin{proof}
Choisissons un ouvert affine $\Spec(A) \subset X$ sur lequel $\mathcal{E}$ se restreint
au module localement libre trivial $\mathcal{O}_{\Spec(A)}^{\oplus r}$.
Alors $P \times_X \Spec(A) = \mathbf{P}^{r - 1}_A$. On voit ainsi que
$p$ est propre et lisse ; voir la section \ref{section-projective-space}.
De plus, les classes $c_1^{dR}(\mathcal{O}_P(1))^i$, $i = 0, 1, \ldots, r - 1$,
restreintes à une fibre $X_y = \mathbf{P}^{r - 1}_y$, engendrent librement la
cohomologie de de Rham $H^*_{dR}(X_y/y)$ sur $\kappa(y)$ ; voir
le lemme \ref{lemma-de-rham-cohomology-projective-space}. Nous avons donc vérifié les
hypothèses de la proposition \ref{proposition-global-generation-on-fibres},
d'où le résultat.
\end{proof}

\begin{remark}
\label{remark-projective-space-bundle-formula}
Dans la situation de la
proposition \ref{proposition-projective-space-bundle-formula},
on obtient en outre que le morphisme
$$
\tilde \xi :
\bigoplus\nolimits_{t = 0, \ldots, r - 1}
\Omega^\bullet_{X/S}[-2t]
\longrightarrow
Rp_*\Omega^\bullet_{P/S}
$$
est un isomorphisme dans $D(X, (X \to S)^{-1}\mathcal{O}_X)$, ce qui résulte
immédiatement de l'application de la
proposition \ref{proposition-global-generation-on-fibres}.
Notons que la flèche pour $t = 0$ est simplement le morphisme canonique
$c_{P/X} : \Omega^\bullet_{X/S} \to Rp_*\Omega^\bullet_{P/S}$
de la section \ref{section-de-rham-complex}.
On peut en fait préciser davantage ce morphisme dans ce cas particulier.
Considérons le morphisme canonique
$$
\xi' : \Omega^\bullet_{P/S} \to \Omega^\bullet_{P/S}[2]
$$
de la remarque \ref{remark-cup-product-as-a-map} correspondant à
$c_1^{dR}(\mathcal{O}_P(1))$. Alors
$$
\xi'[2(t - 1)] \circ \ldots \circ \xi'[2] \circ \xi' : 
\Omega^\bullet_{P/S} \to \Omega^\bullet_{P/S}[2t]
$$
est le morphisme de la remarque \ref{remark-cup-product-as-a-map} correspondant à
$c_1^{dR}(\mathcal{O}_P(1))^t$. En suivant les choix effectués dans la
démonstration de la proposition \ref{proposition-global-generation-on-fibres},
on obtient la valeur
$$
\tilde \xi|_{\Omega^\bullet_{X/S}[-2t]} =
Rp_*\xi'[-2] \circ \ldots \circ Rp_*\xi'[-2(t - 1)] \circ
Rp_*\xi'[-2t] \circ c_{P/X}[-2t]
$$
pour la restriction de notre isomorphisme au facteur direct
$\Omega^\bullet_{X/S}[-2t]$. Cela a la conséquence simple
suivante, que nous utiliserons plus bas : posons
$$
M = \bigoplus\nolimits_{t = 1, \ldots, r - 1} \Omega^\bullet_{X/S}[-2t]
\quad\text{et}\quad
K = \bigoplus\nolimits_{t = 0, \ldots, r - 2} \Omega^\bullet_{X/S}[-2t]
$$,
considérés comme des sous-complexes de la source de la flèche $\tilde \xi$.
Il résulte formellement de la discussion ci-dessus que
$$
c_{P/X} \oplus
\tilde \xi|_M :
\Omega^\bullet_{X/S} \oplus M \longrightarrow
Rp_*\Omega^\bullet_{P/S}
$$
est un isomorphisme et que le diagramme
$$
\xymatrix{
K \ar[d]_{\tilde \xi|_K} \ar[r]_{\text{id}} &
M[2] \ar[d]^{(\tilde \xi|_M)[2]} \\
Rp_*\Omega^\bullet_{P/S} \ar[r]^{Rp_*\xi'} &
Rp_*\Omega^\bullet_{P/S}[2]
}
$$
est commutatif, où $\text{id} : K \to M[2]$ identifie le facteur
correspondant à $t$ dans la décomposition de $K$ au facteur
correspondant à $t + 1$ dans la décomposition de $M$.
\end{remark}









\section{Pôles logarithmiques le long d'un diviseur}
\label{section-divisor}

\noindent
Soit $X \to S$ un morphisme de schémas. Soit $Y \subset X$ un
diviseur de Cartier effectif. Si $X$ est, étale-localement le long de $Y$,
de la forme $Y \times \mathbf{A}^1$, il existe une suite exacte courte canonique
de complexes
$$
0 \to \Omega^\bullet_{X/S} \to
\Omega^\bullet_{X/S}(\log Y) \to
\Omega^\bullet_{Y/S}[-1] \to 0
$$
possédant de nombreuses propriétés utiles, que nous étudierons dans cette section. Il existe une
variante de cette construction dans laquelle on part d'un diviseur à croisements
normaux
(Morphismes étales, définition \ref{etale-definition-strict-normal-crossings}),
que nous étudierons ailleurs (insérer ici une référence ultérieure).

\begin{definition}
\label{definition-local-product}
Soit $X \to S$ un morphisme de schémas. Soit $Y \subset X$ un
diviseur de Cartier effectif. On dit que le
{\it complexe de de Rham à pôles logarithmiques est défini pour $Y \subset X$ sur $S$}
si, pour tout $y \in Y$ et toute équation locale $f \in \mathcal{O}_{X, y}$
de $Y$, on a
\begin{enumerate}
\item $\mathcal{O}_{X, y} \to \Omega_{X/S, y}$, $g \mapsto g \text{d}f$,
est une injection scindée, et
\item $\Omega^p_{X/S, y}$ est sans $f$-torsion pour tout $p$.
\end{enumerate}
\end{definition}

\noindent
Un calcul local facile montre qu'il suffit, pour tout $y \in Y$,
de trouver une équation locale $f$ pour laquelle les conditions (1) et (2) sont satisfaites.

\begin{lemma}
\label{lemma-log-complex}
Soit $X \to S$ un morphisme de schémas. Soit $Y \subset X$ un
diviseur de Cartier effectif.
Supposons que le complexe de de Rham à pôles logarithmiques soit défini pour $Y \subset X$ sur $S$.
Il existe une suite exacte courte canonique
de complexes
$$
0 \to \Omega^\bullet_{X/S} \to
\Omega^\bullet_{X/S}(\log Y) \to
\Omega^\bullet_{Y/S}[-1] \to 0
$$
\end{lemma}

\begin{proof}
Notre hypothèse est que, pour tout $y \in Y$ et toute équation locale
$f \in \mathcal{O}_{X, y}$ de $Y$, on a
$$
\Omega_{X/S, y} = \mathcal{O}_{X, y}\text{d}f \oplus M
\quad\text{et}\quad
\Omega^p_{X/S, y} = \wedge^{p - 1}(M)\text{d}f \oplus \wedge^p(M)
$$
pour un certain module $M$ dont les puissances extérieures sont sans $f$-torsion, à savoir $\wedge^p(M)$.
Il en résulte que
$$
\Omega^p_{Y/S, y} = \wedge^p(M/fM) = \wedge^p(M)/f\wedge^p(M)
$$
Nous utiliserons tacitement ces faits dans la suite.
En particulier, les faisceaux $\Omega^p_{X/S}$ n'ont pas de sections locales non nulles
à support dans $Y$, et l'on dispose d'une inclusion canonique
$$
\Omega^p_{X/S} \subset \Omega^p_{X/S}(Y)
$$ ;
voir Compléments de platitude, section \ref{flat-section-eta}. Soit $U = \Spec(A)$
un sous-schéma ouvert affine tel que $Y \cap U = V(f)$ pour un certain
élément non diviseur de zéro $f \in A$. Considérons le $\mathcal{O}_U$-sous-module
de $\Omega^p_{X/S}(Y)|_U$ engendré par
$\Omega^p_{X/S}|_U$ et $\text{d}\log(f) \wedge \Omega^{p - 1}_{X/S}$,
où $\text{d}\log(f) = f^{-1}\text{d}(f)$.
Il est indépendant du choix de $f$, car un autre générateur de
l'idéal de $Y$ sur $U$ est égal à $uf$ pour une unité $u \in A$, et l'on obtient
$$
\text{d}\log(uf) - \text{d}\log(f) = \text{d}\log(u) = u^{-1}\text{d}u
$$,
qui est une section de $\Omega_{X/S}$ sur $U$. Ces faisceaux locaux
se recollent en un sous-module quasi-cohérent
$$
\Omega^p_{X/S} \subset \Omega^p_{X/S}(\log Y) \subset \Omega^p_{X/S}(Y)
$$
Convenons de considérer $\Omega^p_{Y/S}$ comme un
$\mathcal{O}_X$-module quasi-cohérent. Il existe un unique morphisme surjectif
$\mathcal{O}_X$-linéaire
$$
\text{Res} : \Omega^p_{X/S}(\log Y) \to \Omega^{p - 1}_{Y/S}
$$
défini par la règle
$$
\text{Res}(\eta' + \text{d}\log(f) \wedge \eta) = \eta|_{Y \cap U}
$$
pour tout ouvert $U$ comme ci-dessus et tous
$\eta' \in \Omega^p_{X/S}(U)$ et $\eta \in \Omega^{p - 1}_{X/S}(U)$.
Si une forme $\eta$ sur $U$ se restreint à zéro sur $Y \cap U$, alors
$\eta = \text{d}f \wedge \eta' + f\eta''$ pour certaines formes $\eta'$ et $\eta''$
sur $U$. On en déduit
une suite exacte courte
$$
0 \to \Omega^p_{X/S} \to \Omega^p_{X/S}(\log Y) \to \Omega^{p - 1}_{Y/S} \to 0
$$
pour tout $p$. Il reste à définir les différentielles
$\Omega^p_{X/S}(\log Y) \to \Omega^{p + 1}_{X/S}(\log Y)$.
Sur le sous-faisceau $\Omega^p_{X/S}$, on utilise la différentielle du
complexe de de Rham de $X$ sur $S$. Enfin, on pose
$\text{d}(\text{d}\log(f) \wedge \eta) = -\text{d}\log(f) \wedge \text{d}\eta$.
Le signe est imposé par la règle de Leibniz (sur $\Omega^\bullet_{X/S}$)
et est compatible avec la différentielle sur $\Omega^\bullet_{Y/S}[-1]$,
qui est bien $-\text{d}_{Y/S}$ selon notre convention de signe dans
Homologie, définition \ref{homology-definition-shift-cochain}.
On obtient ainsi une suite exacte courte
de complexes comme dans l'énoncé du lemme.
\end{proof}

\begin{definition}
\label{definition-log-complex}
Soit $X \to S$ un morphisme de schémas. Soit $Y \subset X$ un
diviseur de Cartier effectif. Supposons que le complexe de de Rham à pôles logarithmiques
soit défini pour $Y \subset X$ sur $S$. Alors le complexe
$$
\Omega^\bullet_{X/S}(\log Y)
$$
construit dans le lemme \ref{lemma-log-complex} est le
{\it complexe de de Rham à pôles logarithmiques pour $Y \subset X$ sur $S$}.
\end{definition}

\noindent
Ce complexe possède de nombreuses propriétés utiles.

\begin{lemma}
\label{lemma-multiplication-log}
Soit $p : X \to S$ un morphisme de schémas. Soit $Y \subset X$ un
diviseur de Cartier effectif. Supposons que le complexe de de Rham à pôles logarithmiques
soit défini pour $Y \subset X$ sur $S$.
\begin{enumerate}
\item Les applications
$\wedge : \Omega^p_{X/S} \times \Omega^q_{X/S} \to \Omega^{p + q}_{X/S}$
s'étendent de manière unique en des applications $\mathcal{O}_X$-bilinéaires
$$
\wedge : \Omega^p_{X/S}(\log Y) \times \Omega^q_{X/S}(\log Y)
\to \Omega^{p + q}_{X/S}(\log Y)
$$
satisfaisant à la règle de Leibniz
$
\text{d}(\omega \wedge \eta) = \text{d}(\omega) \wedge \eta +
(-1)^{\deg(\omega)} \omega \wedge \text{d}(\eta)$,
\item pour la multiplication de (1), le morphisme
$\Omega^\bullet_{X/S} \to \Omega^\bullet_{X/S}(\log(Y)$
est un homomorphisme de $\mathcal{O}_S$-algèbres différentielles graduées,
\item via les applications de (1), on a $\Omega^p_{X/S}(\log Y) =
\wedge^p(\Omega^1_{X/S}(\log Y))$, et
\item le morphisme
$\text{Res} : \Omega^\bullet_{X/S}(\log Y) \to \Omega^\bullet_{Y/S}[-1]$
vérifie
$$
\text{Res}(\omega \wedge \eta) = \text{Res}(\omega) \wedge \eta|_Y
$$
pour toute section locale $\omega$ de $\Omega^p_{X/S}(\log Y)$ et toute section locale $\eta$
de $\Omega^q_{X/S}$.
\end{enumerate}
\end{lemma}

\begin{proof}
Cela résulte par un calcul direct de la construction locale
du complexe dans la démonstration du lemme \ref{lemma-log-complex}.
Nous omettons les détails.
\end{proof}

\noindent
Considérons un diagramme commutatif
$$
\xymatrix{
X' \ar[r]_f \ar[d] & X \ar[d] \\
S' \ar[r] & S
}
$$
de schémas. Soit $Y \subset X$ un diviseur de Cartier effectif
dont l'image inverse $Y' = f^*Y$ est définie
(Diviseurs, définition
\ref{divisors-definition-pullback-effective-Cartier-divisor}).
Supposons que
le complexe de de Rham à pôles logarithmiques soit défini pour $Y \subset X$ sur $S$
et que
le complexe de de Rham à pôles logarithmiques soit défini pour $Y' \subset X'$ sur $S'$.
On obtient alors un morphisme de suites exactes courtes de complexes
$$
\xymatrix{
0 \ar[r] &
f^{-1}\Omega^\bullet_{X/S} \ar[r] \ar[d] &
f^{-1}\Omega^\bullet_{X/S}(\log Y) \ar[r] \ar[d] &
f^{-1}\Omega^\bullet_{Y/S}[-1] \ar[r] \ar[d] &
0 \\
0 \ar[r] &
\Omega^\bullet_{X'/S'} \ar[r] &
\Omega^\bullet_{X'/S'}(\log Y') \ar[r] &
\Omega^\bullet_{Y'/S'}[-1] \ar[r] &
0
}
$$
Par linéarisation, on obtient, pour tout $p$, un morphisme linéaire
$f^*\Omega^p_{X/S}(\log Y) \to \Omega^p_{X'/S'}(\log Y')$.

\begin{lemma}
\label{lemma-gysin-via-log-complex}
Soit $f : X \to S$ un morphisme de schémas. Soit $Y \subset X$ un diviseur
de Cartier effectif. Supposons que le complexe de de Rham à pôles logarithmiques soit défini pour
$Y \subset X$ sur $S$. Notons
$$
\delta : \Omega^\bullet_{Y/S} \to \Omega^\bullet_{X/S}[2]
$$
dans $D(X, f^{-1}\mathcal{O}_S)$ le morphisme de « bord » provenant de la
suite exacte courte du lemme \ref{lemma-log-complex}. Notons
$$
\xi' : \Omega^\bullet_{X/S} \to \Omega^\bullet_{X/S}[2]
$$
dans $D(X, f^{-1}\mathcal{O}_S)$ le morphisme de la
remarque \ref{remark-cup-product-as-a-map}
correspondant à $\xi = c_1^{dR}(\mathcal{O}_X(-Y))$. Notons
$$
\zeta' : \Omega^\bullet_{Y/S} \to \Omega^\bullet_{Y/S}[2]
$$
dans $D(Y, f|_Y^{-1}\mathcal{O}_S)$ le morphisme de la
remarque \ref{remark-cup-product-as-a-map} correspondant à
$\zeta = c_1^{dR}(\mathcal{O}_X(-Y)|_Y)$. Alors le diagramme
$$
\xymatrix{
\Omega^\bullet_{X/S} \ar[d]_{\xi'} \ar[r] &
\Omega^\bullet_{Y/S} \ar[d]^{\zeta'} \ar[ld]_\delta \\
\Omega^\bullet_{X/S}[2] \ar[r] &
\Omega^\bullet_{Y/S}[2]
}
$$
est commutatif dans $D(X, f^{-1}\mathcal{O}_S)$.
\end{lemma}

\begin{proof}
Plus précisément, on définit $\delta$ comme le morphisme de bord correspondant à la
suite exacte courte décalée
$$
0 \to \Omega^\bullet_{X/S}[1] \to
\Omega^\bullet_{X/S}(\log Y)[1] \to
\Omega^\bullet_{Y/S} \to 0
$$
Il suffit de démontrer la commutativité de chaque triangle. Posons
$\mathcal{L} = \mathcal{O}_X(-Y)$. Notons $\pi : L \to X$ le fibré en droites
tel que $\pi_*\mathcal{O}_L = \bigoplus_{n \geq 0} \mathcal{L}^{\otimes n}$.

\medskip\noindent
Commutativité du triangle supérieur gauche.
D'après le lemme \ref{lemma-the-complex-for-L-star-gives-chern-class},
le morphisme $\xi'$ est le morphisme de bord du triangle donné dans le
lemme \ref{lemma-the-complex-for-L-star}.
Par fonctorialité, il suffit de prouver qu'il existe un morphisme de
suites exactes courtes
$$
\xymatrix{
0 \ar[r] &
\Omega^\bullet_{X/S}[1] \ar[r] \ar[d] &
\Omega^\bullet_{L^\star/S, 0}[1] \ar[r] \ar[d] &
\Omega^\bullet_{X/S} \ar[r] \ar[d] &
0 \\
0 \ar[r] &
\Omega^\bullet_{X/S}[1] \ar[r] &
\Omega^\bullet_{X/S}(\log Y)[1] \ar[r] &
\Omega^\bullet_{Y/S} \ar[r] &
0
}
$$
où les flèches verticales de gauche et de droite sont les flèches évidentes.
On peut définir la flèche verticale du milieu par la règle
$$
\omega' + \text{d}\log(s) \wedge \omega \longmapsto
\omega' + \text{d}\log(f) \wedge \omega
$$
où $\omega', \omega$ sont des sections locales de $\Omega^\bullet_{X/S}$,
où $s$ est un générateur local de $\mathcal{L}$ et
où $f \in \mathcal{O}_X(-Y)$ est la section correspondante du faisceau
d'idéaux de $Y$ dans $X$. Puisque les constructions des morphismes dans les
lemmes \ref{lemma-the-complex-for-L-star} et \ref{lemma-log-complex}
se correspondent exactement, cela convient.

\medskip\noindent
Commutativité du triangle inférieur droit. Notons
$\overline{L}$ la restriction de $L$ à $Y$.
D'après le lemme \ref{lemma-the-complex-for-L-star-gives-chern-class},
le morphisme $\zeta'$ est le morphisme de bord du triangle donné dans le
lemme \ref{lemma-the-complex-for-L-star}, appliqué au fibré en droites
$\overline{L}$ sur $Y$.
Par fonctorialité, il suffit de prouver qu'il existe un morphisme de
suites exactes courtes
$$
\xymatrix{
0 \ar[r] &
\Omega^\bullet_{X/S}[1] \ar[r] \ar[d] &
\Omega^\bullet_{X/S}(\log Y)[1] \ar[r] \ar[d] &
\Omega^\bullet_{Y/S} \ar[r] \ar[d] &
0 \\
0 \ar[r] &
\Omega^\bullet_{Y/S}[1] \ar[r] &
\Omega^\bullet_{\overline{L}^\star/S, 0}[1] \ar[r] &
\Omega^\bullet_{Y/S} \ar[r] &
0 \\
}
$$
où les flèches verticales de gauche et de droite sont les flèches évidentes.
On peut définir la flèche verticale du milieu par la règle
$$
\omega' + \text{d}\log(f) \wedge \omega \longmapsto
\omega'|_Y + \text{d}\log(\overline{s}) \wedge \omega|_Y
$$
où $\omega', \omega$ sont des sections locales de $\Omega^\bullet_{X/S}$,
où $f$ est un générateur local de $\mathcal{O}_X(-Y)$ considéré comme
une fonction sur $X$, et où $\overline{s}$ désigne $f|_Y$ considéré comme une
section de $\mathcal{L}|_Y = \mathcal{O}_X(-Y)|_Y$.
Puisque les constructions des morphismes dans les
lemmes \ref{lemma-the-complex-for-L-star} et \ref{lemma-log-complex}
se correspondent exactement, cela convient.
\end{proof}

\begin{lemma}
\label{lemma-log-complex-consequence}
Soit $X \to S$ un morphisme de schémas. Soit $Y \subset X$ un diviseur
de Cartier effectif. Supposons que le complexe de de Rham à pôles logarithmiques soit défini pour
$Y \subset X$ sur $S$. Soit $b \in H^m_{dR}(X/S)$ une classe de cohomologie de de Rham
dont la restriction à $Y$ est nulle. Alors
$c_1^{dR}(\mathcal{O}_X(Y)) \cup b = 0$ dans $H^{m + 2}_{dR}(X/S)$.
\end{lemma}

\begin{proof}
Cela résulte immédiatement du lemme \ref{lemma-gysin-via-log-complex}.
En effet, on a
$$
c_1^{dR}(\mathcal{O}_X(Y)) \cup b =
-c_1^{dR}(\mathcal{O}_X(-Y)) \cup b = -\xi'(b) = -\delta(b|_Y) = 0
$$
comme voulu. Pour la deuxième égalité, voir
la remarque \ref{remark-cup-product-as-a-map}.
\end{proof}

\begin{lemma}
\label{lemma-check-log-smooth}
Soient $X \to T \to S$ des morphismes de schémas. Soit $Y \subset X$ un diviseur
de Cartier effectif. Si $X \to T$ et $Y \to T$ sont tous deux lisses, alors
le complexe de de Rham à pôles logarithmiques est défini pour $Y \subset X$ sur $S$.
\end{lemma}

\begin{proof}
Soit $y \in Y$ un point.
D'après Compléments sur les morphismes, lemme \ref{more-morphisms-lemma-etale-local-structure},
il existe un entier $0 \geq m$ et un diagramme commutatif
$$
\xymatrix{
Y \ar[d] &
V \ar[l] \ar[d] \ar[r] &
\mathbf{A}^m_T
\ar[d]^{(a_1, \ldots, a_m) \mapsto (a_1, \ldots, a_m, 0)} \\
X &
U \ar[l] \ar[r]^-\pi &
\mathbf{A}^{m + 1}_T
}
$$
où $U \subset X$ est ouvert, $V = Y \cap U$,
$\pi$ est étale, $V = \pi^{-1}(\mathbf{A}^m_T)$ et $y \in V$.
Notons $z \in \mathbf{A}^m_T$ l'image de $y$. On a alors
$$
\Omega^p_{X/S, y} = \Omega^p_{\mathbf{A}^{m + 1}_T/S, z}
\otimes_{\mathcal{O}_{\mathbf{A}^{m + 1}_T, z}} \mathcal{O}_{X, x}
$$
d'après le lemme \ref{lemma-etale}. Notons $x_1, \ldots, x_{m + 1}$
les fonctions coordonnées sur $\mathbf{A}^{m + 1}_T$.
Puisque les conditions (1) et (2) de la définition \ref{definition-local-product}
ne dépendent pas du choix de la coordonnée locale,
il suffit de les vérifier lorsque $f$ est
l'image de $x_{m + 1}$ par l'homomorphisme local plat d'anneaux
$\mathcal{O}_{\mathbf{A}^{m + 1}_T, z} \to \mathcal{O}_{X, x}$.
Il suffit donc de vérifier les conditions (1) et (2)
pour $\mathbf{A}^m_T \subset \mathbf{A}^{m + 1}_T$ et le point $z$.
Pour traiter ce cas, on peut supposer $S = \Spec(A)$ et $T = \Spec(B)$
affines. Soit $A \to B$ le morphisme d'anneaux correspondant au morphisme
$T \to S$, et posons $P = B[x_1, \ldots, x_{m + 1}]$, de sorte que
$\mathbf{A}^{m + 1}_T = \Spec(P)$. On a
$$
\Omega_{P/A} = \Omega_{B/A} \otimes_B P \oplus
\bigoplus\nolimits_{j = 1, \ldots, m} P \text{d}x_j \oplus
P \text{d}x_{m + 1}
$$
Ainsi, l'application $P \to \Omega_{P/A}$, $g \mapsto g \text{d}x_{m + 1}$,
est une injection scindée, et $x_{m + 1}$ est non diviseur de zéro sur
$\Omega^p_{P/A}$ pour tout $p \geq 0$. La localisation en l'idéal premier
correspondant à $z$ achève la démonstration.
\end{proof}

\begin{remark}
\label{remark-check-log-completion-1}
Soit $S$ un schéma localement noethérien. Soit $X$ localement de type
fini sur $S$. Soit $Y \subset X$ un diviseur de Cartier effectif.
Si le morphisme
$$
\mathcal{O}_{X, y}^\wedge \longrightarrow \mathcal{O}_{Y, y}^\wedge
$$
admet une section pour tout $y \in Y$, alors
le complexe de de Rham à pôles logarithmiques est défini pour $Y \subset X$ sur $S$.
Si nous avons besoin de ce résultat, nous en formulerons un énoncé précis et
en ajouterons ici une démonstration.
\end{remark}

\begin{remark}
\label{remark-check-log-completion-2}
Soit $S$ un schéma localement noethérien. Soit $X$ localement de type
fini sur $S$. Soit $Y \subset X$ un diviseur de Cartier effectif.
Si, pour tout $y \in Y$, on peut trouver un diagramme de schémas sur $S$
$$
X \xleftarrow{\varphi} U \xrightarrow{\psi} V
$$
où $\varphi$ est étale et $\psi|_{\varphi^{-1}(Y)} : \varphi^{-1}(Y) \to V$
est étale, alors le complexe de de Rham à pôles logarithmiques est défini pour
$Y \subset X$ sur $S$. C'est notamment le cas lorsque le couple $(X, Y)$
est, étale-localement, de la forme $(V \times \mathbf{A}^1, V \times \{0\})$.
Si nous avons besoin de ce résultat, nous en formulerons
un énoncé précis et en ajouterons ici une démonstration.
\end{remark}

















\section{Calculs}
\label{section-calculations}

\noindent
Dans cette section, nous calculons certains groupes de cohomologie de Hodge
et de de Rham pour un éclatement standard.

\medskip\noindent
Fixons un anneau $R$ et posons $S = \Spec(R)$. Fixons des entiers $0 \leq m$ et
$1 \leq n$. Considérons l'immersion fermée
$$
Z = \mathbf{A}^m_S \longrightarrow \mathbf{A}^{m + n}_S = X,\quad
(a_1, \ldots, a_m) \mapsto (a_1, \ldots, a_m, 0, \ldots 0).
$$
Nous allons considérer l'éclatement $L$ de $X$
le long du sous-schéma fermé $Z$. Écrivons
$$
X =
\mathbf{A}^{m + n}_S =
\Spec(A)
\quad\text{avec}\quad
A = R[x_1, \ldots, x_m, y_1, \ldots, y_n]
$$
Nous considérerons $A = R[x_1, \ldots, x_m, y_1, \ldots, y_n]$ comme une
$R$-algèbre graduée en posant $\deg(x_i) = 0$ et $\deg(y_j) = 1$.
Pour cette graduation, on a
$$
P =
\text{Proj}(A) =
\mathbf{A}^m_S \times_S \mathbf{P}^{n - 1}_S =
Z \times_S \mathbf{P}^{n - 1}_S =
\mathbf{P}^{n - 1}_Z
$$
Observons que l'idéal définissant $Z$ dans $X$ est $A_+$.
Ainsi, $L$ est le Proj de l'algèbre de Rees
$$
A \oplus A_+ \oplus (A_+)^2 \oplus \ldots =
\bigoplus\nolimits_{d \geq 0} A_{\geq d}
$$
Par conséquent, $L$ est un exemple du phénomène étudié
plus généralement dans Compléments sur les morphismes, section
\ref{more-morphisms-section-proj-spec} ;
nous utiliserons les observations qui y ont été faites sans plus les mentionner.
En particulier, on dispose d'un diagramme commutatif
$$
\xymatrix{
P \ar[r]_0 \ar[d]_p &
L \ar[r]_-\pi \ar[d]^b &
P \ar[d]^p \\
Z \ar[r]^i &
X \ar[r] &
Z
}
$$
tel que $\pi : L \to P$ soit un fibré en droites sur
$P = Z \times_S \mathbf{P}^{n - 1}_S$,
de section nulle $0$, dont l'image $E = 0(P) \subset L$
est le diviseur exceptionnel de l'éclatement $b$.

\begin{lemma}
\label{lemma-comparison}
Pour $a \geq 0$, on a
\begin{enumerate}
\item le morphisme
$\Omega^a_{X/S} \to b_*\Omega^a_{L/S}$ est un isomorphisme,
\item le morphisme $\Omega^a_{Z/S} \to p_*\Omega^a_{P/S}$ est un isomorphisme,
et
\item le morphisme $Rb_*\Omega^a_{L/S} \to i_*Rp_*\Omega^a_{P/S}$ induit un isomorphisme
sur les faisceaux de cohomologie en degrés $\geq 1$.
\end{enumerate}
\end{lemma}

\begin{proof}
Démontrons d'abord (2). Puisque
$P = Z \times_S \mathbf{P}^{n - 1}_S$,
on voit que
$$
\Omega^a_{P/S} = \bigoplus\nolimits_{a = r + s}
\text{pr}_1^*\Omega^r_{Z/S} \otimes
\text{pr}_2^*\Omega^s_{\mathbf{P}^{n - 1}_S/S}
$$
En rappelant que $p = \text{pr}_1$, la formule de projection
(Cohomologie, lemme \ref{cohomology-lemma-projection-formula})
donne
$$
p_*\Omega^a_{P/S} = \bigoplus\nolimits_{a = r + s}
\Omega^r_{Z/S} \otimes
\text{pr}_{1, *}\text{pr}_2^*\Omega^s_{\mathbf{P}^{n - 1}_S/S}
$$
D'après les calculs de la section \ref{section-projective-space},
en particulier ceux de la
démonstration du lemme \ref{lemma-hodge-cohomology-projective-space},
on a $\text{pr}_{1, *}\text{pr}_2^*\Omega^s_{\mathbf{P}^{n - 1}_S/S} = 0$,
sauf si $s = 0$, auquel cas on obtient
$\text{pr}_{1, *}\mathcal{O}_P = \mathcal{O}_Z$.
Cela démontre (2).

\medskip\noindent
D'après les résultats de la section \ref{section-line-bundle}, en particulier
le lemme \ref{lemma-push-omega-a}, on a
$\pi_*\Omega^a_{L/S} = \Omega^a_{P/S} \oplus
\bigoplus_{k \geq 1} \Omega^a_{L/S, k}$.
Puisque la composée $\pi \circ 0$ dans le diagramme ci-dessus
est l'identité de $P$, il suffit, pour démontrer (3), de montrer que
les groupes de cohomologie supérieure de $\Omega^a_{L/S, k}$ sont nuls pour $k > 0$.
D'après les lemmes \ref{lemma-the-complex-for-L-star} et \ref{lemma-push-omega-a},
on dispose de suites exactes courtes
$$
0 \to \Omega^a_{P/S} \otimes \mathcal{O}_P(k)
\to \Omega^a_{L/S, k} \to
\Omega^{a - 1}_{P/S} \otimes \mathcal{O}_P(k) \to 0
$$,
où $\Omega^{a - 1}_{P/S} = 0$ si $a = 0$. Puisque
$P = Z \times_S \mathbf{P}^{n - 1}_S$, on a
$$
\Omega^a_{P/S} = \bigoplus\nolimits_{i + j = a}
\Omega^i_{Z/S} \boxtimes \Omega^j_{\mathbf{P}^{n - 1}_S/S}
$$
d'après le lemme \ref{lemma-de-rham-complex-product}.
Comme $\Omega^i_{Z/S}$ est libre de rang fini,
on voit qu'il suffit de montrer que les groupes de cohomologie supérieure de
$\mathcal{O}_Z \boxtimes \Omega^j_{\mathbf{P}^{n - 1}_S/S}(k)$
sont nuls pour $k > 0$. Cela résulte du
lemme \ref{lemma-twisted-hodge-cohomology-projective-space}
appliqué à $P = Z \times_S \mathbf{P}^{n - 1}_S = \mathbf{P}^{n - 1}_Z$,
ce qui achève la démonstration de (3).

\medskip\noindent
Il reste à démontrer (1). Si $n = 1$, on éclate
un diviseur de Cartier effectif ; $b$ est donc un isomorphisme,
d'où (1). Si $n > 1$, la composée
$$
\Gamma(X, \Omega^a_{X/S})
\to
\Gamma(L, \Omega^a_{L/S})
\to
\Gamma(L \setminus E, \Omega^a_{L/S})
=
\Gamma(X \setminus Z, \Omega^a_{X/S})
$$
est un isomorphisme, puisque $\Omega^a_{X/S}$ est libre de rang fini
(nous omettons ce détail). Ainsi, (1) ne peut être faux que s'il
existe des éléments non nuls de $\Gamma(L, \Omega^a_{L/S})$ qui s'annulent
en dehors de $E = 0(P)$. Puisque $L$ est un fibré en droites sur $P$
de section nulle $0 : P \to L$, il suffit de montrer que,
sur un fibré en droites, il n'existe pas de section non nulle d'un faisceau
de différentielles qui s'annule identiquement hors de la section nulle.
Le lecteur le vérifiera soit, de préférence, par un calcul local,
soit en utilisant $\Omega_{L/S, k} \subset \Omega_{L^\star/S, k}$
(voir les références ci-dessus).
\end{proof}

\noindent
Nous suggérons au lecteur de passer maintenant à la section suivante.

\begin{lemma}
\label{lemma-comparison-bis}
Pour $a \geq 0$, il existe des morphismes canoniques
$$
b^*\Omega^a_{X/S} \longrightarrow
\Omega^a_{L/S} \longrightarrow
b^*\Omega^a_{X/S} \otimes_{\mathcal{O}_L} \mathcal{O}_L((n - 1)E)
$$
dont la composée est induite par l'inclusion
$\mathcal{O}_L \subset \mathcal{O}_L((n - 1)E)$.
\end{lemma}

\begin{proof}
La première flèche de la formule affichée est
décrite à la section \ref{section-de-rham-complex}.
Pour obtenir la deuxième flèche, il faut montrer que, si l'on considère
une section locale de $\Omega^a_{L/S}$ comme une « section méromorphe »
de $b^*\Omega^a_{X/S}$, elle a un pôle d'ordre au plus
$n - 1$ le long de $E$.
Pour le voir, travaillons dans des cartes affines locales
de $L$. Rappelons que $L$ est recouvert par les spectres des
algèbres d'éclatement affines $A[\frac{I}{y_i}]$, où $I = A_{+}$
est l'idéal engendré par $y_1, \ldots, y_n$. Voir
Algèbre, section \ref{algebra-section-blow-up} et
Diviseurs, lemme \ref{divisors-lemma-blowing-up-affine}.
Par symétrie, il suffit de travailler dans la
carte correspondant à $i = 1$. Alors
$$
A[\frac{I}{y_1}] = R[x_1, \ldots, x_m, y_1, t_2, \ldots, t_n]
$$
où $t_i = y_i/y_1$ ; voir
Compléments d'algèbre, lemme \ref{more-algebra-lemma-blowup-regular-sequence}.
Ainsi, le module $\Omega^1_{L/S}$ est librement engendré, sur l'ouvert
affine correspondant, par
$\text{d}x_1, \ldots, \text{d}x_m$, $\text{d}y_1$ et
$\text{d}t_1, \ldots, \text{d}t_n$.
Bien entendu, les $m + 1$ premiers de ces générateurs proviennent de
$b^*\Omega^1_{X/S}$, et pour les $n - 1$ restants, on a
$$
\text{d}t_j =
\text{d}\frac{y_j}{y_1} =
\frac{1}{y_1}\text{d}y_j - \frac{y_j}{y_1^2}\text{d}y_1 =
\frac{\text{d}y_j - t_j \text{d}y_1}{y_1}
$$,
qui a un pôle d'ordre $1$ le long de $E$, puisque $E$ est défini par $y_1$
sur cette carte. Les produits extérieurs de $a$ de ces éléments forment une base
de $\Omega^a_{L/S}$ sur cette carte, et au plus
$n - 1$ des $\text{d}t_j$ interviennent, ce qui achève la démonstration.
\end{proof}

\begin{lemma}
\label{lemma-blowup-twist-same-cohomology}
Soit $E = 0(P)$ le diviseur exceptionnel de l'éclatement $b$.
Pour tout $\mathcal{O}_X$-module localement libre $\mathcal{E}$ et
tout $0 \leq i \leq n - 1$, le morphisme
$$
\mathcal{E}
\longrightarrow
Rb_*(b^*\mathcal{E} \otimes_{\mathcal{O}_L} \mathcal{O}_L(iE))
$$
est un isomorphisme dans $D(\mathcal{O}_X)$.
\end{lemma}

\begin{proof}
D'après la formule de projection, il suffit de le montrer pour
$\mathcal{E} = \mathcal{O}_X$ ; voir Cohomologie, lemme
\ref{cohomology-lemma-projection-formula}.
Puisque $X$ est affine, il suffit de montrer que les morphismes
$$
H^0(X, \mathcal{O}_X) \to
H^0(L, \mathcal{O}_L) \to
H^0(L, \mathcal{O}_L(iE))
$$
sont des isomorphismes et que $H^j(X, \mathcal{O}_L(iE)) = 0$
pour $j > 0$ et $0 \leq i \leq n - 1$ ; voir Cohomologie des schémas, lemme
\ref{coherent-lemma-quasi-coherence-higher-direct-images-application}.
Puisque $\pi$ est affine, on peut calculer les sections globales et la
cohomologie après application de $\pi_*$ ; voir Cohomologie des schémas, lemme
\ref{coherent-lemma-relative-affine-cohomology}. Si $n = 1$,
$L \to X$ est un isomorphisme et $i = 0$, d'où la première assertion.
Si $n > 1$, considérons la composée
$$
H^0(X, \mathcal{O}_X) \to H^0(L, \mathcal{O}_L) \to
H^0(L, \mathcal{O}_L(iE)) \to H^0(L \setminus E, \mathcal{O}_L) =
H^0(X \setminus Z, \mathcal{O}_X)
$$
Puisque
$H^0(X \setminus Z, \mathcal{O}_X) = H^0(X, \mathcal{O}_X)$ dans ce
cas, car $Z$ est de codimension $n \geq 2$ dans $X$ (détails omis), on en déduit
la première assertion. Pour la seconde, rappelons que
$\mathcal{O}_L(E) = \mathcal{O}_L(-1)$ ; voir Diviseurs, lemme
\ref{divisors-lemma-blowing-up-gives-effective-Cartier-divisor}.
On a donc
$$
\pi_*\mathcal{O}_L(iE) =
\pi_*\mathcal{O}_L(-i) =
\bigoplus\nolimits_{k \geq -i} \mathcal{O}_P(k)
$$,
comme expliqué dans
Compléments sur les morphismes, section \ref{more-morphisms-section-proj-spec}.
On conclut par l'annulation des groupes de cohomologie des faisceaux tordus
du faisceau structural sur $P = \mathbf{P}^{n - 1}_Z$,
établie dans Cohomologie des schémas, lemme
\ref{coherent-lemma-cohomology-projective-space-over-ring}.
\end{proof}























\section{Éclatements et cohomologie de de Rham}
\label{section-blowing-up}

\noindent
Fixons un schéma de base $S$, un morphisme lisse $X \to S$ et un sous-schéma fermé
$Z \subset X$, lui aussi lisse sur $S$. Notons $b : X' \to X$
l'éclatement de $X$ le long de $Z$. Notons $E \subset X'$ le diviseur
exceptionnel. Voici le diagramme
\begin{equation}
\label{equation-blowup}
\vcenter{
\xymatrix{
E \ar[r]_j \ar[d]_p & X' \ar[d]^b \\
Z \ar[r]^i & X
}
}
\end{equation}
Notre but dans cette section est de démontrer que l'application
$b^* : H_{dR}^*(X/S) \longrightarrow H_{dR}^*(X'/S)$
est injective (bien que l'on puisse en dire beaucoup plus).

\begin{lemma}
\label{lemma-comparison-general}
Avec les notations de Compléments sur les morphismes, lemme \ref{more-morphisms-lemma-blowup},
pour $a \geq 0$, on a
\begin{enumerate}
\item le morphisme
$\Omega^a_{X/S} \to b_*\Omega^a_{X'/S}$ est un isomorphisme,
\item le morphisme $\Omega^a_{Z/S} \to p_*\Omega^a_{E/S}$ est un isomorphisme,
\item le morphisme $Rb_*\Omega^a_{X'/S} \to i_*Rp_*\Omega^a_{E/S}$ induit un isomorphisme
sur les faisceaux de cohomologie en degrés $\geq 1$.
\end{enumerate}
\end{lemma}

\begin{proof}
Soit $\epsilon : X_1 \to X$ un morphisme étale surjectif. Notons
$i_1 : Z_1 \to X_1$, $b_1 : X'_1 \to X_1$, $E_1 \subset X'_1$ et
$p_1 : E_1 \to Z_1$ les changements de base des objets considérés dans
Compléments sur les morphismes, lemme \ref{more-morphisms-lemma-blowup}.
Observons que $i_1$ est une immersion fermée
de schémas lisses sur $S$, et que $b_1$ est l'éclatement de centre
$Z_1$, d'après Diviseurs, lemme \ref{divisors-lemma-flat-base-change-blowing-up}.
Supposons que l'on puisse démontrer (1), (2) et (3)
pour les morphismes $b_1$, $p_1$ et $i_1$. Alors, d'après
le lemme \ref{lemma-etale}, les images inverses par $\epsilon$
des morphismes de (1), (2) et (3) sont des isomorphismes. Puisque $\epsilon$
est un morphisme plat surjectif, on conclut.
Ainsi, en travaillant étale-localement et en utilisant
Compléments sur les morphismes, lemme \ref{more-morphisms-lemma-etale-local-structure},
on peut supposer être dans la situation étudiée à la
section \ref{section-calculations}. Le lemme est alors
identique au lemme \ref{lemma-comparison}.
\end{proof}

\begin{lemma}
\label{lemma-distinguished-triangle-blowup}
Avec les notations de Compléments sur les morphismes, lemme \ref{more-morphisms-lemma-blowup},
et en notant $f : X \to S$ le morphisme structural, il existe un
triangle distingué canonique
$$
\Omega^\bullet_{X/S} \to
Rb_*(\Omega^\bullet_{X'/S}) \oplus i_*\Omega^\bullet_{Z/S} \to
i_*Rp_*(\Omega^\bullet_{E/S}) \to
\Omega^\bullet_{X/S}[1]
$$
dans $D(X, f^{-1}\mathcal{O}_S)$, où les quatre morphismes
$$
\begin{matrix}
\Omega^\bullet_{X/S} & \to & Rb_*(\Omega^\bullet_{X'/S}), \\
\Omega^\bullet_{X/S} & \to & i_*\Omega^\bullet_{Z/S}, \\
Rb_*(\Omega^\bullet_{X'/S}) & \to & i_*Rp_*(\Omega^\bullet_{E/S}), \\
i_*\Omega^\bullet_{Z/S} & \to & i_*Rp_*(\Omega^\bullet_{E/S})
\end{matrix}
$$
sont les morphismes canoniques (section \ref{section-de-rham-complex}),
au changement de signe près de l'un d'entre eux.
\end{lemma}

\begin{proof}
Choisissons un triangle distingué
$$
C \to Rb_*\Omega^\bullet_{X'/S} \oplus i_*\Omega^\bullet_{Z/S}
\to i_*Rp_*\Omega^\bullet_{E/S} \to C[1]
$$
dans $D(X, f^{-1}\mathcal{O}_S)$. Il suffit de montrer que
$\Omega^\bullet_{X/S}$ est isomorphe à $C$ d'une manière compatible
avec les morphismes canoniques. Les axiomes des catégories triangulées
fournissent un morphisme de triangles distingués
$$
\xymatrix{
C' \ar[r] \ar[d] &
b_*\Omega^\bullet_{X'/S} \oplus i_*\Omega^\bullet_{Z/S} \ar[r] \ar[d] &
i_*p_*\Omega^\bullet_{E/S} \ar[r] \ar[d] &
C'[1] \ar[d] \\
C \ar[r] &
Rb_*\Omega^\bullet_{X'/S} \oplus i_*\Omega^\bullet_{Z/S} \ar[r] &
i_*Rp_*\Omega^\bullet_{E/S} \ar[r] &
C[1]
}
$$
D'après le lemme \ref{lemma-comparison-general}, partie (3), et
Catégories dérivées, proposition \ref{derived-proposition-9}, on en déduit que
$C' \to C$ est un isomorphisme. D'après le lemme \ref{lemma-comparison-general}, partie (2),
le morphisme $i_*\Omega^\bullet_{Z/S} \to i_*p_*\Omega^\bullet_{E/S}$
est un isomorphisme. Ainsi, $C' = b_*\Omega^\bullet_{X'/S}$
dans la catégorie dérivée. Enfin, le lemme \ref{lemma-comparison-general},
partie (1), montre que cet objet est égal à $\Omega^\bullet_{X/S}$.
Nous omettons la vérification de la compatibilité avec les morphismes canoniques.
\end{proof}

\begin{proposition}
\label{proposition-blowup-split}
Avec les notations de Compléments sur les morphismes, lemme \ref{more-morphisms-lemma-blowup},
le morphisme $\Omega^\bullet_{X/S} \to Rb_*\Omega^\bullet_{X'/S}$
admet une rétraction dans $D(X, (X \to S)^{-1}\mathcal{O}_S)$.
\end{proposition}

\begin{proof}
Considérons le triangle construit dans le
lemme \ref{lemma-distinguished-triangle-blowup}.
Nous affirmons que le morphisme
$$
Rb_*(\Omega^\bullet_{X'/S}) \oplus i_*\Omega^\bullet_{Z/S} \to
i_*Rp_*(\Omega^\bullet_{E/S})
$$
admet une section dont l'image contient le facteur direct $i_*\Omega^\bullet_{Z/S}$.
D'après Catégories dérivées, lemme \ref{derived-lemma-split}, cela montrera que
la première flèche du triangle admet une rétraction qui s'annule sur
le facteur direct $i_*\Omega^\bullet_{Z/S}$, ce qui démontre le lemme.
Pour prouver l'assertion, nous décomposerons $Rp_*\Omega^\bullet_{E/S}$
en une somme directe dont le premier facteur correspond à
$\Omega^\bullet_{Z/S}$ et dont le second se relève
par $Rb_*\Omega^\bullet_{X'/S}$.

\medskip\noindent
Démonstration de l'assertion. On peut décomposer $X$ en sous-schémas ouverts et fermés
de dimension relative constante sur $S$ ; voir
Morphismes, lemme \ref{morphisms-lemma-smooth-omega-finite-locally-free}.
Puisque la catégorie dérivée $D(X, f^{-1}\mathcal{O})_S)$ se décompose alors
en un produit de catégories, on peut supposer $X$ de
dimension relative constante $N$ sur $S$. On peut décomposer
$Z = \coprod Z_m$ en sous-schémas ouverts et fermés de dimension
relative $m \geq 0$ sur $S$. La restriction $i_m : Z_m \to X$ de
$i$ à $Z_m$ est une immersion régulière de codimension $N - m$ ; voir Diviseurs, lemme
\ref{divisors-lemma-immersion-smooth-into-smooth-regular-immersion}.
Soit $E = \coprod E_m$ la décomposition correspondante, c'est-à-dire
que l'on pose $E_m = p^{-1}(Z_m)$. Si $p_m : E_m \to Z_m$ désigne la
restriction de $p$ à $E_m$, on dispose d'un isomorphisme canonique
$$
\tilde \xi_m :
\bigoplus\nolimits_{t = 0, \ldots, N - m - 1}
\Omega^\bullet_{Z_m/S}[-2t]
\longrightarrow
Rp_{m, *}\Omega^\bullet_{E_m/S}
$$
dans $D(Z_m, (Z_m \to S)^{-1}\mathcal{O}_S)$,
où, en degré $0$, on a le morphisme canonique
$\Omega^\bullet_{Z_m/S} \to Rp_{m, *}\Omega^\bullet_{E_m/S}$.
Voir la remarque \ref{remark-projective-space-bundle-formula}.
On dispose donc d'un isomorphisme
$$
\tilde \xi :
\bigoplus\nolimits_m
\bigoplus\nolimits_{t = 0, \ldots, N - m - 1}
\Omega^\bullet_{Z_m/S}[-2t]
\longrightarrow
Rp_*(\Omega^\bullet_{E/S})
$$
dans $D(Z, (Z \to S)^{-1}\mathcal{O}_S)$
dont la restriction au facteur direct
$\Omega^\bullet_{Z/S} = \bigoplus \Omega^\bullet_{Z_m/S}$ de la source
est le morphisme canonique $\Omega^\bullet_{Z/S} \to Rp_*(\Omega^\bullet_{E/S})$.
Considérons les sous-complexes $M_m$ et $K_m$ du complexe
$\bigoplus\nolimits_{t = 0, \ldots, N - m - 1} \Omega^\bullet_{Z_m/S}[-2t]$
introduits dans la remarque \ref{remark-projective-space-bundle-formula}.
Posons
$$
M = \bigoplus M_m
\quad\text{et}\quad
K = \bigoplus K_m
$$
On a $M = K[-2]$ et, par construction, le morphisme
$$
c_{E/Z} \oplus \tilde \xi|_M :
\Omega^\bullet_{Z/S} \oplus M
\longrightarrow
Rp_*(\Omega^\bullet_{E/S})
$$
est un isomorphisme (voir la remarque citée ci-dessus).

\medskip\noindent
Considérons le morphisme
$$
\delta : \Omega^\bullet_{E/S}[-2] \longrightarrow \Omega^\bullet_{X'/S}
$$
dans $D(X', (X' \to S)^{-1}\mathcal{O}_S)$ du
lemme \ref{lemma-gysin-via-log-complex},
qui a la propriété que la composée
$$
\Omega^\bullet_{E/S}[-2] \longrightarrow \Omega^\bullet_{X'/S}
\longrightarrow
\Omega^\bullet_{E/S}
$$
est le morphisme $\theta'$ de la remarque \ref{remark-cup-product-as-a-map} pour
$c_1^{dR}(\mathcal{O}_{X'}(-E))|_E) = c_1^{dR}(\mathcal{O}_E(1))$.
La dernière assertion de la remarque \ref{remark-projective-space-bundle-formula}
montre que le diagramme
$$
\xymatrix{
K[-2] \ar[d]_{(\tilde \xi|_K)[-2]} \ar[r]_{\text{id}} &
M \ar[d]^{\tilde x|_M} \\
Rp_*\Omega^\bullet_{E/S}[-2] \ar[r]^-{Rp_*\theta'} &
Rp_*\Omega^\bullet_{E/S}
}
$$
est commutatif. La section cherchée dans notre
assertion est donc le morphisme
\begin{align*}
Rp_*(\Omega^\bullet_{E/S})
& \xrightarrow{(c_{E/Z} \oplus \tilde \xi|_M)^{-1}}
\Omega^\bullet_{Z/S} \oplus M \\
& \xrightarrow{\text{id} \oplus \text{id}^{-1}}
\Omega^\bullet_{Z/S} \oplus K[-2] \\
& \xrightarrow{\text{id} \oplus (\tilde \xi|_K)[-2]}
\Omega^\bullet_{Z/S} \oplus Rp_*\Omega^\bullet_{E/S}[-2] \\
& \xrightarrow{\text{id} \oplus Rb_*\delta}
\Omega^\bullet_{Z/S} \oplus Rb_*\Omega^\bullet_{X'/S}
\end{align*}
La relation entre $\theta'$ et $\delta$ énoncée ci-dessus,
ainsi que le diagramme commutatif faisant intervenir $\theta'$, $\tilde \xi|_K$
et $\tilde \xi|_M$ ci-dessus, donnent exactement ce qu'il faut
pour montrer qu'il s'agit d'une section du morphisme canonique
$\Omega^\bullet_{Z/S} \oplus Rb_*(\Omega^\bullet_{X'/S}) \to
Rp_*(\Omega^\bullet_{E/S})$, ce qui achève la démonstration de l'assertion.
\end{proof}

\noindent
Le lemme \ref{lemma-splitting-on-omega-a}
montre qu'il est nettement plus facile de construire la rétraction en cohomologie
de Hodge que d'établir le résultat de la
proposition \ref{proposition-blowup-split}.
Nous invitons vivement le lecteur à passer à la section suivante.

\begin{lemma}
\label{lemma-ext-zero}
Soit $i : Z \to X$ une immersion fermée de schémas, régulière de
codimension $c$. Alors $\Ext^q_{\mathcal{O}_X}(i_*\mathcal{F}, \mathcal{E}) = 0$
pour $q < c$, lorsque $\mathcal{E}$ est localement libre sur $X$ et $\mathcal{F}$
est un $\mathcal{O}_Z$-module quelconque.
\end{lemma}

\begin{proof}
D'après la suite spectrale du local au global pour $\Ext$, il suffit
de le démontrer localement sur les ouverts affines de $X$. Voir
Cohomologie, section \ref{cohomology-section-ext}.
On peut donc supposer $X = \Spec(A)$
et l'existence d'une suite régulière $f_1, \ldots, f_c$ dans $A$
telle que $Z = \Spec(A/(f_1, \ldots, f_c))$. On peut supposer $c \geq 1$.
On voit alors que $f_1 : \mathcal{E} \to \mathcal{E}$
est injectif. Puisque $i_*\mathcal{F}$ est annulé par $f_1$,
cela montre que le lemme est vrai pour $i = 0$ et que l'on dispose
d'une surjection
$$
\Ext^{q - 1}_{\mathcal{O}_X}(i_*\mathcal{F}, \mathcal{E}/f_1\mathcal{E})
\longrightarrow
\Ext^q_{\mathcal{O}_X}(i_*\mathcal{F}, \mathcal{E})
$$
Il suffit donc de montrer que la source de cette flèche est nulle.
Répétons l'argument : si $c \geq 2$, le morphisme
$f_2 : \mathcal{E}/f_1\mathcal{E} \to \mathcal{E}/f_1\mathcal{E}$
est injectif. Puisque $i_*\mathcal{F}$ est annulé par $f_2$,
cela montre que le lemme est vrai pour $q = 1$ et que l'on dispose d'une
surjection
$$
\Ext^{q - 2}_{\mathcal{O}_X}(i_*\mathcal{F},
\mathcal{E}/f_1\mathcal{E} + f_2\mathcal{E})
\longrightarrow
\Ext^{q - 1}_{\mathcal{O}_X}(i_*\mathcal{F}, \mathcal{E}/f_1\mathcal{E})
$$
En poursuivant de cette manière, on démontre le lemme.
\end{proof}

\begin{lemma}
\label{lemma-splitting-on-omega-a}
Avec les notations de Compléments sur les morphismes, lemme \ref{more-morphisms-lemma-blowup},
pour $a \geq 0$, il existe une unique flèche
$Rb_*\Omega^a_{X'/S} \to \Omega^a_{X/S}$ dans $D(\mathcal{O}_X)$
dont la composée avec $\Omega^a_{X/S} \to Rb_*\Omega^a_{X'/S}$
est l'identité de $\Omega^a_{X/S}$.
\end{lemma}

\begin{proof}
On peut décomposer $X$ en sous-schémas ouverts et fermés
de dimension relative constante sur $S$ ; voir
Morphismes, lemme \ref{morphisms-lemma-smooth-omega-finite-locally-free}.
Puisque la catégorie dérivée $D(X, f^{-1}\mathcal{O})_S)$ se décompose alors
en un produit de catégories, on peut supposer $X$ de
dimension relative constante $N$ sur $S$. On peut décomposer
$Z = \coprod Z_m$ en sous-schémas ouverts et fermés de dimension
relative $m \geq 0$ sur $S$. La restriction $i_m : Z_m \to X$ de
$i$ à $Z_m$ est une immersion régulière de codimension $N - m$ ; voir Diviseurs, lemme
\ref{divisors-lemma-immersion-smooth-into-smooth-regular-immersion}.
Soit $E = \coprod E_m$ la décomposition correspondante, c'est-à-dire
que l'on pose $E_m = p^{-1}(Z_m)$. Nous affirmons qu'il existe des morphismes naturels
$$
b^*\Omega^a_{X/S} \to \Omega^a_{X'/S} \to
b^*\Omega^a_{X/S} \otimes_{\mathcal{O}_{X'}}
\mathcal{O}_{X'}(\sum (N - m - 1)E_m)
$$
dont la composée est induite par l'inclusion
$\mathcal{O}_{X'} \to \mathcal{O}_{X'}(\sum (N - m - 1)E_m)$.
Pour le démontrer, il suffit de montrer que le
conoyau de la première flèche est, localement sur $X'$, annulé par
une équation locale du diviseur de Cartier effectif $\sum (N - m - 1)E_m$.
Pour cela, on peut travailler étale-localement sur $X$ comme dans la
démonstration du lemme \ref{lemma-comparison-general}, et appliquer
le lemme \ref{lemma-comparison-bis}.
En calculant étale-localement à l'aide du lemme \ref{lemma-blowup-twist-same-cohomology},
on voit que la composée induite
$$
\Omega^a_{X/S} \to Rb_*\Omega^a_{X'/S} \to
Rb_*\left(b^*\Omega^a_{X/S} \otimes_{\mathcal{O}_{X'}}
\mathcal{O}_{X'}(\sum (N - m - 1)E_m)\right)
$$
est un isomorphisme dans $D(\mathcal{O}_X)$,
d'où l'existence du morphisme du lemme.

\medskip\noindent
Pour l'unicité, il suffit de montrer qu'il n'existe pas de morphisme non nul de
$\tau_{\geq 1}Rb_*\Omega_{X'/S}$ vers $\Omega^a_{X/S}$ dans $D(\mathcal{O}_X)$.
Il suffit pour cela de montrer qu'il n'existe pas de morphisme non nul
de $R^qb_*\Omega_{X'/s}[-q]$ vers $\Omega^a_{X/S}$ dans $D(\mathcal{O}_X)$
pour $q \geq 1$ (détails omis). D'après
le lemme \ref{lemma-comparison-general},
$R^qb_*\Omega_{X'/s} \cong i_*R^qp_*\Omega^a_{E/S}$
est l'image directe d'un module sur $Z = \coprod Z_m$.
Observons en outre que la restriction de $R^qp_*\Omega^a_{E/S}$
à $Z_m$ n'est non nulle que pour $q < N - m$. En effet, les fibres de
$E_m \to Z_m$ sont de dimension $N - m - 1$, et l'on peut appliquer Limites, lemme
\ref{limits-lemma-higher-direct-images-zero-above-dimension-fibre}.
L'annulation voulue résulte donc du lemme \ref{lemma-ext-zero}.
\end{proof}














\section{Comparaison de faisceaux de formes différentielles}
\label{section-quasi-finite-syntomic}

\noindent
Le but de cette section est de construire, pour tout morphisme syntomique
localement quasi-fini de schémas $f : Y \to X$, des morphismes
$$
c^p_{Y/X} :
\Omega^p_{Y/\mathbf{Z}}
\longrightarrow
f^*\Omega^p_{X/\mathbf{Z}} \otimes_{\mathcal{O}_Y} \det(\NL_{Y/X})
$$
pour tout $p \geq 0$, satisfaisant aux propriétés suivantes :
\begin{enumerate}
\item
\label{item-degree-zero}
$c_{Y/X}^0(1) = \delta(\NL_{Y/X})$ ; voir
Discriminants, section \ref{discriminant-section-tate-map},
\item
\label{item-multiplicative}
$c_{Y/X}^{q + p}(\omega \wedge \eta) = \omega \wedge c_{Y/X}^p(\eta)$
pour les sections locales $\omega$ de $f^*\Omega^q_{X/\mathbf{Z}}$
et $\eta$ de $\Omega^p_{Y/\mathbf{Z}}$,
\item
\label{item-base-change}
la formation de $c^p_{Y/X}$ est compatible avec la restriction aux ouverts
et le changement de base
(au sens du lemme \ref{lemma-base-change-Garel-upstairs}).
\end{enumerate}
Notons que ces conditions impliquent que les morphismes $c^p_{Y/X}$ sont
$\mathcal{O}_Y$-linéaires et que leur composée avec
$f^*\Omega^p_{X/\mathbf{Z}} \to \Omega^p_{Y/\mathbf{Z}}$
est la multiplication par $\delta(\NL_{Y/X})$.
Nous montrerons en fait qu'il existe une unique famille d'opérateurs
$c^p_{Y/X}$ possédant les propriétés (1), (2) et (3). Ce résultat sera appliqué à la
section \ref{section-trace} pour construire le morphisme trace sur
les complexes de de Rham lorsque $f$ est fini.

\begin{lemma}
\label{lemma-funny-map}
Soit $R$ un anneau, et considérons un diagramme commutatif
$$
\xymatrix{
0 \ar[r] &
K^0 \ar[r] &
L^0 \ar[r] &
M^0 \ar[r] & 0 \\
& & L^{-1} \ar[u]_\partial \ar@{=}[r] &
M^{-1} \ar[u]
}
$$
de $R$-modules dont la ligne supérieure est exacte, et où $M^0$ et $M^{-1}$
sont libres de même rang fini. Pour $p \geq 0$, il existe des morphismes canoniques
$$
c^p : 
\wedge^p(H^0(L^\bullet))
\longrightarrow
\wedge^p(K^0) \otimes_R \det(M^\bullet)
$$
possédant les propriétés suivantes :
\begin{enumerate}
\item $c^0(1) = \delta(M^\bullet)$, et
\item $c^{q + p}(\omega \wedge \eta) = \omega \wedge c^p(\eta)$
pour $\omega \in \wedge^q(K^0)$ et $\eta \in \wedge^p(H^0(L^\bullet))$.
\end{enumerate}
En particulier, la composée de $c^p$ avec le morphisme
$\wedge^p(K^0) \to \wedge^p(H^0(L^\bullet))$ est la multiplication
par $\delta(M^\bullet)$.
\end{lemma}

\begin{proof}
Supposons $M^0$ et $M^{-1}$ libres de rang $n$. Pour tout $p \geq 0$,
il existe une unique surjection
$$
\pi_p :
\wedge^{p + n}(L^0)
\longrightarrow
\wedge^p(K^0) \otimes \wedge^n(M^0)
$$
possédant les deux propriétés suivantes : (i) on a
$
\pi_p(k_1 \wedge \ldots \wedge k_p \wedge l_1 \wedge \ldots \wedge l_n) =
k_1 \wedge \ldots \wedge k_p \otimes
m_1 \wedge \ldots \wedge m_n
$
si $k_1, \ldots, k_p \in K^0$ et si $l_1, \ldots, l_n \in L^0$ ont pour images
$m_1, \ldots, m_n \in M^0$, et (ii) le noyau de $\pi_p$ est le
sous-module engendré par les produits extérieurs
$l_1 \wedge \ldots \wedge l_{p + n}$ tels qu'un nombre $> p$ des
$l_j$ appartiennent à $K^0$. Choisissons une base $e_1, \ldots, e_n$ de $L^{-1} = M^{-1}$.
On définit alors notre morphisme par la règle
$$
\eta \mapsto
\pi_p(\tilde \eta \wedge \partial(e_1) \wedge \ldots \wedge \partial(e_n))
\otimes (e_1 \wedge \ldots \wedge e_n)^{\otimes -1}
$$
où $\eta \in \wedge^p(H^0(L^\bullet))$ et où $\tilde \eta \in \wedge^p(L^0)$
est un relèvement de $\eta$. Le membre de droite a un sens, car
$\det(M^\bullet) = \wedge^n(M^0) \otimes \wedge^n(M^{-1})^{\otimes -1}$.
Il est indépendant du choix de $\tilde \eta$ puisque
$\partial(\xi) \wedge \partial(e_1) \wedge \ldots \wedge \partial(e_n) = 0$
pour tout $\xi \in L^{-1}$. Un calcul direct montre que le morphisme
est indépendant du choix de la base de $L^{-1} = M^{-1}$.
Il résulte immédiatement de la définition de $\delta(M^\bullet)$ dans
Compléments d'algèbre, section \ref{more-algebra-section-determinants-complexes},
que $c^0(1) = \delta(M^\bullet)$. Enfin, la règle de l'assertion (2)
se vérifie directement à partir de la définition.
\end{proof}

\begin{lemma}
\label{lemma-funny-map-functorial}
Soit $R_1 \to R_2$ un homomorphisme d'anneaux. Pour $i = 1, 2$, considérons
des diagrammes commutatifs
$$
\xymatrix{
0 \ar[r] &
K_i^0 \ar[r] &
L_i^0 \ar[r] &
M_i^0 \ar[r] & 0 \\
& & L_i^{-1} \ar[u]_{\partial_i} \ar@{=}[r] &
M_i^{-1} \ar[u]
}
$$
de $R_i$-modules comme dans le lemme \ref{lemma-funny-map}. Supposons donnés
des morphismes $K_1^0 \to K_2^0$, $L_1^j \to L_2^j$ et $M_1^j \to M_2^j$
compatibles avec le morphisme d'anneaux $R_1 \to R_2$ et
avec les morphismes des diagrammes affichés. Si les morphismes
$M_1^j \otimes_{R_1} R_2 \to M_2^j$ sont des isomorphismes, alors
les diagrammes
$$
\xymatrix{ 
\wedge^p(H^0(L_1^\bullet)) \ar[d]
\ar[r]_-{c^p} &
\wedge^p(K_1^0) \otimes_{R_1} \det(M_1^\bullet) \ar[d] \\
\wedge^p(H^0(L_2^\bullet))
\ar[r]^-{c^p} &
\wedge^p(K_2^0) \otimes_{R_2} \det(M_2^\bullet)
}
$$
sont commutatifs.
\end{lemma}

\begin{proof}
Cela résulte de la description explicite du morphisme donnée
dans la démonstration du lemme \ref{lemma-funny-map}. Notons que l'hypothèse
selon laquelle les $M_1^j \otimes_{R_1} R_2 \to M_2^j$ sont des isomorphismes
est nécessaire pour voir que la base utilisée pour $M_1^{-1}$ s'envoie sur une base
de $M_2^{-1}$ (et donc que l'on peut utiliser la « même » base
dans les deux constructions pour en vérifier la compatibilité).
\end{proof}

\begin{remark}
\label{remark-local-description}
Soit $A$ un anneau. Soit $P = A[x_1, \ldots, x_n]$. Soient
$f_1, \ldots, f_n \in P$, et posons $B = P/(f_1, \ldots, f_n)$.
Supposons $A \to B$ quasi-fini. Alors
$B$ est une intersection complète globale relative sur $A$ (Algèbre, définition
\ref{algebra-definition-relative-global-complete-intersection}), et
$(f_1, \ldots, f_n)/(f_1, \ldots, f_n)^2$ est libre, engendré par
les classes $\overline{f}_i$, d'après Algèbre, lemme
\ref{algebra-lemma-relative-global-complete-intersection-conormal}.
Considérons le diagramme suivant
$$
\xymatrix{
\Omega_{A/\mathbf{Z}} \otimes_A B \ar[r] &
\Omega_{P/\mathbf{Z}} \otimes_P B \ar[r] &
\Omega_{P/A} \otimes_P B \\
&
(f_1, \ldots, f_n)/(f_1, \ldots, f_n)^2 \ar[u] \ar@{=}[r] &
(f_1, \ldots, f_n)/(f_1, \ldots, f_n)^2 \ar[u]
}
$$
La colonne de droite représente $\NL_{B/A}$ dans $D(B)$ ; sa cohomologie est donc
$\Omega_{B/A}$ en degré $0$. La ligne supérieure est la suite exacte courte scindée
$0 \to \Omega_{A/\mathbf{Z}} \otimes_A B \to
\Omega_{P/\mathbf{Z}} \otimes_P B \to \Omega_{P/A} \otimes_P B \to 0$.
La colonne du milieu a pour cohomologie $\Omega_{B/\mathbf{Z}}$ en degré $0$,
d'après Algèbre, lemme \ref{algebra-lemma-differential-seq}.
Ainsi, le lemme \ref{lemma-funny-map} fournit des morphismes canoniques de $B$-modules
$$
\Omega^p_{B/\mathbf{Z}} \longrightarrow
\Omega^p_{A/\mathbf{Z}} \otimes_A \det(\NL_{B/A})
$$
pour $p \geq 0$, satisfaisant aux propriétés (1) et (2) et, en particulier,
tels que la composée avec
$\Omega^p_{A/\mathbf{Z}} \to \Omega^p_{B/\mathbf{Z}}$
soit la multiplication par $\delta(\NL_{B/A})$.
\end{remark}

\begin{lemma}
\label{lemma-base-change-Garel-upstairs}
Considérons un diagramme commutatif
$$
\xymatrix{
Y' \ar[d]_{f'} \ar[r]_b & Y \ar[d]^f \\
X' \ar[r]^a & X
}
$$
de schémas qui induit un isomorphisme de $Y'$ sur un sous-schéma ouvert de
$X' \times_X Y$. Supposons $f$ localement quasi-fini et syntomique, et
supposons donnés des morphismes $c^p_{Y/X} : \Omega^p_{Y/\mathbf{Z}} \to
f^*\Omega^p_{X/\mathbf{Z}} \otimes_{\mathcal{O}_Y} \det(\NL_{Y/X})$
pour $p \geq 0$ satisfaisant à
(\ref{item-degree-zero}) et (\ref{item-multiplicative}).
Il existe alors au plus une famille de morphismes
$c^p_{Y'/X'} : \Omega^p_{Y'/\mathbf{Z}} \to
(f')^*\Omega^p_{X'/\mathbf{Z}} \otimes_{\mathcal{O}_{Y'}} \det(\NL_{Y'/X'})$
pour $p \geq 0$
satisfaisant à (\ref{item-degree-zero}) et (\ref{item-multiplicative}),
telle que les diagrammes
$$
\xymatrix{
b^*\Omega^p_{Y/\mathbf{Z}} \ar[rr]_-{b^*c^p_{Y/X}} \ar[d] & &
b^*(f^*\Omega^p_{X/\mathbf{Z}} \otimes_{\mathcal{O}_Y} \det(\NL_{Y/X}))
\ar@{=}[r] &
(f')^*a^*\Omega^p_{X/\mathbf{Z}} \otimes_{\mathcal{O}_Y} b^*\det(\NL_{Y/X}))
\ar[d] \\
\Omega^p_{Y'/\mathbf{Z}} \ar[rrr]^-{c^p_{Y'/X'}} & & &
(f')^*\Omega^p_{X'/\mathbf{Z}} \otimes_{\mathcal{O}_{Y'}} \det(\NL_{Y'/X'})
}
$$
soient commutatifs pour tout $p \geq 0$. Les flèches verticales utilisent ici les morphismes
$b^*\Omega^p_{Y/\mathbf{Z}} \to \Omega^p_{Y'/\mathbf{Z}}$ et
$a^*\Omega^p_{X/\mathbf{Z}} \to \Omega^p_{X'/\mathbf{Z}}$
de la section \ref{section-de-rham-complex},
ainsi que l'identification $b^*\det(\NL_{Y/X}) = \det(\NL_{Y'/X'})$ de
Discriminants, section \ref{discriminant-section-tate-map}.
\end{lemma}

\begin{proof}
Le morphisme
$$
(f')^*\Omega_{X'/\mathbf{Z}} \oplus b^*\Omega_{Y/\mathbf{Z}}
\longrightarrow
\Omega_{Y'/\mathbf{Z}}
$$
est surjectif, car $Y'$ est un sous-schéma ouvert du produit fibré ;
l'énoncé algébrique correspondant affirme que $\Omega_{B \otimes_A C/\mathbf{Z}}$
est engendré, comme module, par les images de $\Omega_{B/\mathbf{Z}}$ et de
$\Omega_{C/\mathbf{Z}}$. Ainsi, pour tout $p$, le morphisme
$$
\bigoplus\nolimits_{i = 0, \ldots, p}
(f')^*\Omega^i_{X'/\mathbf{Z}} \otimes
b^*\Omega^{p - i}_{Y/\mathbf{Z}}
\longrightarrow
\Omega^p_{Y'/\mathbf{Z}}
$$
est surjectif. Les conditions (\ref{item-degree-zero}) et
(\ref{item-multiplicative}), combinées à la commutativité
du diagramme du lemme, impliquent que le morphisme $c^p_{Y'/X'}$
est déterminé (mais son existence n'est pas immédiate).
\end{proof}

\begin{lemma}
\label{lemma-existence-base-change-Garel-upstairs}
Considérons un diagramme commutatif
$$
\xymatrix{
Y' \ar[d]_{f'} \ar[r]_b & Y \ar[d]^f \\
X' \ar[r]^a & X
}
$$
de schémas qui induit un isomorphisme de $Y'$ sur un sous-schéma ouvert de
$X' \times_X Y$. Supposons $f$ localement quasi-fini et syntomique. Alors, pour
tout $y' \in Y'$, on peut trouver des ouverts $V' \subset Y'$,
$V \subset Y$, $U' \subset X$ et $U \subset X$ tels que $y' \in V'$,
$f'(V') \subset U'$, $b(V') \subset V$, $a(U') \subset U$,
$f(V) \subset U$, et tels que, pour $p \geq 0$, il existe des morphismes
$c^p_{V/U}$ et $c^p_{V'/U'}$ satisfaisant à (\ref{item-degree-zero})
et (\ref{item-multiplicative}), compatibles avec le diagramme
$$
\xymatrix{
V' \ar[d] \ar[r]_b & V \ar[d] \\
U' \ar[r] & U
}
$$
au sens précisé dans le lemme \ref{lemma-base-change-Garel-upstairs}.
\end{lemma}

\begin{proof}
Il est clair que l'on peut remplacer $Y'$ par $X' \times_X Y$ pour le démontrer.
Choisissons des ouverts affines $V \subset Y$ et $U \subset X$, où $V$ contient
l'image de $y'$, comme dans
Discriminants, lemme \ref{discriminant-lemma-syntomic-quasi-finite}, partie (5).
Choisissons un ouvert affine $U' \subset X'$ contenant l'image de $y'$
et s'envoyant dans $U$. En remplaçant $X, Y, X', Y'$ par
$U, V, U', U' \times_U V$, on peut supposer être dans la situation
décrite au paragraphe suivant.

\medskip\noindent
Supposons que $X' = \Spec(A')$, $X = \Spec(A)$, $Y = \Spec(B)$ et
$Y' = \Spec(A' \otimes_A B)$, où $B = A[x_1, \ldots, x_n]/(f_1, \ldots, f_n)$
est une intersection complète globale relative sur $A$. Alors
$B' = A'[x_1, \ldots, x_n]/(f'_1, \ldots, f'_n)$ est aussi une intersection
complète globale relative sur $A'$, où $f'_i$ est l'image de $f_i$
dans $A'[x_1, \ldots, x_n]$ ; voir Algèbre, lemme
\ref{algebra-lemma-base-change-relative-global-complete-intersection}.
La construction de la remarque \ref{remark-local-description}
fournit les morphismes $c^p_{V/U}$ et $c^p_{V'/U'}$.
Ils sont compatibles d'après le lemme \ref{lemma-funny-map-functorial}
et la compatibilité des présentations de $B$ et de $B'$ sur
$A$ et $A'$. Nous omettons certains détails.
\end{proof}

\begin{lemma}
\label{lemma-Garel-upstairs}
Il existe une unique règle qui, à tout morphisme de schémas syntomique
localement quasi-fini $f : Y \to X$, associe des morphismes de $\mathcal{O}_Y$-modules
$$
c^p_{Y/X} :
\Omega^p_{Y/\mathbf{Z}}
\longrightarrow
f^*\Omega^p_{X/\mathbf{Z}} \otimes_{\mathcal{O}_Y} \det(\NL_{Y/X})
$$
pour $p \geq 0$, satisfaisant à (\ref{item-degree-zero}), (\ref{item-multiplicative})
et (\ref{item-base-change}). En particulier, la composée de
$c^p_{Y/X}$ avec $f^*\Omega^p_{X/\mathbf{Z}} \to \Omega^p_{Y/\mathbf{Z}}$
est la multiplication par $\delta(\NL_{Y/X})$.
\end{lemma}

\begin{proof}
Cette démonstration est très proche de celle de
Discriminants, proposition \ref{discriminant-proposition-tate-map},
que nous suggérons au lecteur de consulter d'abord.

\medskip\noindent
Reformulons l'énoncé. Considérons la catégorie
$\mathcal{C}$ dont les objets, notés $Y/X$, sont les morphismes syntomiques localement quasi-finis
$f : Y \to X$ de schémas, et dont les morphismes
$b/a : Y'/X' \to Y/X$ sont les diagrammes commutatifs
$$
\xymatrix{
Y' \ar[d]_{f'} \ar[r]_b & Y \ar[d]^f \\
X' \ar[r]^a & X
}
$$
qui induisent un isomorphisme de $Y'$ sur un sous-schéma ouvert de
$X' \times_X Y$. Le lemme signifie que, pour tout objet
$Y/X$ de $\mathcal{C}$, on dispose de morphismes $c^p_{Y/X}$, $p \geq 0$,
possédant les propriétés (\ref{item-degree-zero}) et (\ref{item-multiplicative}),
et que, pour tout morphisme $b/a : Y'/X' \to Y/X$ de $\mathcal{C}$,
$c^p_{Y/X}$ et $c^p_{Y'/X'}$ sont compatibles comme dans le
lemme \ref{lemma-base-change-Garel-upstairs}.

\medskip\noindent
Étant donnés $Y/X$ dans $\mathcal{C}$ et $y \in Y$, on peut trouver
des ouverts affines $V \subset Y$ et $U \subset X$ avec $f(V) \subset U$
tels qu'il existe des morphismes
$$
\Omega^p_{Y/\mathbf{Z}}|_V
\longrightarrow
\left(
f^*\Omega^p_{X/\mathbf{Z}} \otimes_{\mathcal{O}_Y} \det(\NL_{Y/X})
\right)|_V
$$
possédant les propriétés (\ref{item-degree-zero}) et (\ref{item-multiplicative}).
Il suffit pour cela de choisir des ouverts affines comme dans
Discriminants, lemme \ref{discriminant-lemma-syntomic-quasi-finite}, partie (5),
puis d'utiliser la construction de la remarque \ref{remark-local-description}.

\medskip\noindent
Notons que le lieu étale de $f$ est exactement l'ensemble des points où
$\delta(\NL_{Y/X})$ ne s'annule pas ; voir la discussion dans
Discriminants, section \ref{discriminant-section-tate-map}.
En particulier,
$f^*\Omega^p_{X/\mathbf{Z}} \to \Omega^p_{Y/\mathbf{Z}}$
devient un isomorphisme après inversion de $\delta(\NL_{Y/X})$.
On en déduit que, si $\Omega^p_{X/\mathbf{Z}}$ est localement libre de rang fini et si
l'annulateur de $\delta(\NL_{Y/X})$ dans $\mathcal{O}_Y$ est nul,
les morphismes locaux construits au paragraphe précédent
sont uniques et se recollent automatiquement !

\medskip\noindent
Notons $\mathcal{C}_{nice} \subset \mathcal{C}$ la sous-catégorie pleine
formée des $Y/X$ tels que
\begin{enumerate}
\item $\Omega_{X/\mathbf{Z}}$ est localement libre, et
\item l'annulateur de $\delta(\NL_{Y/X})$ dans $\mathcal{O}_Y$ est nul.
\end{enumerate}
D'après les observations du paragraphe précédent, pour tout
objet $Y/X$ de $\mathcal{C}_{nice}$, il existe des morphismes uniques
$c^p_{Y/X}$, $p \geq 0$, satisfaisant aux conditions (\ref{item-degree-zero})
et (\ref{item-multiplicative}). Si $b/a : Y'/X' \to Y/X$
est un morphisme de $\mathcal{C}_{nice}$, les morphismes
$c^p_{Y/X}$ et $c^p_{Y'/X'}$ sont compatibles comme dans le
lemme \ref{lemma-base-change-Garel-upstairs} : en effet,
des morphismes compatibles existent localement d'après le
lemme \ref{lemma-existence-base-change-Garel-upstairs},
et l'unicité que nous venons de mentionner montre qu'ils coïncident
avec $c^p_{Y/X}$ et $c^p_{Y'/X'}$.
Autrement dit, nous avons résolu le problème
sur la sous-catégorie pleine $\mathcal{C}_{nice}$. Pour $Y/X$ dans $\mathcal{C}_{nice}$,
nous continuerons à noter $c^p_{Y/X}$ la solution ainsi obtenue.

\medskip\noindent
Plus généralement, supposons donné un morphisme $b/a : Y'/X' \to Y/X$
de $\mathcal{C}$ tel que $Y/X$ soit un objet de $\mathcal{C}_{nice}$.
Le même argument, utilisant cette fois l'unicité du
lemme \ref{lemma-base-change-Garel-upstairs}, montre qu'il existe des
morphismes $c^p_{Y'/X'}$, $p \geq 0$, satisfaisant aux conditions (\ref{item-degree-zero})
et (\ref{item-multiplicative}), compatibles avec les morphismes déjà construits
$c^p_{Y/X}$. Appelons-les les changements de base
$(b/a)^*c^p_{Y/X}$.

\medskip\noindent
Considérons des morphismes
$$
Y_1/X_1 \xleftarrow{b_1/a_1} Y/X \xrightarrow{b_2/a_2} Y_2/X_2
$$
dans $\mathcal{C}$ tels que $Y_1/X_1$ et $Y_2/X_2$ soient des objets
de $\mathcal{C}_{nice}$.
{\bf Assertion.} Les deux changements de base $(b_i/a_i)^*c^p_{Y_i/X_i}$, $i = 1, 2$,
sont égaux. Nous montrerons d'abord que cette assertion entraîne le lemme,
puis nous la démontrerons.

\medskip\noindent
Soient $d, n \geq 1$, et considérons le morphisme syntomique
localement quasi-fini $Y_{n, d} \to X_{n, d}$
construit dans Discriminants, exemple
\ref{discriminant-example-universal-quasi-finite-syntomic}.
Alors $Y_{n, d}$ et $Y_{n, d}$ sont des schémas irréductibles, de type fini et
lisses sur $\mathbf{Z}$. En effet, $X_{n, d}$ est le spectre d'un
anneau de polynômes sur $\mathbf{Z}$, et $Y_{n, d}$ est un sous-schéma ouvert
d'un tel spectre. Le morphisme $Y_{n, d} \to X_{n, d}$ est syntomique localement quasi-fini
et étale au-dessus d'un ouvert dense ; voir Discriminants, lemme
\ref{discriminant-lemma-universal-quasi-finite-syntomic-etale}.
Ainsi, $\delta(\NL_{Y_{n, d}/X_{n, d}})$ est non nul : on dispose, par exemple,
de la description locale de $\delta(\NL_{Y/X})$ dans
Discriminants, remarque \ref{discriminant-remark-local-description-delta},
et de la description locale des morphismes étales dans
Morphismes, lemme \ref{morphisms-lemma-etale-at-point}, partie (8).
Or une section non nulle d'un module inversible sur un schéma
régulier irréductible a un annulateur nul. Ainsi,
$Y_{n, d}/X_{n, d}$ est un objet de $\mathcal{C}_{nice}$.

\medskip\noindent
Soit $Y/X$ un objet quelconque de $\mathcal{C}$. Soit $y \in Y$.
D'après Discriminants, lemme \ref{discriminant-lemma-locally-comes-from-universal},
on peut trouver $n, d \geq 1$ et des morphismes
$$
Y/X \leftarrow V/U \xrightarrow{b/a} Y_{n, d}/X_{n, d}
$$
de $\mathcal{C}$ tels que $V \subset Y$ et $U \subset X$ soient ouverts.
On dispose alors du changement de base $c^p_{V/U} = (b/a)^*c^p_{Y_{n, d}/X_{n, d}}$.
L'assertion garantit que ces morphismes construits localement $c^p_{V/U}$ se recollent !
On obtient donc des morphismes globaux bien définis $c^p_{Y/X}$ possédant les propriétés
(\ref{item-degree-zero}) et (\ref{item-multiplicative}).
Enfin, soit $b/a : Y'/X' \to Y/X$ un morphisme quelconque de $\mathcal{C}$.
On dispose des morphismes $c^p_{Y'/X'}$, $p \geq 0$ et $c^p_{Y/X}$, $p \geq 0$,
construits dans ce paragraphe. Pour vérifier leur compatibilité au sens du
lemme \ref{lemma-base-change-Garel-upstairs}, on peut travailler localement
sur $Y'/X'$ et $Y/X$. On peut donc supposer qu'il existe un morphisme
$Y/X \to Y_{n, d}/X_{n, d}$. Par construction, la famille
$c^p_{Y'/X'}$ et la famille $c^p_{X/Y}$ sont toutes deux compatibles avec
le morphisme vers $Y_{n, d}/X_{n, d}$. Il en résulte aisément
(en utilisant l'unicité du lemme \ref{lemma-base-change-Garel-upstairs})
que $c^p_{Y'/X'}$ est compatible avec $c^p_{Y/X}$.
Il reste donc à démontrer l'assertion.

\medskip\noindent
Le reste de la démonstration est consacré à l'assertion. On peut choisir un point
$y \in Y$ et montrer que les morphismes coïncident dans un voisinage ouvert de $y$.
On peut donc remplacer $Y_1$ et $Y_2$ par des voisinages ouverts de
l'image de $y$ dans $Y_1$ et $Y_2$. On peut ainsi supposer
$Y, X, Y_1, X_1, Y_2, X_2$ affines, et les écrire comme spectres
d'anneaux $B, A, B_1, A_1, B_2, A_2$. Voici le diagramme
$$
\xymatrix{
B_1 \ar[r] & B & B_2 \ar[l] \\
A_1 \ar[u] \ar[r] & A \ar[u] & A_2 \ar[l] \ar[u]
}
$$
Par hypothèse, le spectre de $B$ est un ouvert affine
du spectre de $A \otimes_{A_1} B_1$ comme de celui de $A \otimes_{A_2} B_2$.
En restreignant encore les ouverts, on peut supposer qu'il existe des éléments
$g_i \in A \otimes_{A_i} B_i$ tels que nos morphismes donnent des isomorphismes
$(A \otimes_{A_i} B_i)_{g_i} = B$ ; voir
Propriétés, lemme \ref{properties-lemma-standard-open-two-affines}.
Soient $x_\alpha$ et $y_\beta$ des familles de
variables suffisamment grandes pour que l'on puisse choisir des surjections
$$
A'_1 = A_1[x_\alpha] \to A
\quad\text{et}\quad
A'_2 = A_2[y_\beta] \to A
$$
de $A_1$-algèbres et de $A_2$-algèbres, respectivement. On peut alors choisir des relèvements
$h_i \in A'_i \otimes_{A_i} B_i$ de $g_i \in A \otimes_{A_i} B_i$,
et considérer le diagramme
$$
\xymatrix{
(A'_1 \otimes_{A_1} B_1)_{h_1} \ar[r] & B &
(A'_2 \otimes_{A_1} B_2)_{h_2} \ar[l] \\
A'_1 \ar[u] \ar[r] & A \ar[u] & A'_2 \ar[l] \ar[u]
}
$$
Par construction, les deux carrés sont cocartésiens. Considérons ensuite le
morphisme d'anneaux
$$
A' = A'_1 \times_A A'_2 \longrightarrow
B' =
(A'_1 \otimes_{A_1} B_1)_{h_1} \times_B
(A'_2 \otimes_{A_1} B_2)_{h_2}
$$
D'après
Compléments d'algèbre, lemme \ref{more-algebra-lemma-module-over-fibre-product},
on a $A'_1 \otimes_{A'} B' = (A'_1 \otimes_{A_1} B_1)_{h_1}$
et $A'_2 \otimes_{A'} B' = (A'_2 \otimes_{A_2} B_2)_{h_2}$. En particulier,
les fibres du morphisme $Y' = \Spec(B') \to \Spec(A') = X'$ sont
des sous-schémas ouverts de changements de base des fibres des morphismes
$Y_i \to X_i$ ; elles sont donc finies et sont des intersections complètes
locales (voir
Discriminants, lemme \ref{discriminant-lemma-syntomic-quasi-finite}).
D'après Compléments d'algèbre, lemme
\ref{more-algebra-lemma-properties-algebras-over-fibre-product},
le morphisme d'anneaux $A' \to B'$ est plat et de présentation finie.
Ainsi, Discriminants, lemme
\ref{discriminant-lemma-syntomic-quasi-finite}, partie (3),
montre que $Y'/X'$ est un objet de $\mathcal{C}$.
Considérons maintenant le diagramme commutatif
$$
\xymatrix{
& & Y/X \ar[ld] \ar@/_2pc/[lldd]_{b_1/a_1} \ar[rd] \ar@/^2pc/[rrdd]^{b_2/a_2} \\
& Y'_1/X'_1 \ar[rd] \ar[ld] & & Y'_2/X'_2 \ar[ld] \ar[rd] \\
Y_1/X_1 & & Y'/X' & & Y_2/X_2
}
$$
où $Y'_i/X'_i$ correspond à $A'_i \to (A'_i \otimes_{A_i} B_i)_{h_i}$.
Notons que $Y'_i/X'_i$ est un objet de $\mathcal{C}_{nice}$, car il
s'obtient en prenant un sous-schéma ouvert du produit d'un
espace affine de dimension infinie avec $Y_i/X_i$ ; nous omettons ce détail.
En particulier, l'image inverse de $c^p_{Y_i/X_i}$ par $(b_i/a_i)$ coïncide
avec celle de $c^p_{Y'_i/X'_i}$ par $Y/X \to Y'_i/X'_i$.
On pourrait conclure si $Y'/X'$ était un objet de $\mathcal{C}_{nice}$,
mais ce n'est presque jamais le cas. En effet, en prenant les images inverses de $c^p_{Y'/X'}$
le long des deux chemins du carré, on obtiendrait la conclusion voulue.
Pour contourner le fait que $Y'/X'$ n'appartient pas à $\mathcal{C}_{nice}$,
notons que les arguments précédents montrent que, quitte à restreindre tous
les schémas $X, Y, X'_1, Y'_1, X'_2, Y'_2, X', Y'$ à des ouverts, on peut trouver
$n, d \geq 1$ et prolonger le diagramme comme suit :
$$
\xymatrix{
& Y/X \ar[ld] \ar[rd] \\
Y'_1/X'_1 \ar[rd] & & Y'_2/X'_2 \ar[ld] \\
& Y'/X' \ar[d] \\
& Y_{n, d}/X_{n, d}
}
$$
On peut alors reprendre l'argument déjà donné en prenant
l'image inverse de $c^p_{Y_{n, d}/X_{n, d}}$. Cela achève la démonstration.
\end{proof}








\section{Morphismes trace sur les complexes de de Rham}
\label{section-trace}

\noindent
Une référence pour une partie des résultats de cette section est \cite{Garel}.
Soit $S$ un schéma. Soit $f : Y \to X$ un morphisme fini localement libre
de schémas sur $S$. Il existe alors un morphisme trace
$\text{Trace}_f : f_*\mathcal{O}_Y \to \mathcal{O}_X$ ; voir
Discriminants, section \ref{discriminant-section-discriminant}.
Dans cette situation, un morphisme trace sur les complexes de de Rham est un morphisme
de complexes
$$
\Theta_{Y/X} : f_*\Omega^\bullet_{Y/S} \longrightarrow \Omega^\bullet_{X/S}
$$
tel que $\Theta_{Y/X}$ soit égal à $\text{Trace}_f$ en degré $0$
et vérifie
$$
\Theta_{Y/X}(\omega \wedge \eta) = \omega \wedge \Theta_{Y/X}(\eta)
$$
pour les sections locales $\omega$ de $\Omega^\bullet_{X/S}$ et $\eta$
de $f_*\Omega^\bullet_{Y/S}$. Nous ne savons pas si un tel morphisme trace
$\Theta_{Y/X}$ existe pour tout morphisme fini localement libre $Y \to X$ ;
veuillez écrire à
\href{mailto:stacks.project@gmail.com}{stacks.project@gmail.com}
si vous disposez d'un contre-exemple ou d'une démonstration.

\begin{example}
\label{example-no-trace}
Voici un exemple où il n'existe pas de morphisme trace sur les complexes de de Rham.
Considérons la $\mathbf{C}$-algèbre $B = \mathbf{C}[x, y]$ munie de
l'action de $G = \{\pm 1\}$ définie par $x \mapsto -x$ et $y \mapsto -y$.
Les invariants $A = B^G$ forment un anneau intègre normal de type fini sur $\mathbf{C}$,
engendré par $x^2, xy, y^2$. Nous affirmons que, pour l'inclusion $A \subset B$,
il n'existe pas de morphisme trace raisonnable
$\Omega_{B/\mathbf{C}} \to \Omega_{A/\mathbf{C}}$
sur les formes de degré $1$. Considérons en effet l'élément
$\omega = x \text{d} y \in \Omega_{B/\mathbf{C}}$.
Puisque $\omega$ est invariant sous l'action de $G$, si un morphisme trace « raisonnable »
existait, $2\omega$ devrait appartenir à l'image de
$\Omega_{A/\mathbf{C}} \to \Omega_{B/\mathbf{C}}$. Ce n'est
pas le cas : il est impossible d'écrire $2\omega$ comme combinaison linéaire
de $\text{d}(x^2)$, $\text{d}(xy)$ et $\text{d}(y^2)$,
même à coefficients dans $B$.
Cet exemple contredit le théorème principal de
\cite{Zannier}.
\end{example}

\begin{lemma}
\label{lemma-Garel}
Il existe une unique règle qui, à tout morphisme syntomique fini
de schémas $f : Y \to X$, associe des morphismes de $\mathcal{O}_X$-modules
$$
\Theta^p_{Y/X} :
f_*\Omega^p_{Y/\mathbf{Z}}
\longrightarrow
\Omega^p_{X/\mathbf{Z}}
$$
satisfaisant aux propriétés suivantes :
\begin{enumerate}
\item la composée avec
$\Omega^p_{X/\mathbf{Z}} \otimes_{\mathcal{O}_X} f_*\mathcal{O}_Y
\to f_*\Omega^p_{Y/\mathbf{Z}}$ est égale à
$\text{id} \otimes \text{Trace}_f$,
où $\text{Trace}_f : f_*\mathcal{O}_Y \to \mathcal{O}_X$
est le morphisme de
Discriminants, section \ref{discriminant-section-discriminant},
\item la règle est compatible avec le changement de base.
\end{enumerate}
\end{lemma}

\begin{proof}
Supposons d'abord $X$ localement noethérien. D'après
le lemme \ref{lemma-Garel-upstairs}, on dispose d'un morphisme canonique
$$
c^p_{Y/X} : \Omega_{Y/\mathbf{Z}}^p
\longrightarrow
f^*\Omega_{X/\mathbf{Z}}^p \otimes_{\mathcal{O}_Y} \det(\NL_{Y/X})
$$
D'après Discriminants, proposition \ref{discriminant-proposition-tate-map},
on dispose d'un isomorphisme canonique
$$
c_{Y/X} : \det(\NL_{Y/X}) \to \omega_{Y/X}
$$
envoyant $\delta(\NL_{Y/X})$ sur $\tau_{Y/X}$. En combinant ces morphismes, on obtient
$$
c^p_{Y/X} \otimes c_{Y/X} :
\Omega_{Y/\mathbf{Z}}^p
\longrightarrow
f^*\Omega_{X/\mathbf{Z}}^p \otimes_{\mathcal{O}_Y} \omega_{Y/X}
$$
D'après Discriminants, section \ref{discriminant-section-finite-morphisms},
cela équivaut à la donnée d'un morphisme
$$
\Theta_{Y/X}^p :
f_*\Omega_{Y/\mathbf{Z}}^p
\longrightarrow
\Omega_{X/\mathbf{Z}}^p
$$
Rappelons que la relation entre $c^p_{Y/X} \otimes c_{Y/X}$
et $\Theta_{Y/X}^p$ utilise le morphisme d'évaluation
$f_*\omega_{Y/X} \to \mathcal{O}_X$
qui envoie $\tau_{Y/X}$ sur $\text{Trace}_f(1)$ ; voir
Discriminants, section \ref{discriminant-section-finite-morphisms}.
La propriété (1) est donc satisfaite. La propriété (2) l'est pour les changements de base par
$X' \to X$ avec $X'$ localement noethérien, car $c^p_{Y/X}$ et
$c_{Y/X}$ sont tous deux compatibles avec ces changements de base. Pour $f : Y \to X$
syntomique fini et $X$ localement noethérien,
nous continuerons à noter $\Theta^p_{Y/X}$ la solution ainsi obtenue.

\medskip\noindent
Unicité. Supposons donné un morphisme syntomique fini
$f: Y \to X$ tel que $X$ soit lisse sur $\Spec(\mathbf{Z})$
et que $f$ soit étale au-dessus d'un ouvert dense de $X$. Nous affirmons que,
dans ce cas, $\Theta^p_{Y/X}$ est déterminé de manière unique par la propriété (1).
Considérons en effet les morphismes
$$
\Omega^p_{X/\mathbf{Z}} \otimes_{\mathcal{O}_X} f_*\mathcal{O}_Y \to
f_*\Omega^p_{Y/\mathbf{Z}} \to
\Omega^p_{X/\mathbf{Z}}
$$
Le faisceau $\Omega^p_{X/\mathbf{Z}}$ est sans torsion (par l'hypothèse de
lissité) ; il suffit donc de vérifier que la restriction de
$\Theta^p_{Y/X}$ est déterminée de manière unique sur l'ouvert dense au-dessus
duquel $f$ est étale, autrement dit, on peut supposer $f$ étale.
Or, si $f$ est étale, alors
$f^*\Omega_{X/\mathbf{Z}} = \Omega_{Y/\mathbf{Z}}$,
de sorte que la première flèche de l'équation affichée est un isomorphisme.
Puisque la composée est déterminée, cela garantit l'unicité.

\medskip\noindent
Soit $f : Y \to X$ un morphisme syntomique fini de schémas localement noethériens.
Soit $x \in X$. D'après Discriminants, lemme
\ref{discriminant-lemma-locally-comes-from-universal-finite},
on peut trouver $d \geq 1$ et un diagramme commutatif
$$
\xymatrix{
Y \ar[d] &
V \ar[d] \ar[l] \ar[r] &
V_d \ar[d] \\
X &
U \ar[l] \ar[r] &
U_d
}
$$
tel que $x \in U \subset X$ soit ouvert, $V = f^{-1}(U)$
et $V = U \times_{U_d} V_d$. Ainsi, $\Theta^p_{Y/X}|_V$
est l'image inverse du morphisme $\Theta^p_{V_d/U_d}$.
Or, d'après la discussion précédente sur l'unicité et
Discriminants, lemmes
\ref{discriminant-lemma-universal-finite-syntomic-smooth} et
\ref{discriminant-lemma-universal-finite-syntomic-etale},
le morphisme $\Theta^p_{V_d/U_d}$ est déterminé de manière unique
par la condition (1). L'unicité est donc établie.

\medskip\noindent
Nous avons maintenant établi l'existence et l'unicité
pour tout morphisme syntomique fini $Y \to X$ avec $X$ localement noethérien.
On pourrait donner un argument analogue à celui de la démonstration du
lemme \ref{lemma-Garel-upstairs} pour passer à $X$ quelconque.
On peut toutefois utiliser directement
l'approximation noethérienne absolue pour conclure.
En effet, pour construire $\Theta^p_{Y/X}$,
il suffit de le faire Zariski-localement sur $X$ (à condition de
démontrer aussi l'unicité). On peut donc supposer $X$ affine (nous omettons
ce détail). On peut alors écrire $X = \lim_{i \in I} X_i$
comme limite de schémas affines noethériens indexée par un ensemble ordonné filtrant $I$.
D'après Algèbre, lemme \ref{algebra-lemma-colimit-category-fp-algebras},
on peut trouver $0 \in I$ et un morphisme de présentation
finie entre schémas affines $f_0 : Y_0 \to X_0$ dont le changement de base à
$X$ est $Y \to X$. En augmentant $0$, on peut supposer $Y_0 \to X_0$
fini et syntomique ; voir
Algèbre, lemmes \ref{algebra-lemma-colimit-lci} et
\ref{algebra-lemma-colimit-finite}. Pour $i \geq 0$, le
changement de base $f_i : Y_i = Y_0 \times_{X_0} X_i \to X_i$ est aussi syntomique fini.
Alors
$$
\Gamma(X, f_*\Omega^p_{Y/\mathbf{Z}}) =
\Gamma(Y, \Omega^p_{Y/\mathbf{Z}}) =
\colim_{i \geq 0} \Gamma(Y_i, \Omega^p_{Y_i/\mathbf{Z}}) =
\colim_{i \geq 0} \Gamma(X_i, f_{i, *}\Omega^p_{Y_i/\mathbf{Z}})
$$
On peut donc, et l'on doit, définir $\Theta^p_{Y/X}$ comme la limite inductive
des morphismes $\Theta^p_{Y_i/X_i}$. Ce morphisme est compatible avec tout
diagramme cartésien
$$
\xymatrix{
Y' \ar[r] \ar[d] & Y \ar[d] \\
X' \ar[r] & X
}
$$
où $X'$ est affine, car nous le savons dans le cas des schémas affines noethériens
d'après les arguments précédents (nous omettons ce détail ; indication : si l'on
écrit aussi $X' = \lim_{j \in J} X'_j$, alors, pour tout $i \in I$, il existe $j \in J$
et un morphisme $X'_j \to X_i$ compatible avec le morphisme $X' \to X$).
Cela achève la démonstration.
\end{proof}

\begin{proposition}
\label{proposition-Garel}
\begin{reference}
\cite{Garel}
\end{reference}
Soit $f : Y \to X$ un morphisme syntomique fini de schémas.
Les morphismes $\Theta^p_{Y/X}$ du lemme \ref{lemma-Garel} définissent un morphisme de complexes
$$
\Theta_{Y/X} :
f_*\Omega^\bullet_{Y/\mathbf{Z}}
\longrightarrow
\Omega^\bullet_{X/\mathbf{Z}}
$$
possédant les propriétés suivantes :
\begin{enumerate}
\item en degré $0$, on obtient
$\text{Trace}_f : f_*\mathcal{O}_Y \to \mathcal{O}_X$ ; voir
Discriminants, section \ref{discriminant-section-discriminant},
\item on a
$\Theta_{Y/X}(\omega \wedge \eta) = \omega \wedge \Theta_{Y/X}(\eta)$
pour $\omega$ dans $\Omega^\bullet_{X/\mathbf{Z}}$ et $\eta$
dans $f_*\Omega^\bullet_{Y/\mathbf{Z}}$,
\item si $f$ est un morphisme sur un schéma de base $S$,
$\Theta_{Y/X}$ induit un morphisme de complexes
$f_*\Omega^\bullet_{Y/S} \to \Omega^\bullet_{X/S}$.
\end{enumerate}
\end{proposition}

\begin{proof}
D'après Discriminants, lemme
\ref{discriminant-lemma-locally-comes-from-universal-finite},
pour tout $x \in X$, on peut trouver $d \geq 1$ et un diagramme commutatif
$$
\xymatrix{
Y \ar[d] &
V \ar[d] \ar[l] \ar[r] &
V_d \ar[d] \ar[r] &
Y_d = \Spec(B_d) \ar[d] \\
X &
U \ar[l] \ar[r] &
U_d \ar[r] &
X_d = \Spec(A_d)
}
$$
tel que $x \in U \subset X$ soit un ouvert affine, $V = f^{-1}(U)$
et $V = U \times_{U_d} V_d$. Écrivons $U = \Spec(A)$ et $V = \Spec(B)$,
observons que $B = A \otimes_{A_d} B_d$, et rappelons que
$B_d = A_d e_1 \oplus \ldots \oplus A_d e_d$. Supposons donnés
$a_1, \ldots, a_r \in A$ et $b_1, \ldots, b_s \in B$.
On peut écrire $b_j = \sum a_{j, l} e_l$ avec $a_{j, l} \in A$.
Posons $N = r + sd$ et considérons les factorisations
$$
\xymatrix{
V \ar[r] \ar[d] &
V' = \mathbf{A}^N \times V_d \ar[r] \ar[d] &
V_d \ar[d] \\
U \ar[r]&
U' = \mathbf{A}^N \times U_d \ar[r] &
U_d
}
$$
Ici, la flèche horizontale inférieure droite est donnée par le morphisme
$U \to U_d$ (du diagramme précédent) et le morphisme
$U \to \mathbf{A}^N$ défini par $a_1, \ldots, a_r, a_{1, 1}, \ldots, a_{s, d}$.
On voit alors que les fonctions $a_1, \ldots, a_r$ appartiennent à l'image de
$\Gamma(U', \mathcal{O}_{U'}) \to \Gamma(U, \mathcal{O}_U)$,
et que les fonctions $b_1, \ldots, b_s$ appartiennent à l'image de
$\Gamma(V', \mathcal{O}_{V'}) \to \Gamma(V, \mathcal{O}_V)$.
On voit ainsi que, pour toute famille finie d'éléments\footnote{Ces
éléments sont en effet des sommes finies d'éléments de la forme
$a_0 \text{d}a_1 \wedge \ldots \wedge \text{d}a_i$ avec
$a_0, \ldots, a_i \in A$, ou des sommes finies d'éléments de la forme
$b_0 \text{d}b_1 \wedge \ldots \wedge \text{d}b_j$ avec
$b_0, \ldots, b_j \in B$.} des groupes
$$
\Gamma(V, \Omega^i_{Y/\mathbf{Z}}),\quad i = 0, 1, 2, \ldots
\quad\text{et}\quad
\Gamma(U, \Omega^j_{X/\mathbf{Z}}),\quad j = 0, 1, 2, \ldots
$$,
on peut trouver des factorisations $V \to V' \to V_d$ et
$U \to U' \to U_d$ avec $V' = \mathbf{A}^N \times V_d$ et
$U' = \mathbf{A}^N \times U_d$ comme ci-dessus,
telles que ces sections soient les images inverses de sections de
$$
\Gamma(V', \Omega^i_{V'/\mathbf{Z}}),\quad i = 0, 1, 2, \ldots
\quad\text{et}\quad
\Gamma(U', \Omega^j_{U'/\mathbf{Z}}),\quad j = 0, 1, 2, \ldots
$$
Il en résulte que, pour vérifier
$\text{d} \circ \Theta_{Y/X} = \Theta_{Y/X} \circ \text{d}$,
il suffit de vérifier cette égalité pour $\Theta_{V'/U'}$.
Il en va de même de la propriété (2) du lemme.

\medskip\noindent
D'après Discriminants, lemmes
\ref{discriminant-lemma-universal-finite-syntomic-smooth} et
\ref{discriminant-lemma-universal-finite-syntomic-etale},
le schéma $U_d$ est lisse et le morphisme $V_d \to U_d$
est étale au-dessus d'un ouvert dense de $U_d$.
Il en va donc de même du morphisme
$V' \to U'$. Puisque $\Omega_{U'/\mathbf{Z}}$ est localement libre et donc
$\Omega^p_{U'/\mathbf{Z}}$ sans
torsion, il suffit de vérifier les relations voulues
après restriction à l'ouvert au-dessus duquel $V'$ est fini étale.
On peut ensuite vérifier ces relations après un changement de base étale
surjectif. On peut donc décomposer le revêtement fini étale
et supposer que l'on considère un morphisme de la forme
$$
\coprod\nolimits_{i = 1, \ldots, d} W \longrightarrow W
$$
où $W$ est lisse sur $\mathbf{Z}$.
Dans ce cas, toutes les propriétés locales de notre construction se vérifient trivialement
(pourvu qu'elles soient vraies). Cela achève la démonstration de (1) et (2).

\medskip\noindent
Enfin, (3) résulte de (2), car $\Omega_{Y/S}$
est le quotient de $\Omega_{Y/\mathbf{Z}}$ par le sous-module
engendré par les images inverses de sections locales de $\Omega_{S/\mathbf{Z}}$.
\end{proof}

\begin{example}
\label{example-Garel}
Soit $A$ un anneau. Soit $f = x^d + \sum_{0 \leq i < d} a_{d - i} x^i \in A[x]$.
Soit $B = A[x]/(f)$. D'après la proposition \ref{proposition-Garel},
on dispose d'un morphisme de complexes
$$
\Theta_{B/A} : \Omega^\bullet_B \longrightarrow \Omega^\bullet_A
$$
En particulier, si $t \in B$ désigne l'image de $x \in A[x]$,
on peut considérer les éléments
$$
\Theta_{B/A}(t^i\text{d}t) \in \Omega^1_A,\quad i = 0, \ldots, d - 1
$$
Quels sont ces éléments ? Par le même principe que dans la démonstration de la
proposition \ref{proposition-Garel}, il suffit de les calculer
dans le cas universel, c'est-à-dire lorsque $A = \mathbf{Z}[a_1, \ldots, a_d]$,
ou même après remplacement de $A$ par le corps des fractions
$\mathbf{Q}(a_1, \ldots, a_d)$. En écrivant symboliquement
$$
f = \prod\nolimits_{i = 1, \ldots, d} (x - \alpha_i)
$$,
on voit que, sur $\mathbf{Q}(\alpha_1, \ldots, \alpha_d)$,
l'algèbre $B$ se décompose :
$$
\mathbf{Q}(a_0, \ldots, a_{d - 1})[x]/(f) 
\longrightarrow
\prod\nolimits_{i = 1, \ldots, d} \mathbf{Q}(\alpha_1, \ldots, \alpha_d),
\quad
t \longmapsto (\alpha_1, \ldots, \alpha_d)
$$
Ainsi, par exemple,
$$
\Theta(\text{d}t) = \sum \text{d} \alpha_i = - \text{d}a_1
$$
On a ensuite
$$
\Theta(t\text{d}t) = \sum \alpha_i \text{d}\alpha_i =
a_1 \text{d} a_1 - \text{d}a_2
$$
Puis on a
$$
\Theta(t^2\text{d}t) = \sum \alpha_i^2 \text{d}\alpha_i =
- a_1^2 \text{d} a_1 + a_1 \text{d}a_2 + a_2 \text{d}a_1 - \text{d}a_3
$$
(sauf erreur de calcul), et ainsi de suite. Cela suggère que,
si $f(x) = x^d - a$, alors
$$
\Theta_{B/A}(t^i\text{d}t) =
\left\{
\begin{matrix}
0 & \text{if} & i = 0, \ldots, d - 2 \\
\text{d}a & \text{if} & i = d - 1
\end{matrix}
\right.
$$
dans $\Omega_A$. C'est vrai, car, dans ce cas particulier, on peut effectuer
le calcul pour l'extension $\mathbf{Q}(a)[x]/(x^d - a)$
afin de le vérifier directement.
\end{example}

\begin{lemma}
\label{lemma-Garel-map-frobenius-smooth-char-p}
Soit $p$ un nombre premier. Soit $X \to S$ un morphisme lisse
de dimension relative $d$ de schémas de caractéristique $p$.
Le Frobenius relatif $F_{X/S} : X \to X^{(p)}$ de $X/S$
(Variétés, définition \ref{varieties-definition-relative-frobenius})
est syntomique fini, et le morphisme correspondant
$$
\Theta_{X/X^{(p)}} :
F_{X/S, *}\Omega^\bullet_{X/S} \to \Omega^\bullet_{X^{(p)}/S}
$$
est nul en tout degré sauf en degré $d$, où il définit une
surjection.
\end{lemma}

\begin{proof}
Observons que $F_{X/S}$ est un morphisme fini d'après
Variétés, lemme \ref{varieties-lemma-relative-frobenius-finite}.
Pour montrer que $F_{X/S}$ est plat, il suffit de montrer que
le morphisme $F_{X/S, s} : X_s \to X^{(p)}_s$ entre fibres
est plat pour tout $s \in S$ ; voir Compléments sur les morphismes, théorème
\ref{more-morphisms-theorem-criterion-flatness-fibre}.
La platitude de $X_s \to X^{(p)}_s$ résulte de
Algèbre, lemme \ref{algebra-lemma-CM-over-regular-flat}
(et de la finitude déjà établie).
D'après Compléments sur les morphismes, lemme
\ref{more-morphisms-lemma-lci-permanence},
le morphisme $F_{X/S}$ est localement d'intersection complète.
Ainsi, $F_{X/S}$ est syntomique fini (voir
Compléments sur les morphismes, lemme \ref{more-morphisms-lemma-flat-lci}).

\medskip\noindent
Pour tout point $x \in X$, on peut choisir un diagramme commutatif
$$
\xymatrix{
X \ar[d] & U \ar[l] \ar[d]_\pi \\
S & \mathbf{A}^d_S \ar[l]
}
$$
où $\pi$ est étale et où $x \in U$ est un ouvert de $X$ ; voir
Morphismes, lemme \ref{morphisms-lemma-smooth-etale-over-affine-space}.
Observons que
$\mathbf{A}^d_S \to \mathbf{A}^d_S$, $(x_1, \ldots, x_d) \mapsto
(x_1^p, \ldots, x_d^p)$, est le Frobenius relatif de $\mathcal{A}^d_S$
sur $S$. Le diagramme commutatif
$$
\xymatrix{
U \ar[d]_\pi \ar[r]_{F_{X/S}} & U^{(p)} \ar[d]^{\pi^{(p)}} \\
\mathbf{A}^d_S \ar[r]^{x_i \mapsto x_i^p} & \mathbf{A}^d_S
}
$$
de
Variétés, lemme \ref{varieties-lemma-relative-frobenius-endomorphism-identity},
pour $\pi : U \to \mathbf{A}^d_S$ est cartésien d'après
Morphismes étales, lemme
\ref{etale-lemma-relative-frobenius-etale}.
Puisque la construction de $\Theta$ est compatible avec le changement de base
et que $\Omega_{U/S} = \pi^*\Omega_{\mathbf{A}^d_S/S}$
(lemme \ref{lemma-etale}),
il suffit de démontrer le lemme pour
$\mathbf{A}^d_S$.

\medskip\noindent
Soit $A$ un anneau de caractéristique $p$. Considérons l'unique homomorphisme de $A$-algèbres
$A[y_1, \ldots, y_d] \to A[x_1, \ldots, x_d]$
envoyant $y_i$ sur $x_i^p$. Les arguments précédents
nous ramènent au calcul du morphisme
$$
\Theta^i : \Omega^i_{A[x_1, \ldots, x_d]/A} \to
\Omega^i_{A[y_1, \ldots, y_d]/A}
$$
Nous invitons vivement le lecteur à effectuer lui-même ce calcul.
Comme dans l'exemple \ref{example-Garel}, on peut se ramener au calcul
d'une formule pour $\Theta^i$ dans le cas universel
$$
\mathbf{Z}[y_1, \ldots, y_d] \to \mathbf{Z}[x_1, \ldots, x_d],\quad
y_i \mapsto x_i^p
$$
On peut à son tour trouver la formule pour $\Theta^i$ en calculant dans le cas
complexe, c'est-à-dire pour le morphisme de $\mathbf{C}$-algèbres
$$
\mathbf{C}[y_1, \ldots, y_d] \to \mathbf{C}[x_1, \ldots, x_d],\quad
y_i \mapsto x_i^p
$$
On peut même inverser $x_1, \ldots, x_d$ et $y_1, \ldots, y_d$.
Dans ce cas, on a $\text{d}x_i = p^{-1} x_i^{- p + 1}\text{d}y_i$.
On voit donc que
\begin{align*}
\Theta^i(
x_1^{e_1} \ldots x_d^{e_d} \text{d}x_1 \wedge \ldots \wedge \text{d}x_i)
& =
p^{-i} \Theta^i(
x_1^{e_1 - p + 1} \ldots x_i^{e_i - p + 1} x_{i + 1}^{e_{i + 1}} \ldots
x_d^{e_d} \text{d}y_1 \wedge \ldots \wedge \text{d}y_i ) \\
& =
p^{-i} \text{Trace}(x_1^{e_1 - p + 1} \ldots x_i^{e_i - p + 1}
x_{i + 1}^{e_{i + 1}} \ldots x_d^{e_d})
\text{d}y_1 \wedge \ldots \wedge \text{d}y_i
\end{align*}
d'après les propriétés de $\Theta^i$. Un calcul élémentaire montre
que la trace figurant dans cette expression est nulle, sauf si
$e_1, \ldots, e_i$ sont congrus à $-1$ modulo $p$
et si $e_{i + 1}, \ldots, e_d$ sont divisibles par $p$.
De plus, dans ce cas, on obtient
$$
p^{d - i} y_1^{(e_1 - p + 1)/p} \ldots y_i^{(e_i - p + 1)/p}
y_{i + 1}^{e_{i + 1}/p} \ldots y_d^{e_d/p}
\text{d}y_1 \wedge \ldots \wedge \text{d}y_i
$$
On en déduit que le résultat est nul en caractéristique $p$, sauf si $d = i$,
auquel cas on obtient toutes les formes de degré $d$.
\end{proof}








\section{Dualité de Poincaré}
\label{section-poincare-duality}

\noindent
Dans cette section, nous démontrons la dualité de Poincaré pour la cohomologie de de Rham
d'un schéma propre et lisse sur un corps. Expliquons d'abord comment elle
se présente pour la cohomologie de Hodge.

\begin{lemma}
\label{lemma-duality-hodge}
Soit $k$ un corps. Soit $X$ un schéma non vide, lisse et propre sur $k$,
équidimensionnel de dimension $d$. Il existe une application $k$-linéaire
$$
t : H^d(X, \Omega^d_{X/k}) \longrightarrow k
$$,
unique à précomposition près par la multiplication par une unité de
$H^0(X, \mathcal{O}_X)$, possédant la propriété suivante : pour tous $p, q$, l'accouplement
$$
H^q(X, \Omega^p_{X/k}) \times H^{d - q}(X, \Omega^{d - p}_{X/k})
\longrightarrow
k, \quad
(\xi, \xi') \longmapsto t(\xi \cup \xi')
$$
est parfait.
\end{lemma}

\begin{proof}
D'après Dualité pour les schémas, lemme \ref{duality-lemma-duality-proper-over-field},
on a $\omega_X^\bullet = \Omega^d_{X/k}[d]$.
Puisque $\Omega_{X/k}$ est localement libre de rang $d$
(Morphismes, lemme \ref{morphisms-lemma-smooth-omega-finite-locally-free}),
on a
$$
\Omega^d_{X/k} \otimes_{\mathcal{O}_X} (\Omega^p_{X/k})^\vee
\cong
\Omega^{d - p}_{X/k}
$$
On obtient donc une application $k$-linéaire $t : H^d(X, \Omega^d_{X/k}) \to k$
pour laquelle l'énoncé est vrai, d'après Dualité pour les schémas, lemme
\ref{duality-lemma-duality-proper-over-field-perfect}.
En particulier, l'accouplement
$H^0(X, \mathcal{O}_X) \times H^d(X, \Omega^d_{X/k}) \to k$
est parfait, ce qui implique que toute application $k$-linéaire
$t' : H^d(X, \Omega^d_{X/k}) \to k$ est de la forme
$\xi \mapsto t(g\xi)$ pour un certain $g \in H^0(X, \mathcal{O}_X)$.
Bien entendu, pour que $t'$ définisse encore une dualité
entre $H^0(X, \mathcal{O}_X)$ et $H^d(X, \Omega^d_{X/k})$,
il faut que $g$ soit une unité. Notons $\langle -, - \rangle_{p, q}$
l'accouplement construit à l'aide de $t$, et $\langle -, - \rangle'_{p, q}$
celui construit à l'aide de $t'$. On a clairement
$$
\langle \xi, \xi' \rangle'_{p, q} =
\langle g\xi, \xi' \rangle_{p, q}
$$
pour $\xi \in H^q(X, \Omega^p_{X/k})$ et
$\xi' \in H^{d - q}(X, \Omega^{d - p}_{X/k})$. Puisque $g$ est une unité, donc
inversible, on voit qu'en utilisant $t'$ au lieu de $t$, on obtient encore des accouplements
parfaits pour tous $p, q$.
\end{proof}

\begin{lemma}
\label{lemma-bottom-part-degenerates}
Soit $k$ un corps. Soit $X$ un schéma lisse et propre sur $k$. L'application
$$
\text{d} : H^0(X, \mathcal{O}_X) \to H^0(X, \Omega^1_{X/k})
$$
est nulle.
\end{lemma}

\begin{proof}
Puisque $X$ est lisse sur $k$, il est géométriquement réduit sur $k$ ; voir
Variétés, lemme \ref{varieties-lemma-smooth-geometrically-normal}.
Ainsi, $H^0(X, \mathcal{O}_X) = \prod k_i$
est un produit fini d'extensions de corps finies
séparables $k_i/k$ ; voir Variétés, lemme
\ref{varieties-lemma-proper-geometrically-reduced-global-sections}.
Il en résulte que $\Omega_{H^0(X, \mathcal{O}_X)/k} = \prod \Omega_{k_i/k} = 0$
(voir par exemple Algèbre, lemme
\ref{algebra-lemma-characterize-separable-algebraic-field-extensions}).
Puisque l'application du lemme se factorise sous la forme
$$
H^0(X, \mathcal{O}_X) \to
\Omega_{H^0(X, \mathcal{O}_X)/k} \to
H^0(X, \Omega_{X/k})
$$
par fonctorialité du complexe de de Rham
(voir la section \ref{section-de-rham-complex}), on conclut.
\end{proof}

\begin{lemma}
\label{lemma-top-part-degenerates}
Soit $k$ un corps. Soit $X$ un schéma lisse et propre sur $k$,
équidimensionnel de dimension $d$. L'application
$$
\text{d} : H^d(X, \Omega^{d - 1}_{X/k}) \to H^d(X, \Omega^d_{X/k})
$$
est nulle.
\end{lemma}

\begin{proof}
Il est tentant de penser que cela résulte de la combinaison des
lemmes \ref{lemma-bottom-part-degenerates} et \ref{lemma-duality-hodge}.
Cela ne fonctionne toutefois pas, car les morphismes $\mathcal{O}_X \to \Omega^1_{X/k}$
et $\Omega^{d - 1}_{X/k} \to \Omega^d_{X/k}$ ne sont pas $\mathcal{O}_X$-linéaires.
On ne peut donc utiliser la fonctorialité décrite dans
Dualité pour les schémas, remarque
\ref{duality-remark-coherent-duality-proper-over-field},
pour conclure que l'application du lemme \ref{lemma-bottom-part-degenerates}
est duale de celle du présent lemme.

\medskip\noindent
On peut remplacer $X$ par une composante connexe de $X$. On peut donc supposer
$X$ irréductible. D'après
Variétés, lemmes \ref{varieties-lemma-smooth-geometrically-normal} et
\ref{varieties-lemma-proper-geometrically-reduced-global-sections},
$k' = H^0(X, \mathcal{O}_X)$ est une extension finie
séparable $k'/k$. Puisque $\Omega_{k'/k} = 0$
(voir par exemple Algèbre, lemme
\ref{algebra-lemma-characterize-separable-algebraic-field-extensions}),
on voit que $\Omega_{X/k} = \Omega_{X/k'}$
(voir Morphismes, lemme \ref{morphisms-lemma-triangle-differentials}).
On peut donc remplacer $k$ par $k'$ et supposer $H^0(X, \mathcal{O}_X) = k$.

\medskip\noindent
Supposons $H^0(X, \mathcal{O}_X) = k$. On en déduit que
$\dim H^d(X, \Omega^d_{X/k}) = 1$ d'après le lemme \ref{lemma-duality-hodge}.
Supposons d'abord que la caractéristique de $k$ soit un nombre premier $p$.
Notons $F_{X/k} : X \to X^{(p)}$ le Frobenius relatif de $X$ sur $k$ ;
gardons à l'esprit les propriétés de ce morphisme établies dans le
lemme \ref{lemma-Garel-map-frobenius-smooth-char-p}.
Considérons le diagramme commutatif
$$
\xymatrix{
H^d(X, \Omega^{d - 1}_{X/k}) \ar[d] \ar[r] &
H^d(X^{(p)}, F_{X/k, *}\Omega^{d - 1}_{X/k}) \ar[d] \ar[r]_{\Theta^{d - 1}} &
H^d(X^{(p)}, \Omega^{d - 1}_{X^{(p)}/k}) \ar[d] \\
H^d(X, \Omega^d_{X/k}) \ar[r] &
H^d(X^{(p)}, F_{X/k, *}\Omega^d_{X/k}) \ar[r]^{\Theta^d} &
H^d(X^{(p)}, \Omega^d_{X^{(p)}/k})
}
$$
Les deux flèches horizontales de gauche sont des isomorphismes, car $F_{X/k}$ est fini ; voir
Cohomologie des schémas, lemme \ref{coherent-lemma-relative-affine-cohomology}.
Le carré de droite est commutatif, car $\Theta_{X^{(p)}/X}$ est un morphisme de
complexes et $\Theta^{d - 1}$ est nul. Il suffit donc de montrer que
$\Theta^d$ est non nul (car la dimension de la source de
$\Theta^d$ est $1$ d'après la discussion précédente). Or on sait que
$$
\Theta^d : F_{X/k, *}\Omega^d_{X/k} \to \Omega^d_{X^{(p)}/k}
$$
est surjectif, donc encore surjectif après application du foncteur exact à
droite $H^d(X^{(p)}, -)$ (l'exactitude à droite résulte de l'annulation de la cohomologie
au-delà de $d$, conséquence de
Cohomologie, proposition \ref{cohomology-proposition-vanishing-Noetherian}).
Enfin, $H^d(X^{(d)}, \Omega^d_{X^{(d)}/k})$ est non nul, par exemple parce
qu'il est dual de $H^0(X^{(d)}, \mathcal{O}_{X^{(p)}})$ d'après le
lemme \ref{lemma-duality-hodge} appliqué à $X^{(p)}$ sur $k$.
Cela achève la démonstration dans ce cas.

\medskip\noindent
Supposons enfin que la caractéristique de $k$ soit $0$.
On peut écrire $k$ comme limite inductive filtrante de ses
$\mathbf{Z}$-sous-algèbres de type fini $R$. Pour l'une d'elles, on peut trouver un
diagramme cartésien de schémas
$$
\xymatrix{
X \ar[d] \ar[r] & Y \ar[d] \\
\Spec(k) \ar[r] & \Spec(R)
}
$$
tel que $Y \to \Spec(R)$ soit lisse de dimension relative $d$ et propre.
Voir Limites, lemmes \ref{limits-lemma-descend-finite-presentation},
\ref{limits-lemma-descend-smooth}, \ref{limits-lemma-descend-dimension-d} et
\ref{limits-lemma-eventually-proper}.
Les modules $M^{i, j} = H^j(Y, \Omega^i_{Y/R})$ sont des $R$-modules de type fini ; voir
Cohomologie des schémas, lemme
\ref{coherent-lemma-proper-over-affine-cohomology-finite}.
Ainsi, en remplaçant $R$ par un localisé, on peut supposer tous ces
modules libres de rang fini. On a
$M^{i, j} \otimes_R k = H^j(X, \Omega^i_{X/k})$
par changement de base plat (Cohomologie des schémas, lemme
\ref{coherent-lemma-flat-base-change-cohomology}).
Il suffit donc de montrer que $M^{d - 1, d} \to M^{d, d}$
est nul. C'est un morphisme entre modules libres de rang fini sur un anneau intègre ;
il suffit donc de trouver un ensemble dense d'idéaux premiers $\mathfrak p \subset R$
tel que le morphisme devienne nul après tensorisation par $\kappa(\mathfrak p)$.
Puisque $R$ est de type fini sur $\mathbf{Z}$, on peut prendre
l'ensemble des idéaux premiers $\mathfrak p$ dont le corps résiduel
est de caractéristique positive (détails omis). Observons que
$$
M^{d - 1, d} \otimes_R \kappa(\mathfrak p) =
H^d(Y_{\kappa(\mathfrak p)},
\Omega^{d - 1}_{Y_{\kappa(\mathfrak p)}/\kappa(\mathfrak p)})
$$,
par exemple d'après Limites, lemme
\ref{limits-lemma-higher-direct-images-zero-above-dimension-fibre}.
Il en va de même pour $M^{d, d}$. On voit donc que
$M^{d - 1, d} \otimes_R \kappa(\mathfrak p) \to
M^{d, d} \otimes_R \kappa(\mathfrak p)$
est nul d'après le cas de caractéristique positive traité ci-dessus.
\end{proof}

\begin{proposition}
\label{proposition-poincare-duality}
Soit $k$ un corps. Soit $X$ un schéma non vide, lisse et propre sur $k$,
équidimensionnel de dimension $d$. Il existe une application $k$-linéaire
$$
t : H^{2d}_{dR}(X/k) \longrightarrow k
$$,
unique à précomposition près par la multiplication par une unité de
$H^0(X, \mathcal{O}_X)$, possédant la propriété suivante : pour tout $i$, l'accouplement
$$
H^i_{dR}(X/k) \times H_{dR}^{2d - i}(X/k)
\longrightarrow
k, \quad
(\xi, \xi') \longmapsto t(\xi \cup \xi')
$$
est parfait.
\end{proposition}

\begin{proof}
La suite spectrale de Hodge vers de Rham
(section \ref{section-hodge-to-de-rham}), l'annulation
de $\Omega^i_{X/k}$ pour $i > d$, le théorème d'annulation de
Cohomologie, proposition \ref{cohomology-proposition-vanishing-Noetherian},
et les résultats des lemmes \ref{lemma-bottom-part-degenerates} et
\ref{lemma-top-part-degenerates}
montrent que $H^0_{dR}(X/k) = H^0(X, \mathcal{O}_X)$
et $H^d(X, \Omega^d_{X/k}) = H_{dR}^{2d}(X/k)$.
Plus précisément, ces identifications proviennent des morphismes
de complexes
$$
\Omega^\bullet_{X/k} \to \mathcal{O}_X[0]
\quad\text{et}\quad
\Omega^d_{X/k}[-d] \to \Omega^\bullet_{X/k}
$$
Choisissons $t : H_{dR}^{2d}(X/k) \to k$ correspondant, par cette identification,
à une application $t$ comme dans le lemme \ref{lemma-duality-hodge}.
On voit alors, en tout cas, que l'accouplement affiché dans le lemme
est parfait pour $i = 0$.

\medskip\noindent
Notons $\underline{k}$ le faisceau constant de valeur $k$ sur $X$.
Abrégeons $\Omega^\bullet = \Omega^\bullet_{X/k}$.
Considérons le morphisme (\ref{equation-wedge}), qui s'écrit ici
$$
\wedge :
\text{Tot}(\Omega^\bullet \otimes_{\underline{k}} \Omega^\bullet)
\longrightarrow
\Omega^\bullet
$$
Pour tout entier $p = 0, 1, \ldots, d$, ce morphisme
annule le sous-complexe
$\text{Tot}(\sigma_{> p} \Omega^\bullet \otimes_{\underline{k}}
\sigma_{\geq d - p} \Omega^\bullet)$ pour des raisons de degré.
Ainsi, la restriction de $\wedge$ au sous-complexe
$\text{Tot}(\Omega^\bullet \otimes_{\underline{k}}
\sigma_{\geq d - p}\Omega^\bullet)$ se factorise par un morphisme de complexes
$$
\gamma_p :
\text{Tot}(\sigma_{\leq p} \Omega^\bullet \otimes_{\underline{k}}
\sigma_{\geq d - p} \Omega^\bullet)
\longrightarrow
\Omega^\bullet
$$
Par le même procédé qu'à la section \ref{section-cup-product}, on obtient des
cup-produits
$$
H^i(X, \sigma_{\leq p} \Omega^\bullet) \times
H^{2d - i}(X, \sigma_{\geq d - p}\Omega^\bullet)
\longrightarrow
H_{dR}^{2d}(X, \Omega^\bullet)
$$
Nous montrerons par récurrence sur $p$ que ces cup-produits, composés avec $t$,
induisent des accouplements parfaits entre $H^i(X, \sigma_{\leq p} \Omega^\bullet)$
et $H^{2d - i}(X, \sigma_{\geq d - p}\Omega^\bullet)$. Pour $p = d$,
c'est l'assertion de la proposition.

\medskip\noindent
Le cas initial est $p = 0$. On obtient alors simplement l'accouplement
entre $H^i(X, \mathcal{O}_X)$ et $H^{d - i}(X, \Omega^d)$ du
lemme \ref{lemma-duality-hodge}, d'où le résultat.

\medskip\noindent
Étape de récurrence. Supposons le résultat vrai pour $p$.
Considérons alors le triangle distingué
$$
\Omega^{p + 1}[-p - 1] \to
\sigma_{\leq p + 1}\Omega^\bullet \to
\sigma_{\leq p}\Omega^\bullet \to
\Omega^{p + 1}[-p]
$$
et le triangle distingué
$$
\sigma_{\geq d - p}\Omega^\bullet \to
\sigma_{\geq d - p - 1}\Omega^\bullet \to
\Omega^{d - p - 1}[-d + p + 1] \to
(\sigma_{\geq d - p}\Omega^\bullet)[1]
$$
Observons que tous deux sont des triangles distingués dans la catégorie homotopique
des complexes de faisceaux de $\underline{k}$-modules ; en particulier, les
morphismes $\sigma_{\leq p}\Omega^\bullet \to \Omega^{p + 1}[-p]$ et
$\Omega^{d - p - 1}[-d + p + 1] \to (\sigma_{\geq d - p}\Omega^\bullet)[1]$
sont donnés par de véritables morphismes de complexes, à l'aide de la différentielle
$\Omega^p \to \Omega^{p + 1}$ et de la différentielle
$\Omega^{d - p - 1} \to \Omega^{d - p}$.
Considérons les suites exactes longues de cohomologie associées à ces
triangles distingués
$$
\xymatrix{
H^{i - 1}(X, \sigma_{\leq p}\Omega^\bullet) \ar[d]_a \\
H^i(X, \Omega^{p + 1}[-p - 1]) \ar[d]_b \\
H^i(X, \sigma_{\leq p + 1}\Omega^\bullet) \ar[d]_c \\
H^i(X, \sigma_{\leq p}\Omega^\bullet) \ar[d]_d \\
H^{i + 1}(X, \Omega^{p + 1}[-p - 1])
}
\quad\quad
\xymatrix{
H^{2d - i  + 1}(X, \sigma_{\geq d - p}\Omega^\bullet) \\
H^{2d - i}(X, \Omega^{d - p - 1}[-d + p + 1]) \ar[u]_{a'} \\
H^{2d - i}(X, \sigma_{\geq d - p - 1}\Omega^\bullet) \ar[u]_{b'} \\
H^{2d - i}(X, \sigma_{\geq d - p}\Omega^\bullet) \ar[u]_{c'} \\
H^{2d - i - 1}(X, \Omega^{d - p - 1}[-d + p + 1]) \ar[u]_{d'}
}
$$
Par récurrence et d'après le lemme \ref{lemma-duality-hodge},
on sait que les accouplements construits ci-dessus entre les
$k$-espaces vectoriels des première, deuxième, quatrième et cinquième
lignes sont parfaits. D'après le lemme des $5$, pour montrer que
l'accouplement entre les groupes de cohomologie de la ligne du milieu
est parfait, il suffit de montrer que les couples
$(a, a')$, $(b, b')$, $(c, c')$ et $(d, d')$
sont compatibles avec les accouplements donnés (voir ci-dessous).

\medskip\noindent
Démontrons-le pour le couple $(c, c')$. Il suffit d'observer
que l'on dispose d'un diagramme commutatif
$$
\xymatrix{
\text{Tot}(\sigma_{\leq p} \Omega^\bullet \otimes_{\underline{k}}
\sigma_{\geq d - p} \Omega^\bullet) \ar[d]_{\gamma_p} &
\text{Tot}(\sigma_{\leq p + 1} \Omega^\bullet \otimes_{\underline{k}}
\sigma_{\geq d - p} \Omega^\bullet) \ar[l] \ar[d] \\
\Omega^\bullet &
\text{Tot}(\sigma_{\leq p + 1} \Omega^\bullet \otimes_{\underline{k}}
\sigma_{\geq d - p - 1} \Omega^\bullet) \ar[l]_-{\gamma_{p + 1}}
}
$$
Ainsi, si l'on a $\alpha \in H^i(X, \sigma_{\leq p + 1}\Omega^\bullet)$
et $\beta \in H^{2d - i}(X, \sigma_{\geq d - p}\Omega^\bullet)$,
on obtient
$\gamma_p(\alpha \cup c'(\beta)) = \gamma_{p + 1}(c(\alpha) \cup \beta)$
par fonctorialité du cup-produit.

\medskip\noindent
De même, pour le couple $(b, b')$, on utilise le diagramme commutatif
$$
\xymatrix{
\text{Tot}(\sigma_{\leq p + 1} \Omega^\bullet \otimes_{\underline{k}}
\sigma_{\geq d - p - 1} \Omega^\bullet) \ar[d]_{\gamma_{p + 1}} &
\text{Tot}(\Omega^{p + 1}[-p - 1] \otimes_{\underline{k}}
\sigma_{\geq d - p - 1} \Omega^\bullet) \ar[l] \ar[d] \\
\Omega^\bullet &
\Omega^{p + 1}[-p - 1]
\otimes_{\underline{k}}
\Omega^{d - p - 1}[-d + p + 1] \ar[l]_-\wedge
}
$$
et l'on raisonne de la même manière.

\medskip\noindent
Pour le couple $(d, d')$, on utilise le diagramme commutatif
$$
\xymatrix{
\Omega^{p + 1}[-p] \otimes_{\underline{k}}
\Omega^{d - p - 1}[-d + p] \ar[d] &
\text{Tot}(\sigma_{\leq p}\Omega^\bullet \otimes_{\underline{k}}
\Omega^{d - p - 1}[-d + p]) \ar[l] \ar[d] \\
\Omega^\bullet &
\text{Tot}(\sigma_{\leq p}\Omega^\bullet \otimes_{\underline{k}}
\sigma_{\geq d - p}\Omega^\bullet) \ar[l]
}
$$
et l'on considère les classes de cohomologie dans
$H^i(X, \sigma_{\leq p}\Omega^\bullet)$ et
$H^{2d - i}(X, \Omega^{d - p - 1}[-d + p])$.
En remplaçant $i$ par $i - 1$, on obtient le résultat pour le couple $(a, a')$,
ce qui achève la démonstration du caractère parfait de nos accouplements.

\medskip\noindent
Nous omettons l'argument établissant l'unicité de $t$ à
précomposition près par la multiplication par une unité de $H^0(X, \mathcal{O}_X)$.
\end{proof}







\section{Classes de Chern}
\label{section-chern-classes}

\noindent
Les résultats établis jusqu'ici suffisent pour appliquer la discussion de
Théories cohomologiques de Weil, section \ref{weil-section-chern},
afin de construire des classes de Chern en cohomologie de de Rham.

\begin{lemma}
\label{lemma-chern-classes}
Il existe une unique règle qui associe à tout schéma quasi-compact et
quasi-séparé $X$ une classe de Chern totale
$$
c^{dR} :
K_0(\textit{Vect}(X))
\longrightarrow
\prod\nolimits_{i \geq 0} H^{2i}_{dR}(X/\mathbf{Z})
$$
possédant les propriétés suivantes :
\begin{enumerate}
\item on a $c^{dR}(\alpha + \beta) = c^{dR}(\alpha) c^{dR}(\beta)$
pour $\alpha, \beta \in K_0(\textit{Vect}(X))$,
\item si $f : X \to X'$ est un morphisme de schémas quasi-compacts et
quasi-séparés, alors $c^{dR}(f^*\alpha) = f^*c^{dR}(\alpha)$,
\item pour $\mathcal{L} \in \Pic(X)$, on a
$c^{dR}([\mathcal{L}]) = 1 + c_1^{dR}(\mathcal{L})$
\end{enumerate}
\end{lemma}

\noindent
La construction s'étend facilement à tous les schémas, mais il faut pour cela
généraliser légèrement la discussion de Théories cohomologiques de Weil,
section \ref{weil-section-chern}.

\begin{proof}
Nous allons appliquer Théories cohomologiques de Weil, proposition
\ref{weil-proposition-chern-class}, pour obtenir ce résultat.

\medskip\noindent
Soit $\mathcal{C}$ la catégorie de tous les schémas quasi-compacts et quasi-séparés.
Elle satisfait assurément aux conditions
(1), (2) et (3) (a), (b) et (c) de Théories cohomologiques de Weil,
section \ref{weil-section-chern}.

\medskip\noindent
Notre foncteur contravariant $A$ de $\mathcal{C}$ vers la
catégorie des algèbres graduées enverra $X$ sur
$A(X) = \bigoplus_{i \geq 0} H_{dR}^{2i}(X/\mathbf{Z})$,
muni de son cup-produit.
La fonctorialité est expliquée à la section \ref{section-de-rham-cohomology}
et le cup-produit à la section \ref{section-cup-product}.
Pour les applications additives $c_1^A$, on prend $c_1^{dR}$, construit
à la section \ref{section-first-chern-class}.

\medskip\noindent
On obtient en fait des algèbres commutatives, d'après
le lemme \ref{lemma-cup-product-graded-commutative},
ce qui établit l'axiome (1) pour $A$.

\medskip\noindent
Pour vérifier l'axiome (2) pour $A$, il suffit de vérifier que
$H^*_{dR}(X \coprod Y/\mathbf{Z}) =  H^*_{dR}(X/\mathbf{Z}) \times
H^*_{dR}(Y/\mathbf{Z})$.
Cela résulte du fait que la cohomologie de de Rham
se construit en prenant la cohomologie d'un faisceau d'algèbres différentielles
graduées (pour la topologie de Zariski).

\medskip\noindent
L'axiome (3) pour $A$ affirme simplement que la formation des premières classes
de Chern des modules inversibles est compatible avec les images inverses.
Cela résulte du lemme plus général \ref{lemma-pullback-c1}.

\medskip\noindent
L'axiome (4) pour $A$ est la formule du fibré projectif,
démontrée dans la proposition \ref{proposition-projective-space-bundle-formula}.

\medskip\noindent
Axiome (5). Soit $X$ un schéma quasi-compact et quasi-séparé, et
soit $\mathcal{E} \to \mathcal{F}$ une surjection de
$\mathcal{O}_X$-modules localement libres de rangs $r + 1$ et $r$. Notons
$i : P' = \mathbf{P}(\mathcal{F}) \to \mathbf{P}(\mathcal{E}) = P$
le morphisme d'inclusion correspondant. C'est un morphisme de schémas projectifs lisses
sur $X$ qui présente $P'$ comme un diviseur de Cartier effectif sur $P$.
Ainsi, d'après le lemme \ref{lemma-check-log-smooth}, le complexe à pôles logarithmiques
pour $P' \subset P$ sur $\mathbf{Z}$ est défini.
Par conséquent, pour $a \in A(P)$ tel que $i^*a = 0$, on a
$a \cup c_1^A(\mathcal{O}_P(P')) = 0$ d'après
le lemme \ref{lemma-log-complex-consequence}.
Cela achève la démonstration.
\end{proof}

\begin{remark}
\label{remark-splitting-principle}
Les analogues de Théories cohomologiques de Weil, lemmes
\ref{weil-lemma-splitting-principle} (principe de décomposition) et
\ref{weil-lemma-chern-classes-E-tensor-L} (classes de Chern des produits tensoriels),
sont valables pour les classes de Chern de de Rham sur les schémas quasi-compacts et quasi-séparés.
C'est clair, puisque nous avons montré dans la démonstration du
lemme \ref{lemma-chern-classes}
que tous les axiomes de Théories cohomologiques de Weil, section
\ref{weil-section-chern}, sont satisfaits.
\end{remark}

\noindent
Pour les schémas sur $\mathbf{Q}$, on peut construire un caractère de Chern.

\begin{lemma}
\label{lemma-chern-character}
Il existe une unique règle qui associe à tout schéma quasi-compact et quasi-séparé
$X$ sur $\mathbf{Q}$ un « caractère de Chern »
$$
ch^{dR} : K_0(\textit{Vect}(X)) \longrightarrow
\prod\nolimits_{i \geq 0} H_{dR}^{2i}(X/\mathbf{Q})
$$
possédant les propriétés suivantes :
\begin{enumerate}
\item $ch^{dR}$ est un morphisme d'anneaux pour tout $X$,
\item si $f : X' \to X$ est un morphisme de schémas quasi-compacts et quasi-séparés
sur $\mathbf{Q}$, alors $f^* \circ ch^{dR} =  ch^{dR} \circ f^*$, et
\item pour $\mathcal{L} \in \Pic(X)$,
on a $ch^{dR}([\mathcal{L}]) = \exp(c_1^{dR}(\mathcal{L}))$.
\end{enumerate}
\end{lemma}

\noindent
La construction s'étend facilement à tous les schémas sur $\mathbf{Q}$,
mais il faut pour cela généraliser légèrement la discussion de
Théories cohomologiques de Weil, section \ref{weil-section-chern}.

\begin{proof}
Exactement comme dans la démonstration du lemme \ref{lemma-chern-classes},
on montre que la catégorie des schémas quasi-compacts et quasi-séparés
sur $\mathbf{Q}$, munie du foncteur
$A^*(X) = \bigoplus_{i \geq 0} H_{dR}^{2i}(X/\mathbf{Q})$,
satisfait aux axiomes de
Théories cohomologiques de Weil, section \ref{weil-section-chern}.
De plus, dans ce cas, $A(X)$ est une $\mathbf{Q}$-algèbre pour
tout $X$. Le lemme résulte donc de
Théories cohomologiques de Weil, proposition
\ref{weil-proposition-chern-character}.
\end{proof}







\section{Une théorie cohomologique de Weil}
\label{section-weil}

\noindent
Soit $k$ un corps de caractéristique $0$. Dans cette section, nous montrons que
le foncteur
$$
X \longmapsto H^*_{dR}(X/k)
$$
définit une théorie cohomologique de Weil sur $k$ à coefficients dans $k$, au sens
de Théories cohomologiques de Weil, définition
\ref{weil-definition-weil-cohomology-theory}.
Nous vérifierons que les constructions précédentes de ce
chapitre fournissent des données (D0), (D1) et (D2') satisfaisant
aux axiomes (A1) -- (A9) de
Théories cohomologiques de Weil, section \ref{weil-section-c1}.

\medskip\noindent
Dans toute la suite de cette section, fixons le corps $k$ de caractéristique
$0$ et posons $F = k$. Prenons ensuite les données suivantes :
\begin{enumerate}
\item[(D0)] Pour l'espace vectoriel de dimension $1$ sur $F$, $F(1)$, prenons
$F(1) = F = k$.
\item[(D1)] Pour le foncteur $H^*$, prenons celui qui envoie
un schéma projectif lisse $X$ sur $k$ sur $H^*_{dR}(X/k)$.
La fonctorialité est expliquée à la section \ref{section-de-rham-cohomology}
et le cup-produit à la section \ref{section-cup-product}.
On obtient des $F$-algèbres graduées commutatives d'après
le lemme \ref{lemma-cup-product-graded-commutative}.
\item[(D2')] Pour les applications $c_1^H : \Pic(X) \to H^2(X)(1)$, on
utilise la première classe de Chern de de Rham introduite à la
section \ref{section-first-chern-class}.
\end{enumerate}
Nous allons montrer que les axiomes (A1) -- (A9) sont satisfaits.

\medskip\noindent
Dans ce paragraphe, nous allons ramener la vérification des
axiomes au cas où $k$ est algébriquement clos, en
utilisant Théories cohomologiques de Weil, lemme \ref{weil-lemma-check-over-extension}.
Notons $k'$ une clôture algébrique de $k$.
Posons $F' = k'$. On obtient des données (D0), (D1), (D2') sur $k'$ à
corps de coefficients $F'$, exactement comme ci-dessus.
D'après le lemme \ref{lemma-proper-smooth-de-Rham}, il existe des
isomorphismes fonctoriels
$$
H_{dR}^{2d}(X/k) \otimes_k k'
\longrightarrow
H_{dR}^{2d}(X_{k'}/k')
$$
pour $X$ lisse et projectif sur $k$. De plus, les diagrammes
$$
\xymatrix{
\Pic(X) \ar[r]_{c^{dR}_1} \ar[d] & H_{dR}^2(X/k) \ar[d] \\
\Pic(X_{k'}) \ar[r]^{c^{dR}_1} & H_{dR}^2(X_{k'}/k')
}
$$
sont commutatifs d'après le lemme \ref{lemma-pullback-c1}.
Cela achève la réduction.

\medskip\noindent
Supposons $k$ algébriquement clos de caractéristique nulle.
Nous allons établir les axiomes (A1) -- (A9) pour les données (D0), (D1) et (D2')
ci-dessus.

\medskip\noindent
Axiome (A1). Il faut vérifier que
$H^*_{dR}(X \coprod Y/k) =  H^*_{dR}(X/k) \times H^*_{dR}(Y/k)$.
Cela résulte du fait que la cohomologie de de Rham
se construit en prenant la cohomologie d'un faisceau d'algèbres différentielles
graduées (pour la topologie de Zariski).

\medskip\noindent
Axiome (A2). Il affirme simplement que la formation des premières classes
de Chern des modules inversibles est compatible avec les images inverses.
Cela résulte du lemme plus général \ref{lemma-pullback-c1}.

\medskip\noindent
Axiome (A3). Il résulte de la proposition
plus générale \ref{proposition-projective-space-bundle-formula}.

\medskip\noindent
Axiome (A4). Il résulte du lemme
plus général \ref{lemma-log-complex-consequence}.

\medskip\noindent
Dès à présent, en utilisant
Théories cohomologiques de Weil, lemmes \ref{weil-lemma-chern-classes} et
\ref{weil-lemma-cycle-classes}, on obtient un caractère de Chern et des
applications de classe de cycle
$$
\gamma :
\CH^*(X)
\longrightarrow
\bigoplus\nolimits_{i \geq 0} H^{2i}_{dR}(X/k)
$$
pour $X$ projectif lisse sur $k$, qui sont des homomorphismes d'anneaux gradués
compatibles avec les images inverses par les morphismes $f : X \to Y$
de schémas projectifs lisses sur $k$.

\medskip\noindent
Axiome (A5). On a $H_{dR}^*(\Spec(k)/k) = k = F$ en degré $0$.
On dispose de la formule de K\"unneth pour le produit de deux
$k$-schémas projectifs lisses d'après le lemme \ref{lemma-kunneth-de-rham} (observons que les
produits tensoriels dérivés de l'énoncé ne posent pas de difficulté, puisque la
tensorisation s'effectue sur le corps $k$).

\medskip\noindent
Axiome (A7). Il résulte de la proposition \ref{proposition-blowup-split}.

\medskip\noindent
Axiome (A8). Soit $X$ un schéma projectif lisse sur $k$.
Le texte explicatif accompagnant cet axiome dans
Théories cohomologiques de Weil, section \ref{weil-section-c1},
montre que $k' = H^0(X, \mathcal{O}_X)$ est une
$k$-algèbre finie séparable. Il en résulte que $H_{dR}^*(\Spec(k')/k) = k'$,
concentré en degré $0$, puisque $\Omega_{k'/k} = 0$. D'après
le lemme \ref{lemma-bottom-part-degenerates},
on a aussi $H_{dR}^0(X, \mathcal{O}_X) = k'$, ce qui donne
l'axiome.

\medskip\noindent
Axiome (A6). Soit $X$ un schéma projectif lisse non vide sur $k$,
équidimensionnel de dimension $d$. Notons
$\Delta : X \to X \times_{\Spec(k)} X$
le morphisme diagonal de $X$ sur $k$. Il faut montrer qu'il
existe une application $k$-linéaire
$$
\lambda : H_{dR}^{2d}(X/k) \longrightarrow k
$$
telle que $(1 \otimes \lambda)\gamma([\Delta]) = 1$ dans $H^0_{dR}(X/k)$.
Écrivons
$$
\gamma = \gamma([\Delta]) = \gamma_0 + \ldots + \gamma_{2d}
$$,
où $\gamma_i \in H_{dR}^i(X/k) \otimes_k H_{dR}^{2d - i}(X/k)$
sont les composantes de K\"unneth. Il s'agit de montrer qu'il existe une
application linéaire $\lambda : H_{dR}^{2d}(X/k) \to k$ telle que
$(1 \otimes \lambda)\gamma_0 = 1$ dans $H^0_{dR}(X/k)$.

\medskip\noindent
Soit $X = \coprod X_i$ la décomposition de $X$ en composantes connexes,
donc irréductibles. On a alors la décomposition correspondante
$\Delta = \coprod \Delta_i$ avec $\Delta_i \subset X_i \times X_i$.
Il en résulte que
$$
\gamma([\Delta]) = \sum \gamma([\Delta_i])
$$,
et, de plus, $\gamma([\Delta_i])$ correspond à la classe de
$\Delta_i \subset X_i \times X_i$ par la décomposition
$$
H^*_{dR}(X \times X) = \prod\nolimits_{i, j} H^*_{dR}(X_i \times X_j)
$$
Nous omettons les détails ; une façon de le montrer consiste à utiliser, dans
$\CH^0(X \times X)$, les idempotents $e_{i, j}$ correspondant aux
sous-schémas ouverts et fermés $X_i \times X_j$, ainsi que le fait que
$\gamma$ est un morphisme d'anneaux envoyant $e_{i, j}$ sur
l'idempotent correspondant dans la décomposition en produit de la cohomologie affichée ci-dessus.
Si l'on peut trouver $\lambda_i : H_{dR}^{2d}(X_i/k) \to k$ tel que
$(1 \otimes \lambda_i)\gamma([\Delta_i]) = 1$ dans $H^0_{dR}(X_i/k)$,
alors $\lambda = \sum \lambda_i$ résoudra le problème pour $X$.
On peut donc supposer, et l'on suppose désormais, $X$ irréductible.

\medskip\noindent
Démonstration de l'axiome (A6) pour $X$ irréductible. Puisque $k$ est algébriquement
clos, on a $H^0_{dR}(X/k) = k$, car $H^0(X, \mathcal{O}_X) = k$,
$X$ étant une variété projective sur un corps algébriquement clos
(voir par exemple Variétés, lemme
\ref{varieties-lemma-proper-geometrically-reduced-global-sections}).
Soit $x \in X$ un point fermé quelconque.
Considérons le diagramme cartésien
$$
\xymatrix{
x \ar[d] \ar[r] & X \ar[d]^\Delta \\
X \ar[r]^-{x \times \text{id}} & X \times_{\Spec(k)} X
}
$$
La compatibilité de $\gamma$ avec les images inverses implique que
$\gamma([\Delta])$ s'envoie sur $\gamma([x])$ dans $H_{dR}^{2d}(X/k)$ ;
autrement dit, $\gamma_0 = 1 \otimes \gamma([x])$.
On en déduit deux faits : (a) la classe
$\gamma([x])$ est indépendante de $x$, (b) il suffit
de montrer que $\gamma([x])$ est non nulle, et donc (c)
il suffit de trouver un cycle de dimension zéro $\alpha$ sur $X$ tel que
$\gamma(\alpha) \not = 0$. Pour cela, choisissons un morphisme
fini
$$
f : X \longrightarrow \mathbf{P}^d_k
$$
Pour l'existence d'un tel morphisme, voir
Théorie de l'intersection, section \ref{intersection-section-projection},
en particulier le lemme \ref{intersection-lemma-projection-generically-finite}.
Observons que $f$ est syntomique fini (localement d'intersection complète
d'après Compléments sur les morphismes, lemme \ref{more-morphisms-lemma-lci-permanence},
et plat d'après Algèbre, lemme \ref{algebra-lemma-CM-over-regular-flat}).
D'après la proposition \ref{proposition-Garel}, on dispose d'un morphisme trace
$$
\Theta_f :
f_*\Omega^\bullet_{X/k}
\longrightarrow
\Omega^\bullet_{\mathbf{P}^d_k/k}
$$
dont la composée avec le morphisme canonique
$$
\Omega^\bullet_{\mathbf{P}^d_k/k}
\longrightarrow
f_*\Omega^\bullet_{X/k}
$$
est la multiplication par le degré de $f$. On obtient donc une application
$$
\Theta : H_{dR}^{2d}(X/k) \to H_{dR}^{2d}(\mathbf{P}^d_k/k)
$$
telle que $\Theta \circ f^*$ soit la multiplication par un entier strictement positif.
Ainsi, si l'on peut trouver sur $\mathbf{P}^d_k$ un cycle de dimension zéro dont la classe
est non nulle, on conclut par la compatibilité de $\gamma$
avec les images inverses. Cela résulte du
lemme \ref{lemma-de-rham-cohomology-projective-space} et achève
la démonstration de l'axiome (A6).

\medskip\noindent
Nous utiliserons les faits suivants sans plus les mentionner.
Premièrement, d'après Théories cohomologiques de Weil, remarque \ref{weil-remark-trace},
l'application $\lambda_X : H^{2d}_{dR}(X/k) \to k$ est unique.
Deuxièmement, dans la démonstration de l'axiome (A6), nous avons
vu que $\lambda_X(\gamma([x])) = 1$ lorsque $X$ est irréductible ; autrement dit,
la composée de l'application de classe de cycle
$\gamma : \CH^d(X) \to H_{dR}^{2d}(X/k)$ avec $\lambda_X$
est l'application degré.

\medskip\noindent
Axiome (A9). Soit $Y \subset X$ un diviseur lisse non vide sur un
schéma projectif lisse équidimensionnel non vide $X$ sur $k$,
de dimension $d$. Il faut montrer que le diagramme
$$
\xymatrix{
H_{dR}^{2d - 2}(X/k)
\ar[rrr]_{c^{dR}_1(\mathcal{O}_X(Y)) \cap -} \ar[d]_{restriction} & & &
H_{dR}^{2d}(X) \ar[d]^{\lambda_X} \\
H_{dR}^{2d - 2}(Y/k) \ar[rrr]^-{\lambda_Y} & & & k
}
$$
est commutatif, où $\lambda_X$ et $\lambda_Y$ sont comme dans l'axiome (A6).
Nous avons vu ci-dessus que, si l'on décompose $X = \coprod X_i$ en composantes connexes
(ou, de manière équivalente, irréductibles),
on a la décomposition correspondante $\lambda_X = \sum \lambda_{X_i}$.
De même, si l'on décompose $Y = \coprod Y_j$ en composantes connexes (ou, de manière équivalente,
irréductibles), on a $\lambda_Y = \sum \lambda_{Y_j}$.
De plus, dans ce cas,
$\mathcal{O}_X(Y) = \otimes_j \mathcal{O}_X(Y_j)$, d'où
$$
c_1^{dR}(\mathcal{O}_X(Y)) = \sum\nolimits_j
c^{dR}_1(\mathcal{O}_X(Y_j))
$$
dans $H_{dR}^2(X/k)$. Une chasse au diagramme directe montre qu'il suffit
de prouver la commutativité lorsque $X$ et $Y$ sont tous deux
irréductibles. Alors $H_{dR}^{2d - 2}(Y/k)$ est de dimension $1$,
car on dispose de la dualité de Poincaré pour $Y$ d'après
Théories cohomologiques de Weil, lemme \ref{weil-lemma-poincare-duality}.
D'après l'axiome (A4), le noyau de la restriction (flèche verticale de gauche)
est contenu dans le noyau du cup-produit avec $c^{dR}_1(\mathcal{O}_X(Y))$.
Il suffit donc de trouver une classe de cohomologie
$a \in H_{dR}^{2d - 2}(X)$ dont la restriction à $Y$ soit non nulle
et pour laquelle le diagramme soit commutatif, à savoir pour $a$.
Prenons un module inversible ample quelconque $\mathcal{L}$ et posons
$$
a = c^{dR}_1(\mathcal{L})^{d - 1}
$$
On sait alors que $a|_Y = c^{dR}_1(\mathcal{L}|_Y)^{d - 1}$,
d'où
$$
\lambda_Y(a|_Y) = \deg(c_1(\mathcal{L}|_Y)^{d - 1} \cap [Y])
$$
d'après la description de $\lambda_Y$ ci-dessus. C'est un entier strictement positif,
d'après Homologie de Chow, lemme
\ref{chow-lemma-degrees-and-numerical-intersections}, combiné avec
Variétés, lemme \ref{varieties-lemma-ample-positive}.
De même, on obtient
$$
\lambda_X(c^{dR}_1(\mathcal{O}_X(Y)) \cap a) =
\deg(c_1(\mathcal{O}_X(Y)) \cap c_1(\mathcal{L})^{d - 1} \cap [X])
$$
Puisque l'on sait que $c_1(\mathcal{O}_X(Y)) \cap [X] = [Y]$, presque
par définition, on dispose d'une égalité de cycles de dimension zéro
$$
(Y \to X)_*\left(c_1(\mathcal{L}|_Y)^{d - 1} \cap [Y]\right) =
c_1(\mathcal{O}_X(Y)) \cap c_1(\mathcal{L})^{d - 1} \cap [X]
$$
sur $X$. Ces cycles ont donc le même degré, ce qui achève la démonstration.

\begin{proposition}
\label{proposition-de-rham-is-weil}
Soit $k$ un corps de caractéristique nulle. Le foncteur qui
envoie un schéma projectif lisse $X$ sur $k$ sur $H_{dR}^*(X/k)$
est une théorie cohomologique de Weil au sens de
Théories cohomologiques de Weil, définition
\ref{weil-definition-weil-cohomology-theory}.
\end{proposition}

\begin{proof}
Dans la discussion ci-dessus, nous avons montré que nos données (D0), (D1), (D2')
satisfont aux axiomes (A1) -- (A9) de Théories cohomologiques de Weil, section
\ref{weil-section-c1}. On conclut donc par
Théories cohomologiques de Weil, proposition \ref{weil-proposition-get-weil}.

\medskip\noindent
Ne lisez pas ce qui suit. Dans la démonstration des assertions, nous avons aussi utilisé les
lemmes \ref{lemma-proper-smooth-de-Rham},
\ref{lemma-pullback-c1},
\ref{lemma-log-complex-consequence},
\ref{lemma-kunneth-de-rham},
\ref{lemma-bottom-part-degenerates} et
\ref{lemma-de-rham-cohomology-projective-space},
les propositions
\ref{proposition-projective-space-bundle-formula},
\ref{proposition-blowup-split} et
\ref{proposition-Garel},
Théories cohomologiques de Weil, lemmes
\ref{weil-lemma-check-over-extension},
\ref{weil-lemma-chern-classes},
\ref{weil-lemma-cycle-classes} et
\ref{weil-lemma-poincare-duality},
Théories cohomologiques de Weil, remarque \ref{weil-remark-trace},
Variétés, lemmes
\ref{varieties-lemma-proper-geometrically-reduced-global-sections} et
\ref{varieties-lemma-ample-positive},
Théorie de l'intersection, section \ref{intersection-section-projection} et
lemme \ref{intersection-lemma-projection-generically-finite},
Compléments sur les morphismes, lemme \ref{more-morphisms-lemma-lci-permanence},
Algèbre, lemme \ref{algebra-lemma-CM-over-regular-flat}, et
Homologie de Chow, lemme
\ref{chow-lemma-degrees-and-numerical-intersections}.
\end{proof}

\begin{remark}
\label{remark-hodge-cohomology-is-weil}
Exactement de la même manière que ci-dessus, on peut montrer que
la cohomologie de Hodge $X \mapsto H_{Hodge}^*(X/k)$, munie
de $c_1^{Hodge}$, définit une théorie
cohomologique de Weil. Si nous en avons besoin, nous en donnerons ici
un énoncé précis et une démonstration. Il en résulte la conséquence
amusante suivante : si les nombres de Betti d'une théorie cohomologique
de Weil sont indépendants de la théorie cohomologique de Weil choisie
(sur notre corps $k$ de caractéristique $0$), alors
la suite spectrale de Hodge vers de Rham
dégénère en $E_1$ ! Bien entendu, la dégénérescence de
la suite spectrale de Hodge vers de Rham est connue
(voir par exemple \cite{Deligne-Illusie} pour une remarquable démonstration algébrique),
mais ce résultat n'est nullement facile ! Cela suggère que l'indépendance
des nombres de Betti est également difficile à démontrer
et, à notre connaissance, reste un problème ouvert. Voir
Théories cohomologiques de Weil, remarque
\ref{weil-remark-betti-numbers-in-some-sense} pour une question apparentée.
\end{remark}







\section{Morphismes de Gysin pour les immersions fermées}
\label{section-gysin}

\noindent
Dans cette section, nous définissons le morphisme de Gysin pour les immersions fermées.

\begin{remark}
\label{remark-gysin-equations}
Soit $X \to S$ un morphisme de schémas. Soient
$f_1, \ldots, f_c \in \Gamma(X, \mathcal{O}_X)$ et $Z \subset X$
le sous-schéma fermé défini par $f_1, \ldots, f_c$. Nous étudierons ci-dessous le {\it morphisme de Gysin}
\begin{equation}
\label{equation-gysin}
\gamma^p_{f_1, \ldots, f_c} :
\Omega^p_{Z/S}
\longrightarrow
\mathcal{H}_Z^c(\Omega^{p + c}_{X/S})
\end{equation}
défini comme suit. Étant donnée une section locale $\omega$ de $\Omega^p_{Z/S}$
qui est la restriction d'une section $\tilde \omega$ de $\Omega^p_{X/S}$,
nous posons
$$
\gamma^p_{f_1, \ldots, f_c}(\omega) =
c_{f_1, \ldots, f_c}(\tilde \omega|_Z) \wedge
\text{d}f_1 \wedge \ldots \wedge \text{d}f_c
$$
où $c_{f_1, \ldots, f_c} : \Omega^p_{X/S} \otimes \mathcal{O}_Z \to
\mathcal{H}_Z^c(\Omega^p_{X/S})$ est le morphisme construit dans
Catégories dérivées des schémas, remarque
\ref{perfect-remark-supported-map-c-equations}.
Cette définition est indépendante des choix : pour $\omega$ donné, nous pouvons modifier le choix de
$\tilde \omega$ par des éléments de la forme
$\sum f_i \omega'_i + \sum \text{d}(f_i) \wedge \omega''_i$,
que la construction envoie sur zéro.
\end{remark}

\begin{lemma}
\label{lemma-gysin-differential}
Le morphisme de Gysin (\ref{equation-gysin}) est compatible avec les différentielles de de Rham
sur $\Omega^\bullet_{X/S}$ et $\Omega^\bullet_{Z/S}$.
\end{lemma}

\begin{proof}
Ceci résulte d'un calcul presque immédiat, une fois l'énoncé correctement interprété.
Rappelons d'abord que le foncteur
$\mathcal{H}^c_Z$ calculé dans la catégorie des $\mathcal{O}_X$-modules
coïncide avec le foncteur défini de manière analogue dans la catégorie des faisceaux abéliens
sur $X$ ; voir
Cohomologie, lemme \ref{cohomology-lemma-sections-support-abelian-unbounded}.
La différentielle $\text{d} : \Omega^p_{X/S} \to \Omega^{p + 1}_{X/S}$
induit donc un morphisme
$\mathcal{H}^c_Z(\Omega^p_{X/S}) \to \mathcal{H}^c_Z(\Omega^{p + 1}_{X/S})$.
De plus, la construction du complexe de {\v C}ech alterné augmenté dans
Catégories dérivées des schémas, remarque \ref{perfect-remark-support-c-equations},
s'applique à la catégorie des faisceaux abéliens. Le morphisme
$$
\Coker\left(\bigoplus \mathcal{F}_{1 \ldots \hat i \ldots c} \to
\mathcal{F}_{1 \ldots c}\right)
\longrightarrow
i_*\mathcal{H}^c_Z(\mathcal{F})
$$
utilisé pour construire $c_{f_1, \ldots, f_c}$ dans
Catégories dérivées des schémas, remarque
\ref{perfect-remark-supported-map-c-equations},
est bien défini et
fonctoriel dans la catégorie de tous les faisceaux abéliens sur $X$.
Le lemme résulte donc de l'égalité
$$
\text{d}\left(
\frac{\tilde \omega \wedge \text{d}f_1 \wedge \ldots \wedge
\text{d}f_c}{f_1 \ldots f_c}\right) =
\frac{\text{d}(\tilde \omega) \wedge
\text{d}f_1 \wedge \ldots \wedge \text{d}f_c}{f_1 \ldots f_c}
$$,
qui est claire.
\end{proof}

\begin{lemma}
\label{lemma-gysin-global}
Soit $X \to S$ un morphisme de schémas. Soit $Z \to X$ une immersion fermée
de présentation finie dont le faisceau conormal $\mathcal{C}_{Z/X}$ est
localement libre de rang $c$. Il existe alors un morphisme canonique
$$
\gamma^p : \Omega^p_{Z/S} \to \mathcal{H}^c_Z(\Omega^{p + c}_{X/S})
$$
donné localement par les morphismes $\gamma^p_{f_1, \ldots, f_c}$
de la remarque \ref{remark-gysin-equations}.
\end{lemma}

\begin{proof}
Les hypothèses impliquent que, pour tout $x \in Z \subset X$, il existe un
voisinage ouvert $U$ de $x$ sur lequel $Z$ est défini par $c$
éléments $f_1, \ldots, f_c \in \mathcal{O}_X(U)$. Il suffit donc de montrer que, si $f_1, \ldots, f_c$ et
$g_1, \ldots, g_c$ sont des éléments de $\mathcal{O}_X(U)$ définissant $Z \cap U$,
les morphismes $\gamma^p_{f_1, \ldots, f_c}$
et $\gamma^p_{g_1, \ldots, g_c}$ coïncident. Pour cela, quitte à restreindre
$U$, nous pouvons supposer que $g_j = \sum a_{ji} f_i$ pour certains
$a_{ji} \in \mathcal{O}_X(U)$. Nous avons alors
$c_{f_1, \ldots, f_c} = \det(a_{ji}) c_{g_1, \ldots, g_c}$ d'après
Catégories dérivées des schémas, lemme
\ref{perfect-lemma-supported-map-determinant}.
D'autre part, nous avons
$$
\text{d}(g_1) \wedge \ldots \wedge \text{d}(g_c) \equiv
\det(a_{ji}) \text{d}(f_1) \wedge \ldots \wedge \text{d}(f_c)
\bmod (f_1, \ldots, f_c)\Omega^c_{X/S}
$$
En combinant ces relations, un calcul direct donne
l'égalité voulue.
\end{proof}

\begin{lemma}
\label{lemma-gysin-differential-global}
Soient $X \to S$ et $i : Z \to X$ comme dans le lemme \ref{lemma-gysin-global}.
Le morphisme de Gysin $\gamma^p$ est compatible avec les différentielles de de Rham
sur $\Omega^\bullet_{X/S}$ et $\Omega^\bullet_{Z/S}$.
\end{lemma}

\begin{proof}
La vérification est locale, et l'assertion résulte alors du
lemme \ref{lemma-gysin-differential}.
\end{proof}

\begin{lemma}
\label{lemma-gysin-projection}
Soient $X \to S$ et $i : Z \to X$ comme dans le lemme \ref{lemma-gysin-global}.
Pour $\alpha \in H^q(X, \Omega^p_{X/S})$, nous avons
$\gamma^p(\alpha|_Z) = i^{-1}\alpha \wedge \gamma^0(1)$ dans
$H^q(Z, \mathcal{H}^c_Z(\Omega^{p + c}_{X/S}))$.
Les notations sont expliquées dans la démonstration.
\end{lemma}

\begin{proof}
La restriction $\alpha|_Z$ est l'élément de $H^q(Z, \Omega^p_{Z/S})$
donné par la fonctorialité de la cohomologie de Hodge. En appliquant la fonctorialité
de la cohomologie au morphisme
$\gamma^p : \Omega^p_{Z/S} \to \mathcal{H}^c_Z(\Omega^{p + c}_{X/S})$,
nous obtenons $\gamma^p(\alpha|_Z)$ dans
$H^q(Z, \mathcal{H}^c_Z(\Omega^{p + c}_{X/S}))$.
Cela explique le membre de gauche de la formule.

\medskip\noindent
Pour expliquer le membre de droite, nous prenons d'abord l'image inverse par le morphisme
d'espaces annelés $i : (Z, i^{-1}\mathcal{O}_X) \to (X, \mathcal{O}_X)$,
ce qui donne l'élément $i^{-1}\alpha \in H^q(Z, i^{-1}\Omega^p_{X/S})$.
Soit $\gamma^0(1) \in H^0(Z, \mathcal{H}_Z^c(\Omega^c_{X/S}))$
l'image de $1 \in H^0(Z, \mathcal{O}_Z) = H^0(Z, \Omega^0_{Z/S})$
par $\gamma^0$. Le cup-produit fournit un élément
$$
i^{-1}\alpha \cup \gamma^0(1)
\in
H^{q + c}(Z, 
i^{-1}\Omega^p_{X/S} \otimes_{i^{-1}\mathcal{O}_X}
\mathcal{H}^c_Z(\Omega^c_{X/S}))
$$
En utilisant Cohomologie, remarque \ref{cohomology-remark-support-cup-product},
et le produit extérieur, nous obtenons des morphismes canoniques
$$
i^{-1}\Omega^p_{X/S} \otimes_{i^{-1}\mathcal{O}_X}^\mathbf{L}
R\mathcal{H}_Z(\Omega^c_{X/S}) \to
R\mathcal{H}_Z(\Omega^p_{X/S} \otimes_{\mathcal{O}_X}^\mathbf{L}
\Omega^c_{X/S}) \to
R\mathcal{H}_Z(\Omega^{p + c}_{X/S})
$$
D'après Catégories dérivées des schémas, lemme
\ref{perfect-lemma-supported-trivial-vanishing},
les faisceaux de cohomologie des objets $R\mathcal{H}_Z(\Omega^j_{X/S})$ sont nuls
en degrés $> c$. Ainsi, en prenant les faisceaux de cohomologie
en degré $c$, nous obtenons un morphisme
$$
i^{-1}\Omega^p_{X/S} \otimes_{i^{-1}\mathcal{O}_X}
\mathcal{H}^c_Z(\Omega^c_{X/S}) \longrightarrow
\mathcal{H}^c_Z(\Omega^{p + c}_{X/S})
$$
L'expression $i^{-1}\alpha \wedge \gamma^0(1)$ est l'image
du cup-produit $i^{-1}\alpha \cup \gamma^0(1)$ par la
fonctorialité de la cohomologie.

\medskip\noindent
Le sens de la formule étant ainsi précisé, les
propriétés générales des cup-produits
(Cohomologie, section \ref{cohomology-section-cup-product})
montrent qu'il suffit de prouver que le diagramme
$$
\xymatrix{
i^{-1}\Omega^p_X \otimes \Omega^0_Z \ar[rr]_{\text{id} \otimes \gamma^0}
\ar[d] & &
i^{-1}\Omega^p_X \otimes \mathcal{H}^c_Z(\Omega^c_X) \ar[d]^\wedge \\
\Omega^p_Z \otimes \Omega^0_Z \ar[r]^\wedge &
\Omega^p_Z \ar[r]^{\gamma^p} &
\mathcal{H}^c_Z(\Omega^{p + c}_X)
}
$$
est commutatif dans la catégorie des faisceaux sur $Z$ (avec un abus de
notation évident). Cela se ramène à un calcul simple pour les morphismes
$\gamma^j_{f_1, \ldots, f_c}$, que nous omettons ; ces morphismes sont en fait
choisis précisément pour assurer cette commutativité et pour que $1$ ait pour image
$\frac{\text{d}f_1 \wedge \ldots \wedge \text{d}f_c}{f_1 \ldots f_c}$.
\end{proof}

\begin{lemma}
\label{lemma-gysin-transverse}
Soit $c \geq 0$ un entier. Soit
$$
\xymatrix{
Z' \ar[d]_h \ar[r] & X' \ar[d]_g \ar[r] & S' \ar[d] \\
Z \ar[r] & X \ar[r] & S
}
$$
un diagramme commutatif de schémas.
Supposons que
\begin{enumerate}
\item $Z \to X$ et $Z' \to X'$
vérifient les hypothèses du lemme \ref{lemma-gysin-global} ;
\item le carré de gauche du diagramme soit cartésien ;
\item $h^*\mathcal{C}_{Z/X} \to \mathcal{C}_{Z'/X'}$
(Morphismes, lemme \ref{morphisms-lemma-conormal-functorial})
soit un isomorphisme.
\end{enumerate}
Alors le diagramme
$$
\xymatrix{
h^*\Omega^p_{Z/S} \ar[rr]_-{h^{-1}\gamma^p} \ar[d] & &
\mathcal{O}_{X'}|_{Z'} \otimes_{h^{-1}\mathcal{O}_X|_Z}
h^{-1}\mathcal{H}^c_Z(\Omega^{p + c}_{X/S}) \ar[d] \\
\Omega^p_{Z'/S'} \ar[rr]^{\gamma^p} & &
\mathcal{H}^c_{Z'}(\Omega^{p + c}_{X'/S'})
}
$$
est commutatif. La flèche verticale de gauche est donnée par la fonctorialité des modules de
différentielles et celle de droite fait intervenir
Cohomologie, remarque \ref{cohomology-remark-support-functorial}.
\end{lemma}

\begin{proof}
Plus précisément, considérons le composé
\begin{align*}
\mathcal{O}_{X'}|_{Z'} \otimes_{h^{-1}\mathcal{O}_X|_Z}^\mathbf{L}
h^{-1}R\mathcal{H}_Z(\Omega^{p + c}_{X/S})
& \to
R\mathcal{H}_{Z'}(Lg^*\Omega^{p + c}_{X/S}) \\
& \to
R\mathcal{H}_{Z'}(g^*\Omega^{p + c}_{X/S}) \\
& \to
R\mathcal{H}_{Z'}(\Omega^{p + c}_{X'/S'})
\end{align*}
dont la première flèche est donnée par
Cohomologie, remarque \ref{cohomology-remark-support-functorial},
et la dernière par la fonctorialité des différentielles.
L'annulation des faisceaux de cohomologie en degrés $> c$,
d'après Catégories dérivées des schémas, lemme
\ref{perfect-lemma-supported-trivial-vanishing},
fournit la flèche verticale de droite.
Nous pouvons vérifier la commutativité localement.
Nous pouvons donc supposer que $Z$ est défini par
$f_1, \ldots, f_c \in \Gamma(X, \mathcal{O}_X)$.
Alors $Z'$ est défini par $f'_i = g^\sharp(f_i)$.
Les morphismes $c_{f_1, \ldots, f_c}$ et $c_{f'_1, \ldots, f'_c}$
s'insèrent dans le diagramme commutatif
$$
\xymatrix{
h^*i^*\Omega^p_{X/S} \ar[rr]_-{h^{-1}c_{f_1, \ldots, f_c}} \ar[d] & &
\mathcal{O}_{X'}|_{Z'} \otimes_{h^{-1}\mathcal{O}_X|_Z}
h^{-1}\mathcal{H}^c_Z(\Omega^p_{X/S}) \ar[d] \\
(i')^*\Omega^p_{X'/S'} \ar[rr]^{c_{f'_1, \ldots, f'_c}} & &
\mathcal{H}^c_{Z'}(\Omega^p_{X'/S'})
}
$$
Voir Catégories dérivées des schémas, remarque
\ref{perfect-remark-supported-functorial}.
Rappelons que, pour une $p$-forme $\omega$ sur $Z$, nous définissons
$\gamma^p(\omega)$ en choisissant (localement sur $X$ et $Z$)
une $p$-forme $\tilde \omega$ sur $X$ relevant $\omega$, puis en posant
$\gamma^p(\omega) =
c_{f_1, \ldots, f_c}(\tilde \omega) \wedge
\text{d}f_1 \wedge \ldots \wedge \text{d}f_c$.
Puisque l'image inverse de la forme $\text{d}f_1 \wedge \ldots \wedge \text{d}f_c$
est
$\text{d}f'_1 \wedge \ldots \wedge \text{d}f'_c$, nous concluons.
\end{proof}

\begin{remark}
\label{remark-how-to-use}
Soient $X \to S$, $i : Z \to X$ et $c \geq 0$ comme dans le
lemme \ref{lemma-gysin-global}.
Soit $p \geq 0$ ; supposons que $\mathcal{H}^i_Z(\Omega^{p + c}_{X/S}) = 0$
pour $i = 0, \ldots, c - 1$. Cette annulation a lieu si $X \to S$ est lisse
et si $Z \to X$ est une immersion régulière au sens de Koszul ; voir
Catégories dérivées des schémas, lemme \ref{perfect-lemma-supported-vanishing}.
Nous obtenons alors un morphisme
$$
\gamma^{p, q} :
H^q(Z, \Omega^p_{Z/S})
\longrightarrow
H^{q + c}(X, \Omega^{p + c}_{X/S})
$$
en utilisant d'abord
$\gamma^p : \Omega^p_{Z/S} \to \mathcal{H}^c_Z(\Omega^{p + c}_{X/S})$
pour atteindre
$$
H^q(Z, \mathcal{H}^c_Z(\Omega^{p + c}_{X/S})) =
H^q(Z, R\mathcal{H}_Z(\Omega^{p + c}_{X/S})[c]) =
H^q(X, i_*R\mathcal{H}_Z(\Omega^{p + c}_{X/S})[c])
$$,
puis le morphisme d'adjonction
$i_*R\mathcal{H}_Z(\Omega^{p + c}_{X/S}) \to \Omega^{p + c}_{X/S}$
pour passer au module de cohomologie de Hodge voulu.
\end{remark}

\begin{lemma}
\label{lemma-gysin-differential-hodge}
Soient $X \to S$ et $i : Z \to X$ comme dans le lemme \ref{lemma-gysin-global}.
Supposons que $X \to S$ soit lisse et que $Z \to X$ soit régulière au sens de Koszul.
Les morphismes de Gysin $\gamma^{p, q}$ sont compatibles avec les différentielles de de Rham
sur $\Omega^\bullet_{X/S}$ et $\Omega^\bullet_{Z/S}$.
\end{lemma}

\begin{proof}
Cela résulte immédiatement du
lemme \ref{lemma-gysin-differential-global}.
\end{proof}

\begin{lemma}
\label{lemma-gysin-projection-global}
Soient $X \to S$, $i : Z \to X$ et $c \geq 0$ comme dans le
lemme \ref{lemma-gysin-global}. Supposons que $X \to S$ soit lisse et
que $Z \to X$ soit régulière au sens de Koszul. Pour $\alpha \in H^q(X, \Omega^p_{X/S})$, nous avons
$\gamma^{p, q}(\alpha|_Z) = \alpha \cup \gamma^{0, 0}(1)$ dans
$H^{q + c}(X, \Omega^{p + c}_{X/S})$, où $\gamma^{a, b}$ est défini dans la
remarque \ref{remark-how-to-use}.
\end{lemma}

\begin{proof}
Ce lemme résulte du lemme \ref{lemma-gysin-projection}
et de Cohomologie, lemme \ref{cohomology-lemma-support-cup-product}.
Le lecteur peut omettre la discussion plus détaillée qui suit.

\medskip\noindent
Nous utiliserons sans le rappeler l'égalité
$R\mathcal{H}_Z(\Omega^j_{X/S}) = \mathcal{H}^c_Z(\Omega^j_{X/S})[-c]$
pour tout $j$, signalée dans la remarque \ref{remark-how-to-use}.
Nous utiliserons également sans les mentionner les identifications
$H^{q + c}_Z(X, \Omega^j_{X/S}) = H^{q + c}(Z, R\mathcal{H}_Z(\Omega^j_{X/S}) =
H^q(Z, \mathcal{H}^c_Z(\Omega^j_{X/S}))$ ; voir
Cohomologie, lemme \ref{cohomology-lemma-local-to-global-sections-with-support},
pour la première. Avec ces identifications :
\begin{enumerate}
\item $\gamma^0(1) \in H^c_Z(X, \Omega^c_{X/S})$ a pour image $\gamma^{0, 0}(1)$
dans $H^c(X, \Omega^c_{X/S})$ ;
\item le membre de droite $i^{-1}\alpha \wedge \gamma^0(1)$
de l'égalité du lemme \ref{lemma-gysin-projection}
est l'image par le produit extérieur du cup-produit de
Cohomologie, remarque \ref{cohomology-remark-support-cup-product-global},
des éléments $\alpha$ et $\gamma^0(1)$ ; autrement dit, les constructions
de la démonstration du lemme \ref{lemma-gysin-projection} et de
Cohomologie, remarque \ref{cohomology-remark-support-cup-product-global}, coïncident ;
\item d'après Cohomologie, lemme \ref{cohomology-lemma-support-cup-product},
cet élément a pour image $\alpha \cup \gamma^{0, 0}(1)$ dans
$H^{q + c}(X, \Omega^p_{X/S} \otimes \Omega^c_{X/S})$ ;
\item le membre de gauche $\gamma^p(\alpha|_Z)$ de l'égalité du
lemme \ref{lemma-gysin-projection} a pour image
$\gamma^{p, q}(\alpha|_Z)$.
\end{enumerate}
Cela achève la démonstration.
\end{proof}

\begin{lemma}
\label{lemma-gysin-transverse-global}
Soient $c \geq 0$ et
$$
\xymatrix{
Z' \ar[d]_h \ar[r] & X' \ar[d]_g \ar[r] & S' \ar[d] \\
Z \ar[r] & X \ar[r] & S
}
$$
vérifiant les hypothèses du lemme \ref{lemma-gysin-transverse}. Supposons de plus
que $X \to S$ et $X' \to S'$ soient lisses, et que
$Z \to X$ et $Z' \to X'$ soient des immersions régulières au sens de Koszul.
Alors le diagramme
$$
\xymatrix{
H^q(Z, \Omega^p_{Z/S}) \ar[rr]_-{\gamma^{p, q}} \ar[d] & &
H^{q + c}(X, \Omega^{p + c}_{X/S}) \ar[d] \\
H^q(Z', \Omega^p_{Z'/S'}) \ar[rr]^{\gamma^{p, q}} & &
H^{q + c}(X', \Omega^{p + c}_{X'/S'})
}
$$
est commutatif, où $\gamma^{p, q}$ est défini dans la remarque \ref{remark-how-to-use}.
\end{lemma}

\begin{proof}
Cela résulte de la combinaison du lemme \ref{lemma-gysin-transverse}
et de Cohomologie, lemme \ref{cohomology-lemma-support-functorial}.
\end{proof}

\begin{lemma}
\label{lemma-class-of-a-point}
Soit $k$ un corps. Soit $X$ un schéma irréductible, lisse et propre sur $k$,
de dimension $d$. Soit $Z \subset X$ le sous-schéma fermé réduit constitué
d'un unique point $k$-rationnel $x$. Alors l'image de
$1 \in k = H^0(Z, \mathcal{O}_Z) = H^0(Z, \Omega^0_{Z/k})$
par le morphisme $H^0(Z, \Omega^0_{Z/k}) \to H^d(X, \Omega^d_{X/k})$
de la remarque \ref{remark-how-to-use} est non nulle.
\end{lemma}

\begin{proof}
Le morphisme $\gamma^0 : \mathcal{O}_Z \to
\mathcal{H}^d_Z(\Omega^d_{X/k}) = R\mathcal{H}_Z(\Omega^d_{X/k})[d]$
a pour adjoint un morphisme
$$
g^0 : i_*\mathcal{O}_Z \longrightarrow \Omega^d_{X/k}[d]
$$
dans $D(\mathcal{O}_X)$. Rappelons que $\Omega^d_{X/k} = \omega_X$ est un
faisceau dualisant pour $X/k$ ; voir
Dualité pour les schémas, lemme \ref{duality-lemma-duality-proper-over-field}.
Le dual $k$-linéaire du morphisme de l'énoncé
est donc le morphisme
$$
H^0(X, \mathcal{O}_X) \to \Ext^d_X(i_*\mathcal{O}_Z, \omega_X)
$$
qui envoie $1$ sur $g^0$. Il suffit ainsi de montrer que $g^0$ est non nul.
Nous pouvons le vérifier dans un voisinage quelconque $U$ du point $x$. Choisissons $U$
de sorte qu'il existe $f_1, \ldots, f_d \in \mathcal{O}_X(U)$
s'annulant seulement en $x$ et engendrant l'idéal maximal
$\mathfrak m_x \subset \mathcal{O}_{X, x}$. Nous pouvons supposer
que $U = \Spec(R)$ est affine. En reprenant la
construction de $\gamma^0$, nous voyons que notre extension est donnée par
$$
k \to
(R \to \bigoplus\nolimits_{i_0} R_{f_{i_0}} \to
\bigoplus\nolimits_{i_0 < i_1} R_{f_{i_0}f_{i_1}} \to
\ldots \to R_{f_1\ldots f_r})[d] \to R[d]
$$,
où $1$ a pour image $1/f_1 \ldots f_c$ par le premier morphisme.
Cette extension est non nulle, car $1/f_1 \ldots f_c$ est, dans ce cas, un élément non nul
du groupe de cohomologie locale $H^d_{(f_1, \ldots, f_d)}(R)$.
\end{proof}






\section{Dualité de Poincaré relative}
\label{section-relative-poincare-duality}

\noindent
Dans cette section, nous démontrons la dualité de Poincaré pour la cohomologie de de Rham relative
d'un schéma propre et lisse sur une base. Nous recommandons vivement au lecteur
de consulter d'abord la section \ref{section-poincare-duality}.

\begin{situation}
\label{situation-relative-duality}
Ici, $S$ est un schéma quasi-compact et quasi-séparé, et
$f : X \to S$ est un morphisme propre et lisse de schémas dont toutes les
fibres sont non vides et équidimensionnelles de dimension $n$.
\end{situation}

\begin{lemma}
\label{lemma-relative-bottom-part-degenerates}
Dans la situation \ref{situation-relative-duality}, l'image directe
$f_*\mathcal{O}_X$ est une $\mathcal{O}_S$-algèbre finie étale
et, localement sur $S$, nous avons $Rf_*\mathcal{O}_X = f_*\mathcal{O}_X \oplus P$
dans $D(\mathcal{O}_S)$, où $P$ est parfait d'amplitude de Tor dans $[1, \infty)$.
Le morphisme $\text{d} : f_*\mathcal{O}_X \to f_*\Omega_{X/S}$ est nul.
\end{lemma}

\begin{proof}
La première partie de l'énoncé résulte de
Catégories dérivées des schémas, lemme \ref{perfect-lemma-proper-flat-geom-red}.
En posant $S' = \underline{\Spec}_S(f_*\mathcal{O}_X)$, nous obtenons une factorisation
$X \to S' \to S$ (c'est la factorisation de Stein ; voir
Compléments sur les morphismes, section \ref{more-morphisms-section-stein-factorization},
bien que nous n'en ayons pas besoin ici),
et nous voyons que $\Omega_{X/S} = \Omega_{X/S'}$, par exemple d'après
Morphismes, lemmes \ref{morphisms-lemma-triangle-differentials} et
\ref{morphisms-lemma-etale-at-point}. Cela implique évidemment que
$\text{d} : f_*\mathcal{O}_X \to f_*\Omega_{X/S}$ est nul.
\end{proof}

\begin{lemma}
\label{lemma-relative-duality-hodge}
Dans la situation \ref{situation-relative-duality}, il existe un morphisme de
$\mathcal{O}_S$-modules
$$
t : Rf_*\Omega^n_{X/S}[n] \longrightarrow \mathcal{O}_S
$$,
unique à précomposition près par la multiplication par un élément inversible de
$H^0(X, \mathcal{O}_X)$, ayant la propriété suivante : pour tout $p$, l'accouplement
$$
Rf_*\Omega^p_{X/S}
\otimes_{\mathcal{O}_S}^\mathbf{L}
Rf_*\Omega^{n - p}_{X/S}[n]
\longrightarrow
\mathcal{O}_S
$$
donné par le cup-produit relatif composé avec $t$
est un accouplement parfait de complexes parfaits sur $S$.
\end{lemma}

\begin{proof}
Supposons d'abord que $S$ soit un schéma noethérien.
Soit $\omega^\bullet_{X/S}$ le complexe dualisant relatif de $X$ sur $S$, comme
dans Dualité pour les schémas, remarque \ref{duality-remark-relative-dualizing-complex},
et soit $Rf_*\omega_{X/S}^\bullet \to \mathcal{O}_S$ son morphisme trace. D'après
Dualité pour les schémas, lemme \ref{duality-lemma-smooth-proper}
(c'est ici que nous utilisons que $S$ est noethérien),
il existe un isomorphisme $\omega^\bullet_{X/S} \cong \Omega^n_{X/S}[n]$,
qui nous permet d'obtenir $t$. Les complexes $Rf_*\Omega^p_{X/S}$
sont parfaits d'après le lemme \ref{lemma-proper-smooth-de-Rham}.
Puisque $\Omega^p_{X/S}$ est localement libre et que
$\Omega^p_{X/S} \otimes_{\mathcal{O}_X} \Omega^{n - p}_{X/S} \to
\Omega^n_{X/S}$ définit un isomorphisme $\Omega^p_{X/S} \cong
\SheafHom_{\mathcal{O}_X}(\Omega^{n - p}_{X/S}, \Omega^n_{X/S})$,
l'accouplement induit par le cup-produit relatif est parfait d'après
Dualité pour les schémas, remarque
\ref{duality-remark-relative-dualizing-complex-relative-cup-product}.

\medskip\noindent
Démonstration de l'existence dans le cas général.
Par approximation noethérienne absolue, nous pouvons trouver un diagramme cartésien
$$
\xymatrix{
X \ar[r] \ar[d]^f & X_0 \ar[d]^{f_0} \\
S \ar[r]^g & S_0
}
$$
avec $S_0$ noethérien et $f_0$ propre et lisse,
dont toutes les fibres sont non vides et équidimensionnelles de dimension $n$.
Voir Limites, lemmes \ref{limits-lemma-descend-finite-presentation},
\ref{limits-lemma-descend-smooth}, \ref{limits-lemma-descend-dimension-d},
\ref{limits-lemma-descend-surjective} et
\ref{limits-lemma-eventually-proper}. Choisissons un morphisme
$$
t_0 : Rf_{0, *}\Omega^n_{X_0/S_0}[n] \longrightarrow \mathcal{O}_{S_0}
$$
comme dans l'énoncé. Par cohomologie et changement de base, les
complexes $Rf_*(\Omega^p_{X/S})$ s'identifient à
$Lg^*(Rf_{0, *}\Omega^p_{X_0/S_0})$ pour tout $p \geq 0$ ; voir
le lemme \ref{lemma-proper-smooth-de-Rham}. En particulier, l'image inverse
de $t_0$ par $g$ nous fournit un morphisme $t$ comme dans l'énoncé.
En effet, d'après Cohomologie, lemme \ref{cohomology-lemma-base-change-cup-product},
le morphisme de cup-produit
$$
Rf_*\Omega^p_{X/S}
\otimes_{\mathcal{O}_S}^\mathbf{L}
Rf_*\Omega^{n - p}_{X/S}[n]
\longrightarrow
Rf_*\Omega^n_{X/S}
$$
est l'image inverse par $g$ du morphisme de cup-produit correspondant pour $f_0$.
La propriété selon laquelle $t_0$ donne des accouplements parfaits se transporte donc par image inverse
en la même propriété pour $t$ sur $S$.

\medskip\noindent
Unicité de $t$. Choisissons un triangle distingué
$f_*\mathcal{O}_X \to Rf_*\mathcal{O}_X \to P \to f_*\mathcal{O}_X[1]$.
D'après le lemme \ref{lemma-relative-bottom-part-degenerates},
l'objet $P$ est parfait d'amplitude de Tor dans $[1, \infty)$,
et le triangle est scindé localement sur $S$.
Ainsi, $R\SheafHom_{\mathcal{O}_X}(P, \mathcal{O}_X)$ est parfait
d'amplitude de Tor dans $(-\infty, -1]$. La dualité ci-dessus montre donc que,
localement sur $S$, nous avons
$$
Rf_*\Omega^n_{X/S}[n] \cong
R\SheafHom_{\mathcal{O}_S}(f_*\mathcal{O}_X, \mathcal{O}_S)
\oplus R\SheafHom_{\mathcal{O}_X}(P, \mathcal{O}_X)
$$
Il en résulte que $R^nf_*\Omega^n_{X/S}$ est localement libre de rang fini et
que nous obtenons un accouplement $\mathcal{O}_S$-bilinéaire parfait
$$
f_*\mathcal{O}_X \times R^nf_*\Omega^n_{X/S} \longrightarrow \mathcal{O}_S
$$
au moyen de $t$.
Cela implique que tout morphisme $\mathcal{O}_S$-linéaire
$t' : R^nf_*\Omega^n_{X/S} \to \mathcal{O}_S$ est de la forme
$t' = t \circ g$ pour un certain
$g \in \Gamma(S, f_*\mathcal{O}_X) = \Gamma(X, \mathcal{O}_X)$.
Pour que $t'$ définisse encore un accouplement parfait, $g$ doit
être inversible. Cela achève la démonstration.
\end{proof}

\begin{lemma}
\label{lemma-relative-top-part-degenerates}
Dans la situation \ref{situation-relative-duality}, le morphisme
$\text{d} : R^nf_*\Omega^{n - 1}_{X/S} \to R^nf_*\Omega^n_{X/S}$
est nul.
\end{lemma}

\noindent
Comme nous l'avons signalé dans la démonstration du lemme \ref{lemma-top-part-degenerates},
ce lemme n'est pas une conséquence immédiate des lemmes
\ref{lemma-relative-duality-hodge} et
\ref{lemma-relative-bottom-part-degenerates}.

\begin{proof}[Démonstration lorsque $S$ est réduit]
Supposons $S$ réduit. Observons que
$\text{d} : R^nf_*\Omega^{n - 1}_{X/S} \to R^nf_*\Omega^n_{X/S}$
est un morphisme $\mathcal{O}_S$-linéaire de $\mathcal{O}_S$-modules quasi-cohérents.
Le $\mathcal{O}_S$-module $R^nf_*\Omega^n_{X/S}$ est localement libre de rang fini
(car il est dual du $\mathcal{O}_S$-module localement libre de rang fini
$f_*\mathcal{O}_X$, d'après les lemmes
\ref{lemma-relative-duality-hodge} et
\ref{lemma-relative-bottom-part-degenerates}).
Puisque $S$ est réduit, il suffit de montrer que
le germe de $\text{d}$ en chaque point générique $\eta \in S$
est nul ; on le voit en considérant les sections sur les ouverts affines,
en utilisant que le but de $\text{d}$ est localement libre et en appliquant
Algèbre, lemme \ref{algebra-lemma-reduced-ring-sub-product-fields} (2).
Puisque $S$ est réduit, nous avons $\mathcal{O}_{S, \eta} = \kappa(\eta)$ ; voir
Algèbre, lemme \ref{algebra-lemma-minimal-prime-reduced-ring}.
Ainsi, $\text{d}_\eta$ s'identifie au morphisme
$$
\text{d} :
H^n(X_\eta, \Omega^{n - 1}_{X_\eta/\kappa(\eta)})
\longrightarrow
H^n(X_\eta, \Omega^n_{X_\eta/\kappa(\eta)})
$$,
qui est nul d'après le lemme \ref{lemma-top-part-degenerates}.
\end{proof}

\begin{proof}[Démonstration dans le cas général]
La question est locale sur $S$ pour la topologie plate : si $S' \to S$ est un morphisme
plat surjectif de schémas et si le morphisme considéré est nul après image inverse sur $S'$,
alors il est nul. De plus, sa formation commute aux changements de base
plats par le théorème de changement de base plat (Cohomologie des schémas, lemme
\ref{coherent-lemma-flat-base-change-cohomology}).

\medskip\noindent
Considérons la factorisation de Stein $X \to S' \to S$ de
Compléments sur les morphismes, théorème
\ref{more-morphisms-theorem-stein-factorization-general}.
D'après le lemme \ref{lemma-relative-bottom-part-degenerates}, le morphisme
$\pi : S' \to S$ est fini étale.
Le morphisme $f : X \to S'$ est propre (d'après le théorème),
lisse (d'après Compléments sur les morphismes, lemme
\ref{more-morphisms-lemma-smooth-etale-permanence}) et à fibres géométriquement
connexes d'après le théorème de factorisation de Stein.
Dans la démonstration du lemme \ref{lemma-relative-bottom-part-degenerates},
nous avons vu que $\Omega_{X/S} = \Omega_{X/S'}$ puisque $S' \to S$ est étale.
Ainsi, $\Omega^\bullet_{X/S} = \Omega^\bullet_{X/S'}$.
Nous avons
$$
R^qf_*\Omega^p_{X/S} = \pi_*R^qf'_*\Omega^p_{X/S'}
$$
pour tous $p, q$, d'après la suite spectrale de Leray
(Cohomologie, lemme \ref{cohomology-lemma-relative-Leray}),
le fait que $\pi$ soit fini, donc affine, et
Cohomologie des schémas, lemme \ref{coherent-lemma-relative-affine-vanishing}
(nous utilisons aussi, bien entendu, que $R^qf'_*\Omega^p_{X'/S}$ est
quasi-cohérent).
Le morphisme du lemme s'obtient donc en appliquant $\pi_*$ à
$\text{d} : R^nf'_*\Omega^{n - 1}_{X/S'} \to R^nf'_*\Omega^n_{X/S'}$.
Autrement dit, pour démontrer le lemme, nous pouvons remplacer
$f : X \to S$ par $f' : X \to S'$ et nous ramener au cas étudié
dans le paragraphe suivant.

\medskip\noindent
Supposons que $f$ ait des fibres géométriquement connexes et que
$f_*\mathcal{O}_X = \mathcal{O}_S$.
Pour chaque $s \in S$, nous pouvons choisir un voisinage étale
$(S', s') \to (S, s)$ tel que le changement de base $X' \to S'$ de $S$
admette une section. Voir Compléments sur les morphismes, lemme
\ref{more-morphisms-lemma-etale-nbhd-dominates-smooth}.
D'après les remarques initiales de la démonstration, cela nous ramène au cas
étudié dans le paragraphe suivant.

\medskip\noindent
Supposons que $f$ ait des fibres géométriquement connexes,
que $f_*\mathcal{O}_X = \mathcal{O}_S$ et que nous disposions
d'une section $s : S \to X$ de $f$. Nous pouvons supposer, et nous supposons, que $S = \Spec(A)$
est affine. Le morphisme
$s^* : R\Gamma(X, \mathcal{O}_X) \to R\Gamma(S, \mathcal{O}_S) = A$
est une rétraction du morphisme $A \to R\Gamma(X, \mathcal{O}_X)$. Nous pouvons donc écrire
$$
R\Gamma(X, \mathcal{O}_X) = A \oplus P
$$,
où $P$ est le « noyau » de $s^*$. D'après
le lemme \ref{lemma-relative-bottom-part-degenerates}, l'objet $P$
de $D(A)$ est parfait d'amplitude de Tor dans $[1, n]$. Comme dans la démonstration
du lemme \ref{lemma-relative-duality-hodge}, nous voyons que
$H^n(X, \Omega^n_{X/S})$ est un $A$-module localement libre de rang $1$
(et même le dual de $A$, donc libre de rang $1$ ; nous allons bientôt choisir
un générateur, mais nous ne souhaitons pas vérifier qu'il s'agit du même,
et cela ne sera d'ailleurs pas nécessaire).

\medskip\noindent
Notons $Z \subset X$ l'image de $s$, qui est un sous-schéma fermé de $X$ d'après
Schémas, lemme \ref{schemes-lemma-section-immersion}.
Observons que $Z \to X$ est une immersion fermée régulière (et, a fortiori, régulière au sens de Koszul d'après
Diviseurs, lemme \ref{divisors-lemma-regular-quasi-regular-immersion}),
en vertu de
Diviseurs, lemme \ref{divisors-lemma-section-smooth-regular-immersion}.
Bien entendu, $Z \to X$ est de codimension $n$. D'après
la remarque \ref{remark-how-to-use},
nous pouvons donc considérer le morphisme
$$
\gamma^{0, 0} : H^0(Z, \Omega^0_{Z/S}) \longrightarrow H^n(X, \Omega^n_{X/S})
$$,
et nous posons $\xi = \gamma^{0, 0}(1) \in H^n(X, \Omega^n_{X/S})$.

\medskip\noindent
Nous affirmons que $\xi$ est un élément d'une base. En effet, puisque le changement de base est valable
en degré maximal (voir, par exemple, Limites, lemme
\ref{limits-lemma-higher-direct-images-zero-above-dimension-fibre}),
nous avons
$H^n(X, \Omega^n_{X/S}) \otimes_A k = H^n(X_k, \Omega^n_{X_k/k})$
pour tout morphisme d'anneaux $A \to k$. La compatibilité de
la construction de $\xi$ au changement de base,
voir le lemme \ref{lemma-gysin-transverse-global},
montre que l'image de $\xi$ dans $H^n(X_k, \Omega^n_{X_k/k})$
est non nulle d'après le lemme \ref{lemma-class-of-a-point} lorsque $k$ est un corps.
Ainsi, $\xi$ est une section partout non nulle d'un module inversible,
donc un générateur.

\medskip\noindent
Soit $\theta \in H^n(X, \Omega^{n - 1}_{X/S})$. Nous devons montrer que
$\text{d}(\theta)$ est nul dans $H^n(X, \Omega^n_{X/S})$.
Nous pouvons écrire $\text{d}(\theta) = a \xi$ pour un certain $a \in A$,
puisque $\xi$ est un élément d'une base. Il s'agit alors de montrer que $a = 0$.

\medskip\noindent
Considérons l'immersion fermée
$$
\Delta : X \to X \times_S X
$$
C'est aussi une section d'un morphisme lisse (à savoir l'une ou l'autre des projections),
donc une immersion à la fois régulière et régulière au sens de Koszul, de codimension $n$.
Nous pouvons ainsi considérer les morphismes
$$
\gamma^{p, q} :
H^q(X, \Omega^p_{X/S})
\longrightarrow
H^{q + n}(X \times_S X, \Omega^{p + n}_{X \times_S X/S})
$$
de la remarque \ref{remark-how-to-use}. Considérons l'image
$$
\gamma^{n - 1, n}(\theta) \in
H^{2n}(X \times_S X, \Omega^{2n - 1}_{X \times_S X})
$$
D'après le lemme \ref{lemma-de-rham-complex-product}, nous avons
$$
\Omega^{2n - 1}_{X \times_S X} =
\Omega^{n - 1}_{X/S} \boxtimes \Omega^n_{X/S} \oplus
\Omega^n_{X/S} \boxtimes \Omega^{n - 1}_{X/S}
$$
La formule de Künneth (Catégories dérivées des schémas, lemme \ref{perfect-lemma-kunneth}, ou
Catégories dérivées des schémas, lemme \ref{perfect-lemma-kunneth-single-sheaf})
donne
$$
H^{2n}(X \times_S X, \Omega^{n - 1}_{X/S} \boxtimes \Omega^n_{X/S}) =
H^n(X, \Omega^{n - 1}_{X/S}) \otimes_A H^n(X, \Omega^n_{X/S})
$$
et
$$
H^{2n}(X \times_S X, \Omega^n_{X/S} \boxtimes \Omega^{n - 1}_{X/S}) =
H^n(X, \Omega^n_{X/S}) \otimes_A H^n(X, \Omega^{n - 1}_{X/S})
$$
En effet, puisque nous nous plaçons en degré maximal, aucun groupe Tor supérieur
n'intervient. Comme $\xi$ est un générateur, cela signifie
que nous pouvons écrire
$$
\gamma^{n - 1, n}(\theta) = \theta_1 \otimes \xi + \xi \otimes \theta_2
$$
avec $\theta_1, \theta_2 \in H^n(X, \Omega^{n - 1}_{X/S})$.
En raisonnant exactement de la même manière, nous pouvons écrire
$$
\gamma^{n, n}(\xi) = b \xi \otimes \xi
$$
dans
$H^{2n}(X \times_S X, \Omega^{2n}_{X \times_S X/S}) =
H^n(X, \Omega^n_{X/S}) \otimes_A H^n(X, \Omega^n_{X/S})$
pour un certain $b \in H^0(S, \mathcal{O}_S)$.

\medskip\noindent
{\bf Affirmation :} $\theta_1 = \theta$, $\theta_2 = \theta$ et $b = 1$.
Montrons que cette affirmation implique le résultat voulu, $a = 0$.
En effet, d'après le lemme \ref{lemma-gysin-differential-hodge},
nous avons
$$
\gamma^{n, n}(\text{d}(\theta)) = \text{d}(\gamma^{n - 1, n}(\theta))
$$
Nos choix précédents donnent alors
$$
a \xi \otimes \xi =
\gamma^{n, n}(a\xi) =
\text{d}(\theta \otimes \xi + \xi \otimes \theta) =
a \xi \otimes \xi + (-1)^n a \xi \otimes \xi
$$
La dernière égalité provient du fait que le morphisme
$\text{d} : \Omega^{2n - 1}_{X \otimes_S X/S} \to \Omega^{2n}_{X \times_S X/S}$
est, d'après le lemme \ref{lemma-de-rham-complex-product},
la somme de la différentielle
$\text{d} \boxtimes 1 : \Omega^{n - 1}_{X/S} \boxtimes \Omega^n_{X/S}
\to \Omega^n_{X/S} \boxtimes \Omega^n_{X/S}$
et de la différentielle
$(-1)^n 1 \boxtimes \text{d} : \Omega^n_{X/S} \boxtimes \Omega^{n - 1}_{X/S}
\to \Omega^n_{X/S} \boxtimes \Omega^n_{X/S}$. Pour plus de détails, voir la discussion dans
la section \ref{section-kunneth} et dans
Catégories dérivées des schémas, section
\ref{perfect-section-kunneth-complexes}.
Puisque $\xi \otimes \xi$ est une base du module libre de rang $1$ sur $A$
$H^n(X, \Omega^n_{X/S}) \otimes_A H^n(X, \Omega^n_{X/S})$,
nous concluons que
$$
a = a + (-1)^n a \Rightarrow a = 0
$$,
comme voulu.

\medskip\noindent
Dans la suite de la démonstration, nous établissons l'affirmation ci-dessus. Notons
$\eta = \gamma^{0, 0}(1) \in H^n(X \times_S X, \Omega^n_{X \times_S X/S})$.
Puisque $\Omega^n_{X \times_S X/S} =
\bigoplus_{p + p' = n} \Omega^p_{X/S} \boxtimes \Omega^{p'}_{X/S}$,
nous pouvons écrire
$$
\eta = \eta_0 + \eta_1 + \ldots + \eta_n
$$,
où $\eta_p$ appartient à
$H^n(X \times_S X, \Omega^p_{X/S} \boxtimes \Omega^{n - p}_{X/S})$.
Pour $p = 0$, nous pouvons écrire
\begin{align*}
H^n(X \times_S X, \mathcal{O}_X \boxtimes \Omega^n_{X/S})
& =
H^n(R\Gamma(X, \mathcal{O}_X) \otimes_A^\mathbf{L}
R\Gamma(X, \Omega^n_{X/S})) \\
& =
A \otimes_A H^n(X, \Omega^n_{X/S}) \oplus
H^n(P \otimes_A^\mathbf{L} R\Gamma(X, \Omega^n_{X/S}))
\end{align*}
grâce à la décomposition $R\Gamma(X, \mathcal{O}_X) = A \oplus P$ obtenue plus haut.
Considérons le morphisme $(s, \text{id}) : X \to X \times_S X$.
Alors $(s, \text{id})^{-1}(\Delta) = Z$ au sens schématique.
Ainsi, $(s, \text{id})^*\eta = \xi$ d'après
le lemme \ref{lemma-gysin-transverse-global}. Cela signifie que
$$
\xi = (s, \text{id})^*\eta = (s^* \otimes \text{id})(\eta_0)
$$
Autrement dit, la première composante de $\eta_0$
dans la décomposition en somme directe ci-dessus est $\xi$. Nous pouvons donc écrire
$$
\eta_0 = 1 \otimes \xi + \eta'_0
$$
avec $\eta'_0 \in H^n(P \otimes_A^\mathbf{L} R\Gamma(X, \Omega^n_{X/S}))$.
Exactement de la même manière, pour $p = n$, nous pouvons écrire
\begin{align*}
H^n(X \times_S X, \Omega^n_{X/S} \boxtimes \mathcal{O}_X)
& =
H^n(R\Gamma(X, \Omega^n_{X/S}) \otimes_A^\mathbf{L}
R\Gamma(X, \mathcal{O}_X)) \\
& =
H^n(X, \Omega^n_{X/S}) \otimes_A A \oplus
H^n(R\Gamma(X, \Omega^n_{X/S}) \otimes_A^\mathbf{L} P)
\end{align*},
puis
$$
\eta_n = \xi \otimes 1 + \eta'_n
$$
avec $\eta'_n \in H^n(R\Gamma(X, \Omega^n_{X/S}) \otimes_A^\mathbf{L} P)$.

\medskip\noindent
Observons que $\text{pr}_1^*\theta = \theta \otimes 1$
et $\text{pr}_2^*\theta = 1 \otimes \theta$ sont des
classes de cohomologie de Hodge sur
$X \times_S X$ dont l'image inverse par la diagonale est $\theta$, cette diagonale étant notée $\Delta$.
D'après le lemme \ref{lemma-gysin-projection-global}, nous avons donc
$$
\theta_1 \otimes \xi + \xi \otimes \theta_2 =
\gamma^{n - 1, n}(\theta) =
(\theta \otimes 1) \cup \eta =
(1 \otimes \theta) \cup \eta
$$
dans l'anneau de cohomologie de Hodge de $X \times_S X$ sur $S$.
En termes de la décomposition en somme directe des modules
de différentielles de $X \times_S X/S$, nous obtenons
$$
\theta_1 \otimes \xi =
(\theta \otimes 1) \cup \eta_0
\quad\text{et}\quad
\xi \otimes \theta_2 =
(1 \otimes \theta) \cup \eta_n
$$
La formule $\eta_0 = 1 \otimes \xi + \eta'_0$ obtenue plus haut
montre que, pour établir $\theta_1 = \theta$, il suffit de prouver
$$
(\theta \otimes 1) \cup \eta'_0 = 0
$$
Pour cela, observons que le cup-produit avec $\theta \otimes 1$ est donné
par l'action en cohomologie du morphisme
$$
(P \otimes_A^\mathbf{L} R\Gamma(X, \Omega^n_{X/S}))[-n]
\xrightarrow{\theta \otimes 1}
R\Gamma(X, \Omega^{n - 1}_{X/S}) \otimes_A^\mathbf{L}
R\Gamma(X, \Omega^n_{X/S})
$$
dans la catégorie dérivée ; voir Cohomologie, remarque
\ref{cohomology-remark-cup-with-element-map-total-cohomology}.
Ce morphisme est le produit tensoriel dérivé des deux morphismes
$$
\theta : P[-n] \to R\Gamma(X, \Omega^{n - 1}_{X/S})
\quad\text{et}\quad
1 : R\Gamma(X, \Omega^n_{X/S}) \to R\Gamma(X, \Omega^n_{X/S})
$$
d'après Catégories dérivées des schémas, remarque
\ref{perfect-remark-annoying-compatibility}.
Or le premier est nul dans $D(A)$, car il va
d'un complexe parfait d'amplitude de Tor dans $[n + 1, 2n]$ vers un complexe
dont la cohomologie n'est éventuellement non nulle qu'en degrés $0, 1, \ldots, n$ ; voir
Compléments d'algèbre, lemme \ref{more-algebra-lemma-splitting-unique}.
Un raisonnement analogue établit l'annulation de
$(1 \otimes \theta) \cup \eta'_n$. Enfin,
exactement de la même manière, nous obtenons
$$
b \xi \otimes \xi = \gamma^{n, n}(\xi) = (\xi \otimes 1) \cup \eta_0
$$,
et nous concluons comme précédemment en montrant que
$(\xi \otimes 1) \cup \eta'_0 = 0$ par le même raisonnement que ci-dessus.
Cela achève la démonstration.
\end{proof}

\begin{proposition}
\label{proposition-relative-poincare-duality}
Soit $S$ un schéma quasi-compact et quasi-séparé. Soit $f : X \to S$
un morphisme propre et lisse de schémas dont toutes les fibres sont non vides
et équidimensionnelles de dimension $n$. Il existe un morphisme de
$\mathcal{O}_S$-modules
$$
t : R^{2n}f_*\Omega^\bullet_{X/S} \longrightarrow \mathcal{O}_S
$$,
unique à précomposition près par la multiplication par un élément inversible de
$H^0(X, \mathcal{O}_X)$, ayant la propriété suivante : l'accouplement
$$
Rf_*\Omega^\bullet_{X/S}
\otimes_{\mathcal{O}_S}^\mathbf{L}
Rf_*\Omega^\bullet_{X/S}[2n]
\longrightarrow
\mathcal{O}_S, \quad
(\xi, \xi') \longmapsto t(\xi \cup \xi')
$$
est un accouplement parfait de complexes parfaits sur $S$.
\end{proposition}

\begin{proof}
La démonstration est exactement la même que celle de la
proposition \ref{proposition-poincare-duality}.

\medskip\noindent
La suite spectrale relative de Hodge vers de Rham
$$
E_1^{p, q} = R^qf_*\Omega^p_{X/S} \Rightarrow R^{p + q}f_*\Omega^\bullet_{X/S}
$$
(section \ref{section-hodge-to-de-rham}), l'annulation
de $\Omega^i_{X/S}$ pour $i > n$, l'annulation donnée, par exemple, dans Limites, lemme
\ref{limits-lemma-higher-direct-images-zero-above-dimension-fibre},
et les résultats des lemmes \ref{lemma-relative-bottom-part-degenerates} et
\ref{lemma-relative-top-part-degenerates}
montrent que $R^0f_*\Omega_{X/S} = R^0f_*\mathcal{O}_X$
et $R^nf_*\Omega^n_{X/S} = R^{2n}f_*\Omega^\bullet_{X/S}$.
Plus précisément, ces identifications proviennent des morphismes
de complexes
$$
\Omega^\bullet_{X/S} \to \mathcal{O}_X[0]
\quad\text{et}\quad
\Omega^n_{X/S}[-n] \to \Omega^\bullet_{X/S}
$$
Choisissons $t : R^{2n}f_*\Omega_{X/S} \to \mathcal{O}_S$
qui corresponde, via cette identification, à un morphisme $t$ comme dans le
lemme \ref{lemma-relative-duality-hodge}.

\medskip\noindent
Abrégeons les notations en posant $\Omega^\bullet = \Omega^\bullet_{X/S}$.
Considérons le morphisme (\ref{equation-wedge}), qui s'écrit ici
$$
\wedge :
\text{Tot}(\Omega^\bullet \otimes_{f^{-1}\mathcal{O}_S} \Omega^\bullet)
\longrightarrow
\Omega^\bullet
$$
Pour chaque entier $p = 0, 1, \ldots, n$, ce morphisme s'annule sur le sous-complexe
$\text{Tot}(\sigma_{> p} \Omega^\bullet \otimes_{f^{-1}\mathcal{O}_S}
\sigma_{\geq n - p} \Omega^\bullet)$ pour des raisons de degré.
La restriction de $\wedge$ au sous-complexe
$\text{Tot}(\Omega^\bullet \otimes_{f^{-1}\mathcal{O}_S}
\geq_{n - p}\Omega^\bullet)$ se factorise donc par un morphisme de complexes
$$
\gamma_p :
\text{Tot}(\sigma_{\leq p} \Omega^\bullet \otimes_{f^{-1}\mathcal{O}_S}
\sigma_{\geq n - p} \Omega^\bullet)
\longrightarrow
\Omega^\bullet
$$
En procédant comme dans la section \ref{section-cup-product}, nous obtenons
des cup-produits relatifs
$$
Rf_*\sigma_{\leq p} \Omega^\bullet
\otimes_{\mathcal{O}_S}^\mathbf{L}
Rf_*\sigma_{\geq n - p}\Omega^\bullet
\longrightarrow
Rf_*\Omega^\bullet
$$
Nous montrerons par récurrence sur $p$ que ces cup-produits, composés avec $t$,
induisent des accouplements parfaits entre $Rf_*\sigma_{\leq p} \Omega^\bullet$
et $Rf_*\sigma_{\geq n - p}\Omega^\bullet[2n]$. Pour $p = n$,
il s'agit de l'assertion de la proposition.

\medskip\noindent
Le cas initial est $p = 0$. Nous avons alors
$$
Rf_*\sigma_{\leq p}\Omega^\bullet = Rf_*\mathcal{O}_X
\quad\text{et}\quad
Rf_*\sigma_{\geq n - p}\Omega^\bullet[2n] = Rf_*(\Omega^n[-n])[2n] =
Rf_*\Omega^n[n]
$$
Nous obtenons simplement l'accouplement
entre $Rf_*\mathcal{O}_X$ et $Rf_*\Omega^n[n]$ du
lemme \ref{lemma-relative-duality-hodge}, et le résultat est donc vrai.

\medskip\noindent
Étape de récurrence. Supposons le résultat vrai pour $p$.
Considérons alors le triangle distingué
$$
\Omega^{p + 1}[-p - 1] \to
\sigma_{\leq p + 1}\Omega^\bullet \to
\sigma_{\leq p}\Omega^\bullet \to
\Omega^{p + 1}[-p]
$$
et le triangle distingué
$$
\sigma_{\geq n - p}\Omega^\bullet \to
\sigma_{\geq n - p - 1}\Omega^\bullet \to
\Omega^{n - p - 1}[-n + p + 1] \to
(\sigma_{\geq n - p}\Omega^\bullet)[1]
$$
Observons que ce sont deux triangles distingués dans la catégorie homotopique
des complexes de faisceaux de $f^{-1}\mathcal{O}_S$-modules ; en particulier,
les morphismes $\sigma_{\leq p}\Omega^\bullet \to \Omega^{p + 1}[-p]$ et
$\Omega^{n - p - 1}[-d + p + 1] \to (\sigma_{\geq n - p}\Omega^\bullet)[1]$
sont donnés par de véritables morphismes de complexes, à savoir à l'aide de la différentielle
$\Omega^p \to \Omega^{p + 1}$ et de la différentielle
$\Omega^{n - p - 1} \to \Omega^{n - p}$.
Considérons les triangles distingués obtenus à partir de ceux-ci
en appliquant $Rf_*$ :
$$
\xymatrix{
Rf_*\sigma_{\leq p}\Omega^\bullet \ar[d]_a \\
Rf_*\Omega^{p + 1}[-p - 1] \ar[d]_b \\
Rf_*\sigma_{\leq p + 1}\Omega^\bullet \ar[d]_c \\
Rf_*\sigma_{\leq p}\Omega^\bullet \ar[d]_d \\
Rf_*\Omega^{p + 1}[-p - 1]
}
\quad\quad
\xymatrix{
Rf_*\sigma_{\geq n - p}\Omega^\bullet \\
Rf_*\Omega^{n - p - 1}[-n + p + 1] \ar[u]_{a'} \\
Rf_*\sigma_{\geq n - p - 1}\Omega^\bullet \ar[u]_{b'} \\
Rf_*\sigma_{\geq n - p}\Omega^\bullet \ar[u]_{c'} \\
Rf_*\Omega^{n - p - 1}[-n + p + 1] \ar[u]_{d'}
}
$$
Nous montrerons ci-dessous que les couples $(a, a')$, $(b, b')$, $(c, c')$ et
$(d, d')$ sont compatibles avec les accouplements donnés. Nous obtenons donc un
morphisme du triangle distingué de gauche vers le triangle distingué
obtenu en appliquant $R\SheafHom(-, \mathcal{O}_S)$ au triangle distingué
de droite. Par récurrence et d'après le lemme \ref{lemma-duality-hodge},
les accouplements construits ci-dessus entre les
complexes des première, deuxième, quatrième et cinquième
lignes sont parfaits, c'est-à-dire qu'ils définissent des isomorphismes après dualisation.
D'après Catégories dérivées, lemme \ref{derived-lemma-third-isomorphism-triangle},
nous en concluons que l'accouplement entre les complexes de la ligne médiane
est parfait, comme voulu.

\medskip\noindent
Soient $e : K \to K'$ et $e' : M' \to M$ des morphismes d'objets
de $D(\mathcal{O}_S)$, et soient
$K \otimes_{\mathcal{O}_S}^\mathbf{L} M \to \mathcal{O}_S$ et
$K' \otimes_{\mathcal{O}_S}^\mathbf{L} M' \to \mathcal{O}_S$
des accouplements. Nous disons qu'ils sont compatibles si le
diagramme
$$
\xymatrix{
K' \otimes_{\mathcal{O}_S}^\mathbf{L} M' \ar[d] &
K \otimes_{\mathcal{O}_S}^\mathbf{L} M'
\ar[l]^{e \otimes 1} \ar[d]^{1 \otimes e'} \\
\mathcal{O}_S &
K \otimes_{\mathcal{O}_S}^\mathbf{L} M \ar[l]
}
$$
est commutatif. Cela signifie en effet que le diagramme
$$
\xymatrix{
K \ar[r] \ar[d]_e & R\SheafHom(M, \mathcal{O}_S) \ar[d]^{R\SheafHom(e', -)} \\
K' \ar[r] & R\SheafHom(M', \mathcal{O}_S)
}
$$
est commutatif, ce qui suffit pour nos besoins.

\medskip\noindent
Démontrons cette compatibilité pour le couple $(c, c')$. Il suffit d'observer
que nous avons un diagramme commutatif
$$
\xymatrix{
\text{Tot}(\sigma_{\leq p} \Omega^\bullet \otimes_{f^{-1}\mathcal{O}_S}
\sigma_{\geq n - p} \Omega^\bullet) \ar[d]_{\gamma_p} &
\text{Tot}(\sigma_{\leq p + 1} \Omega^\bullet \otimes_{f^{-1}\mathcal{O}_S}
\sigma_{\geq n - p} \Omega^\bullet) \ar[l] \ar[d] \\
\Omega^\bullet &
\text{Tot}(\sigma_{\leq p + 1} \Omega^\bullet \otimes_{f^{-1}\mathcal{O}_S}
\sigma_{\geq n - p - 1} \Omega^\bullet) \ar[l]_-{\gamma_{p + 1}}
}
$$
La fonctorialité du cup-produit donne alors la commutativité du
diagramme voulu.

\medskip\noindent
De même, pour le couple $(b, b')$, nous utilisons le diagramme commutatif
$$
\xymatrix{
\text{Tot}(\sigma_{\leq p + 1} \Omega^\bullet \otimes_{f^{-1}\mathcal{O}_S}
\sigma_{\geq n - p - 1} \Omega^\bullet) \ar[d]_{\gamma_{p + 1}} &
\text{Tot}(\Omega^{p + 1}[-p - 1] \otimes_{f^{-1}\mathcal{O}_S}
\sigma_{\geq n - p - 1} \Omega^\bullet) \ar[l] \ar[d] \\
\Omega^\bullet &
\Omega^{p + 1}[-p - 1]
\otimes_{f^{-1}\mathcal{O}_S}
\Omega^{n - p - 1}[-n + p + 1] \ar[l]_-\wedge
}
$$

\medskip\noindent
Pour les couples $(d, d')$ et $(a, a')$, nous utilisons le diagramme commutatif
$$
\xymatrix{
\Omega^{p + 1}[-p] \otimes_{f^{-1}\mathcal{O}_S}
\Omega^{n - p - 1}[-n + p] \ar[d] &
\text{Tot}(\sigma_{\leq p}\Omega^\bullet \otimes_{f^{-1}\mathcal{O}_S}
\Omega^{n - p - 1}[-n + p]) \ar[l] \ar[d] \\
\Omega^\bullet &
\text{Tot}(\sigma_{\leq p}\Omega^\bullet \otimes_{f^{-1}\mathcal{O}_S}
\sigma_{\geq n - p}\Omega^\bullet) \ar[l]
}
$$

\medskip\noindent
Nous omettons la démonstration de l'unicité de $t$ à
précomposition près par la multiplication par un élément inversible de $H^0(X, \mathcal{O}_X)$.
\end{proof}











\begin{multicols}{2}[\section{Autres chapitres}]
\noindent
Préliminaires
\begin{enumerate}
\item \hyperref[introduction-section-phantom]{Introduction}
\item \hyperref[conventions-section-phantom]{Conventions}
\item \hyperref[sets-section-phantom]{Théorie des ensembles}
\item \hyperref[categories-section-phantom]{Catégories}
\item \hyperref[topology-section-phantom]{Topologie}
\item \hyperref[sheaves-section-phantom]{Faisceaux sur les espaces}
\item \hyperref[sites-section-phantom]{Sites et faisceaux}
\item \hyperref[stacks-section-phantom]{Champs}
\item \hyperref[fields-section-phantom]{Corps}
\item \hyperref[algebra-section-phantom]{Algèbre commutative}
\item \hyperref[brauer-section-phantom]{Groupes de Brauer}
\item \hyperref[homology-section-phantom]{Algèbre homologique}
\item \hyperref[derived-section-phantom]{Catégories dérivées}
\item \hyperref[simplicial-section-phantom]{Méthodes simpliciales}
\item \hyperref[more-algebra-section-phantom]{Compléments d'algèbre}
\item \hyperref[smoothing-section-phantom]{Lissage des morphismes d'anneaux}
\item \hyperref[modules-section-phantom]{Faisceaux de modules}
\item \hyperref[sites-modules-section-phantom]{Modules sur les sites}
\item \hyperref[injectives-section-phantom]{Objets injectifs}
\item \hyperref[cohomology-section-phantom]{Cohomologie des faisceaux}
\item \hyperref[sites-cohomology-section-phantom]{Cohomologie sur les sites}
\item \hyperref[dga-section-phantom]{Algèbre différentielle graduée}
\item \hyperref[dpa-section-phantom]{Algèbre à puissances divisées}
\item \hyperref[sdga-section-phantom]{Faisceaux différentiels gradués}
\item \hyperref[hypercovering-section-phantom]{Hyperrecouvrements}
\end{enumerate}
Schémas
\begin{enumerate}
\setcounter{enumi}{25}
\item \hyperref[schemes-section-phantom]{Schémas}
\item \hyperref[constructions-section-phantom]{Constructions de schémas}
\item \hyperref[properties-section-phantom]{Propriétés des schémas}
\item \hyperref[morphisms-section-phantom]{Morphismes de schémas}
\item \hyperref[coherent-section-phantom]{Cohomologie des schémas}
\item \hyperref[divisors-section-phantom]{Diviseurs}
\item \hyperref[limits-section-phantom]{Limites de schémas}
\item \hyperref[varieties-section-phantom]{Variétés}
\item \hyperref[topologies-section-phantom]{Topologies sur les schémas}
\item \hyperref[descent-section-phantom]{Descente}
\item \hyperref[perfect-section-phantom]{Catégories dérivées de schémas}
\item \hyperref[more-morphisms-section-phantom]{Compléments sur les morphismes}
\item \hyperref[flat-section-phantom]{Compléments sur la platitude}
\item \hyperref[groupoids-section-phantom]{Schémas en groupoïdes}
\item \hyperref[more-groupoids-section-phantom]{Compléments sur les schémas en groupoïdes}
\item \hyperref[etale-section-phantom]{Morphismes étales de schémas}
\end{enumerate}
Thèmes de théorie des schémas
\begin{enumerate}
\setcounter{enumi}{41}
\item \hyperref[chow-section-phantom]{Homologie de Chow}
\item \hyperref[intersection-section-phantom]{Théorie de l'intersection}
\item \hyperref[pic-section-phantom]{Schémas de Picard des courbes}
\item \hyperref[weil-section-phantom]{Théories cohomologiques de Weil}
\item \hyperref[adequate-section-phantom]{Modules adéquats}
\item \hyperref[dualizing-section-phantom]{Complexes dualisants}
\item \hyperref[duality-section-phantom]{Dualité pour les schémas}
\item \hyperref[discriminant-section-phantom]{Discriminants et différentes}
\item \hyperref[derham-section-phantom]{Cohomologie de de Rham}
\item \hyperref[local-cohomology-section-phantom]{Cohomologie locale}
\item \hyperref[algebraization-section-phantom]{Géométrie algébrique et formelle}
\item \hyperref[curves-section-phantom]{Courbes algébriques}
\item \hyperref[resolve-section-phantom]{Résolution des surfaces}
\item \hyperref[models-section-phantom]{Réduction semi-stable}
\item \hyperref[functors-section-phantom]{Foncteurs et morphismes}
\item \hyperref[equiv-section-phantom]{Catégories dérivées de variétés}
\item \hyperref[pione-section-phantom]{Groupes fondamentaux de schémas}
\item \hyperref[etale-cohomology-section-phantom]{Cohomologie étale}
\item \hyperref[crystalline-section-phantom]{Cohomologie cristalline}
\item \hyperref[proetale-section-phantom]{Cohomologie pro-étale}
\item \hyperref[relative-cycles-section-phantom]{Cycles relatifs}
\item \hyperref[more-etale-section-phantom]{Compléments de cohomologie étale}
\item \hyperref[trace-section-phantom]{La formule des traces}
\end{enumerate}
Espaces algébriques
\begin{enumerate}
\setcounter{enumi}{64}
\item \hyperref[spaces-section-phantom]{Espaces algébriques}
\item \hyperref[spaces-properties-section-phantom]{Propriétés des espaces algébriques}
\item \hyperref[spaces-morphisms-section-phantom]{Morphismes d'espaces algébriques}
\item \hyperref[decent-spaces-section-phantom]{Espaces algébriques décents}
\item \hyperref[spaces-cohomology-section-phantom]{Cohomologie des espaces algébriques}
\item \hyperref[spaces-limits-section-phantom]{Limites d'espaces algébriques}
\item \hyperref[spaces-divisors-section-phantom]{Diviseurs sur les espaces algébriques}
\item \hyperref[spaces-over-fields-section-phantom]{Espaces algébriques sur des corps}
\item \hyperref[spaces-topologies-section-phantom]{Topologies sur les espaces algébriques}
\item \hyperref[spaces-descent-section-phantom]{Descente et espaces algébriques}
\item \hyperref[spaces-perfect-section-phantom]{Catégories dérivées d'espaces}
\item \hyperref[spaces-more-morphisms-section-phantom]{Compléments sur les morphismes d'espaces}
\item \hyperref[spaces-flat-section-phantom]{Platitude sur les espaces algébriques}
\item \hyperref[spaces-groupoids-section-phantom]{Groupoïdes dans les espaces algébriques}
\item \hyperref[spaces-more-groupoids-section-phantom]{Compléments sur les groupoïdes dans les espaces}
\item \hyperref[bootstrap-section-phantom]{Amorçage}
\item \hyperref[spaces-pushouts-section-phantom]{Sommes amalgamées d'espaces algébriques}
\end{enumerate}
Thèmes de géométrie
\begin{enumerate}
\setcounter{enumi}{81}
\item \hyperref[spaces-chow-section-phantom]{Groupes de Chow des espaces}
\item \hyperref[groupoids-quotients-section-phantom]{Quotients de groupoïdes}
\item \hyperref[spaces-more-cohomology-section-phantom]{Compléments sur la cohomologie des espaces}
\item \hyperref[spaces-simplicial-section-phantom]{Espaces simpliciaux}
\item \hyperref[spaces-duality-section-phantom]{Dualité pour les espaces}
\item \hyperref[formal-spaces-section-phantom]{Espaces algébriques formels}
\item \hyperref[restricted-section-phantom]{Algébrisation des espaces formels}
\item \hyperref[spaces-resolve-section-phantom]{Résolution des surfaces revisitée}
\end{enumerate}
Théorie des déformations
\begin{enumerate}
\setcounter{enumi}{89}
\item \hyperref[formal-defos-section-phantom]{Théorie formelle des déformations}
\item \hyperref[defos-section-phantom]{Théorie des déformations}
\item \hyperref[cotangent-section-phantom]{Le complexe cotangent}
\item \hyperref[examples-defos-section-phantom]{Problèmes de déformation}
\end{enumerate}
Champs algébriques
\begin{enumerate}
\setcounter{enumi}{93}
\item \hyperref[algebraic-section-phantom]{Champs algébriques}
\item \hyperref[examples-stacks-section-phantom]{Exemples de champs}
\item \hyperref[stacks-sheaves-section-phantom]{Faisceaux sur les champs algébriques}
\item \hyperref[criteria-section-phantom]{Critères de représentabilité}
\item \hyperref[artin-section-phantom]{Axiomes d'Artin}
\item \hyperref[quot-section-phantom]{Espaces de Quot et de Hilbert}
\item \hyperref[stacks-properties-section-phantom]{Propriétés des champs algébriques}
\item \hyperref[stacks-morphisms-section-phantom]{Morphismes de champs algébriques}
\item \hyperref[stacks-limits-section-phantom]{Limites de champs algébriques}
\item \hyperref[stacks-cohomology-section-phantom]{Cohomologie des champs algébriques}
\item \hyperref[stacks-perfect-section-phantom]{Catégories dérivées de champs}
\item \hyperref[stacks-introduction-section-phantom]{Introduction aux champs algébriques}
\item \hyperref[stacks-more-morphisms-section-phantom]{Compléments sur les morphismes de champs}
\item \hyperref[stacks-geometry-section-phantom]{La géométrie des champs}
\end{enumerate}
Thèmes de théorie des espaces de modules
\begin{enumerate}
\setcounter{enumi}{107}
\item \hyperref[moduli-section-phantom]{Champs de modules}
\item \hyperref[moduli-curves-section-phantom]{Espaces de modules de courbes}
\end{enumerate}
Divers
\begin{enumerate}
\setcounter{enumi}{109}
\item \hyperref[examples-section-phantom]{Exemples}
\item \hyperref[exercises-section-phantom]{Exercices}
\item \hyperref[guide-section-phantom]{Guide bibliographique}
\item \hyperref[desirables-section-phantom]{Desiderata}
\item \hyperref[coding-section-phantom]{Style de programmation}
\item \hyperref[obsolete-section-phantom]{Obsolète}
\item \hyperref[fdl-section-phantom]{Licence de documentation libre GNU}
\item \hyperref[index-section-phantom]{Index généré automatiquement}
\end{enumerate}
\end{multicols}


\bibliography{my}
\bibliographystyle{amsalpha}

\end{document}
